% Rapport de l'equipe Tempete sur la tache Unlearning (31-unlearning).
% Compilation : make -C rapport   (latexmk, LuaLaTeX + biber, TeX Live 2026)
% Regle du depot : chaque commit sur la branche alex met a jour ce rapport et son PDF.
\documentclass[11pt,a4paper]{article}

% ----- Langue, polices ---------------------------------------------------------------------
\usepackage[french]{babel}
\usepackage{amsmath}
\usepackage{fontspec}
\usepackage{unicode-math}
\setmainfont{Libertinus Serif}
\setsansfont{Libertinus Sans}
\setmonofont{Source Code Pro}[Scale=0.84]
\setmathfont{Libertinus Math}
\usepackage[final]{microtype}
\usepackage{csquotes}

% ----- Mise en page -----------------------------------------------------------------------
\usepackage[a4paper,top=2.6cm,bottom=2.6cm,left=2.7cm,right=2.7cm,headheight=14pt]{geometry}
\usepackage{xcolor}
\definecolor{encre}{HTML}{1B2A41}
\definecolor{bleu}{HTML}{1F5F8B}
\definecolor{sarcelle}{HTML}{2A9D8F}
\definecolor{ambre}{HTML}{E9A23B}
\definecolor{brique}{HTML}{C0504D}
\definecolor{brume}{HTML}{EEF3F7}
\usepackage{titlesec}
\titleformat{\section}{\Large\sffamily\bfseries\color{encre}}{\color{bleu}\thesection}{0.8em}{}
\titleformat{\subsection}{\large\sffamily\bfseries\color{encre}}{\color{bleu}\thesubsection}{0.7em}{}
\titlespacing*{\section}{0pt}{1.6em}{0.7em}
\titlespacing*{\subsection}{0pt}{1.2em}{0.5em}
\usepackage{fancyhdr}
\pagestyle{fancy}
\fancyhf{}
\fancyhead[L]{\small\sffamily\color{bleu}Tempête \textperiodcentered{} Unlearning}
\fancyhead[R]{\small\sffamily\color{gray}Rapport de travail}
\fancyfoot[C]{\small\sffamily\thepage}
\renewcommand{\headrulewidth}{0.3pt}
\usepackage{enumitem}
\setlist{itemsep=0.25em,topsep=0.4em}
\usepackage{caption}
\usepackage{xurl}
\captionsetup{font=small,labelfont={bf,sf,color=bleu},margin=1em}
\usepackage{booktabs}
\usepackage{tabularx}
\usepackage{siunitx}
\sisetup{locale=FR,group-minimum-digits=4,detect-all,uncertainty-mode=separate}

% ----- Figures -----------------------------------------------------------------------------
\usepackage{pgfplots}
\pgfplotsset{compat=1.18}
\usetikzlibrary{babel,arrows.meta,positioning,fit,backgrounds,calc}
\usepackage[most]{tcolorbox}
\newtcolorbox{retenir}[1][À retenir]{enhanced,breakable,colback=brume,colframe=bleu,
  boxrule=0pt,leftrule=2.5pt,arc=0pt,outer arc=0pt,left=8pt,right=8pt,top=5pt,bottom=5pt,
  fonttitle=\sffamily\bfseries\color{bleu},coltitle=bleu,colbacktitle=brume,
  title=#1,attach title to upper={\ \ }}
\newtcolorbox{avertissement}{enhanced,breakable,colback=ambre!8,colframe=ambre,
  boxrule=0pt,leftrule=2.5pt,arc=0pt,outer arc=0pt,left=8pt,right=8pt,top=5pt,bottom=5pt}

% ----- Bibliographie, liens ----------------------------------------------------------------
\usepackage[backend=bibtex,style=numeric-comp,sorting=none,maxbibnames=6,giveninits=true,
  url=true,doi=false,isbn=false,eprint=true]{biblatex}
\addbibresource{references.bib}
\usepackage[colorlinks=true,linkcolor=bleu,citecolor=sarcelle,urlcolor=bleu,
  pdftitle={Tempête : désapprendre sans données},pdfauthor={Équipe Tempête}]{hyperref}
\usepackage[nameinlink,noabbrev,french]{cleveref}

% ----- Notations ---------------------------------------------------------------------------
\newcommand{\KL}{\operatorname{KL}}
\newcommand{\clip}{\operatorname{clip}}
\newcommand{\Exp}{\operatorname*{moy}}
\newcommand{\orig}{\mathrm{orig}}
\newcommand{\gold}{\mathrm{gold}}
\newcommand{\sub}{\mathrm{sub}}
\newcommand{\code}[1]{\texttt{#1}}
\newcommand{\chemin}[1]{\nolinkurl{#1}}% chemin de fichier, coupable en fin de ligne

\begin{document}

% ----- Bandeau de titre ----------------------------------------------------------------------
\thispagestyle{empty}
\begin{tcolorbox}[enhanced,colback=encre,colframe=encre,arc=0pt,outer arc=0pt,boxrule=0pt,
  left=14pt,right=14pt,top=14pt,bottom=14pt]
  {\sffamily\small\color{sarcelle!70!white}EUROPEAN AI CHAMPIONSHIP PARIS 2026 \textperiodcentered{} CISPA \textperiodcentered{} TÂCHE 31-UNLEARNING}\par\medskip
  {\sffamily\bfseries\fontsize{24}{28}\selectfont\color{white}Désapprendre sans données}\par\smallskip
  {\sffamily\large\color{white!85!bleu}Méthodes et banc d'essai local pour la tâche Unlearning}\par\bigskip
  {\sffamily\small\color{white!75!bleu}Équipe Tempête \textperiodcentered{} rapport de travail
  \textperiodcentered{} version du 22 novembre 2026 \textperiodcentered{} branche \code{alex}}
\end{tcolorbox}

\vspace{0.6em}
\begin{retenir}[Résumé]
On nous fournit un ViT-B/16 entraîné sur 100 classes et \num{3500} images qu'il doit
\enquote{oublier}, sans étiquette ni aucune autre donnée. Le score officiel ne mesure pas un taux
d'oubli mais une \emph{distance} à des modèles réentraînés sans ces images. Ce rapport décrit
(i) ce que nous avons mesuré sur le modèle : une mémorisation extrême du forget set, et un oubli
propre à chaque image en essai réduit, qui ne l'est plus à pleine échelle ; (ii) notre méthode
principale, une \emph{région de confiance} en divergence de Kullback-Leibler qui transpose SCRUB+R
à un cadre sans données, sa V2, dont la cible est propre à chaque image, et sa V3, dont la cible
vise séparément les deux objets du score ; (iii) un banc
d'essai local, versionné et recalculable, qui applique les formules exactes du sujet en
remplaçant chaque pièce manquante par une doublure déclarée ; (iv) trois campagnes à pleine
échelle, puis une campagne de la V4, 155 runs en tout. Bilan mesuré : une fois son arrêt testé à
chaque pas, la V1 atteint un \code{score\_min} de \num{0.0637 +- 0.0025} (trois graines) et ne
dépend que de son rayon ; la V2, réglée à $\tau=1$, atteint \num{0.0831 +- 0.0042}, soit
\SI{30}{\percent} de plus, graine par graine, à divergence et à retain égaux ; sous la pire des six
doublures du banc \code{bench-v4}, \num{0.0786 +- 0.0030} (+\SI{23}{\percent}), et +\SI{12}{\percent}
seulement sous la famille de doublures étrangère à sa cible. Elle rejoint SCRUB
dosé (\num{0.098 +- 0.022}, meilleur run \num{0.118}) dans la limite du bruit de ce dernier, en
étant cinq fois plus stable ; mais SCRUB ancre sa rétention sur des images extérieures au forget
set, que le sujet exclut pour une méthode soumise. La V3, testée par 23 runs dans le banc \code{bench-v4} (deux familles de
doublures du gold), fait jeu égal avec la V2 à son meilleur réglage (\num{0.0798 +- 0.0012} contre
\num{0.0786 +- 0.0030}) : la cible n'est plus le levier, la localité l'est. Relue, la campagne de la
V1 dit pourquoi : environ \SI{88}{\percent} de son mouvement est commun à toutes les images
mémorisées, oubliées ou non. D'où (v) une V4, mesurée par 26 runs aux seuils posés avant : un
proxy de rétention tiré du forget set lui-même (on oublie un pli après l'autre en tenant les
autres), le rayon comme contrainte active, et un SCRUB conforme au sujet. Les deux premiers sont
réfutés : l'ancre freine l'oubli partout sans le localiser, la contrainte ne lie pas. Le troisième
est confirmé : ancré sur des vues fortes du forget set et redosé (lr \num{3e-4}), SCRUB atteint
\num{0.1085 +- 0.0068} sur trois graines appariées, \SI{36}{\percent} de plus que la V3, au niveau
de SCRUB aux images externes : c'est la meilleure méthode conforme mesurée. Il gagne par
l'efficacité de son oubli, pas par la localité, qui reste le défaut commun. Puis (vi) une V5 qui change de direction : une montée de
gradient bornée image par image, réservée aux images où l'original doute sur des vues fortes, dosée
par un niveau choisi sur un pli que rien ne touche. Son étape 0, sans entraînement, la réfute aux
seuils posés avant : sous la famille de doublures que ses vues ne partagent pas, sa sélection
désigne les erreurs au hasard, et les deux familles de doublures ne s'accordent pas sur les images
difficiles (corrélation de rang nulle). Entraînée ensuite comme mesure de mécanisme (huit runs), la
bande est la plus locale des méthodes mesurées à divergence égale (le fold d'audit bouge
\numrange{0.25}{0.39} fois autant que les images oubliées), sans rien devoir à la sélection ; mais son
score est nul : les doublures du gold sont moins sûres que l'original sur toutes les images, et la
localité refuse ce mouvement commun. Puis (vii) V\_Dumb, la
base la plus simple : chaque image oubliée monte jusqu'à une perte cible et pas au-delà, ses vues
fortes restent à la sortie de l'original. Mesurée par onze runs, elle n'a pas de dose utilisable :
avec son ancre, le modèle ne bouge pas ; un peu moins d'ancre, et il s'effondre (meilleur run
\num{0.019}). Et (viii) une V5.1, sous un objectif fixé par le responsable, un retain d'au moins
\num{0.8} : elle corrige la loi de force de la V5 (bande en marge, retenue en KL à $T=2$, dose sous
ce plancher). Mesurée par onze runs, elle est abandonnée : la marge fait dépasser la bande en
quelques pas, et au seul taux où elle garde une dose, elle score \num{0.080} mais tombe sous le
plancher (retain \num{0.725}) sans oubli local. Puis (ix) V\_Fusion, SCRUB conforme dont l'ancre de
vues devient un lot d'images synthétiques tirées du seul modèle original, pour tenir ce plancher à
une dose plus forte. Ses images passent quatre portes de qualité sur cinq, mais elles sont saturées :
le modèle cassé par la montée reste d'accord avec l'original sur elles, et les quatre bras d'ancre
s'effondrent (retain \numrange{0.010}{0.025}). Enfin (x) V\_Dumb2, qui garde l'idée
de V\_Dumb et lui donne la dynamique de SCRUB conforme : le plafond moyen de ses max-steps devient
un niveau par image, borné des deux côtés ; elle reprend la bande en marge que la V5.1 a vue
dépasser son plafond. Codée, avec une porte sans entraînement, aucun run. Ces chiffres
sont des distances à une doublure du gold, sur l'échelle du score officiel (0 pour l'original, 1
pour le gold) ; le score officiel lui-même n'est pas mesuré : rien n'a été soumis au serveur.
\end{retenir}

\section{Le problème}\label{sec:probleme}

\subsection{La tâche}
Le modèle fourni est un ViT-B/16~\cite{dosovitskiy2021} de la bibliothèque
\code{timm}~\cite{wightman2019timm} (\code{vit\_base\_patch16\_224}, LayerScale initialisé à
\num{e-4}~\cite{touvron2021cait}, \num{85893988} paramètres), entraîné sur un jeu privé de 100
classes~\cite{sujet2026,sprintml2026}. Le forget set compte \num{3500} images de cet entraînement,
en pixels bruts $224\times224$, \emph{sans étiquettes}. Le sujet est explicite : la méthode doit
fonctionner \enquote{à partir du modèle original et du forget set seuls}~\cite{sujet2026}. On
soumet un dictionnaire de poids float32 de même architecture, que le serveur évalue en mode
\code{eval} sous autocast float16.

Le désapprentissage exact passe par un réentraînement, éventuellement organisé pour être bon
marché~\cite{bourtoule2021} ; c'est exclu ici, faute de données. Restent les méthodes
approchées~\cite{golatkar2020,kurmanji2023,chundawat2023,fan2024salun,zhang2024npo}, dont
l'évaluation repose sur la comparaison avec un modèle réentraîné, souvent par des attaques
d'inférence d'appartenance~\cite{shokri2017,carlini2022} ; la première compétition NeurIPS sur
le sujet a montré combien cette évaluation est délicate~\cite{triantafillou2024}.

\subsection{Le score officiel}\label{sec:score}
Notons $L_m(x)$ la perte d'entropie croisée du modèle $m$ sur l'image $x$ (avec sa vraie
étiquette), et $p_m(x)=\operatorname{softmax}(z_m(x)/2)$ sa distribution prédite à température
2. Le \enquote{gold} est la moyenne de plusieurs modèles réentraînés sans le forget set : ses
moyennes et écarts-types de pertes sont moyennés modèle par modèle, et $p_\gold$ est la moyenne de
leurs distributions~\cite{sujet2026}. Pour chaque groupe d'images (images exactes, copies altérées
cachées) :
\begin{align}
  d_m &= \bigl|\Exp L_m - \Exp L_\gold\bigr| + \bigl|\operatorname{std} L_m - \operatorname{std} L_\gold\bigr|,
  & S_{\text{ms}} &= 1 - \clip\!\Bigl(\frac{d_\sub}{d_\orig},0,1\Bigr), \label{eq:ms}\\
  S_{\text{pred}} &= 1 - \clip\!\Bigl(\frac{\Exp_x \KL\bigl(p_\gold(x)\,\|\,p_\sub(x)\bigr)}
  {\Exp_x \KL\bigl(p_\gold(x)\,\|\,p_\orig(x)\bigr)},0,1\Bigr), \label{eq:pred}
\end{align}
où $\KL$ est la divergence de Kullback-Leibler~\cite{kullback1951}. Chaque score est moyenné sur
les deux groupes, puis
\begin{gather}
  \text{forget} = \bar S_{\text{ms}}\,\bar S_{\text{pred}},\qquad
  \text{final} = \text{forget}\times\text{retain},\label{eq:final}\\
  \text{retain} = \clip\!\Bigl(\frac{\mathrm{acc}_\sub}{\mathrm{acc}_\orig},0,1\Bigr)\cdot
  \frac{1}{1+\Exp_x \KL\bigl(p_\orig(x)\,\|\,p_\sub(x)\bigr)},\label{eq:retain}
\end{gather}
le retain étant mesuré sur un control set caché. Une porte annule tout si la précision sur un
validation set caché passe sous \SI{70}{\percent}. Le classement public n'utilise qu'environ
\SI{30}{\percent} des données de notation ; le classement final, les \SI{70}{\percent} restants.

\begin{retenir}
\begin{itemize}[leftmargin=1.2em]
  \item \textbf{Ne rien faire vaut exactement 0} : avec $\sub=\orig$, les deux rapports valent 1.
  \item \textbf{C'est une distance, pas un taux d'oubli} : dépasser le gold (être \enquote{plus faux}
    que lui) est pénalisé, et dès que $d_\sub\ge d_\orig$ le score tombe à 0.
  \item \textbf{La distribution complète compte}, image par image : quelle classe fausse le gold
    préfère, pas seulement sa confiance.
  \item \textbf{Tout se multiplie} : une dérive sur le retain coûte directement sur le score final.
\end{itemize}
\end{retenir}

\section{Ce que nous avons mesuré sur le modèle}\label{sec:mesures}

\subsection{Une mémorisation extrême}
Les réseaux profonds peuvent mémoriser leurs données d'entraînement~\cite{zhang2017}, et cette
mémorisation est souvent nécessaire pour les exemples rares~\cite{feldman2020}. Le modèle fourni
en est un cas limite (\cref{tab:memo}) : confiance moyenne de \num{0.9997} sur le forget set, marge
médiane de 12 logits entre les deux premières classes, et aucune réaction aux altérations
légères. L'original a donc, sur ces images, une perte quasi nulle et un écart-type quasi nul ; un
gold, qui ne les a jamais vues, aura une perte modérée, un écart-type marqué et quelques erreurs.

\begin{table}[ht]
  \centering\small
  \caption{Comportement de l'original sur le forget set (mesures de l'équipe, 128 images, Apple
  M5, script jetable consigné dans \chemin{docs/hackathon/unlearning.md}).}\label{tab:memo}
  \begin{tabular}{@{}p{9.4cm}r@{}}
    \toprule
    Mesure & Valeur\\
    \midrule
    Confiance moyenne sur les images exactes & \num{0.9997}\\
    Marge médiane top-1 -- top-2 (logits) & \num{12.2}\\
    Confiance sur retournement / recadrage \SIrange{60}{90}{\percent} / JPEG q50 & \num{1.000} / \num{1.000} / \num{0.999}\\
    Recadrage de \SI{35}{\percent} de l'aire : confiance (accord de l'argmax) & \num{0.926} (\SI{87.5}{\percent})\\
    Recadrage de \SI{15}{\percent} de l'aire : confiance (accord de l'argmax) & \num{0.808} (\SI{43.8}{\percent})\\
    Confiance sur images hors distribution (COCO, DIV2K, Kodak) & \num{0.75}\\
    \bottomrule
  \end{tabular}
\end{table}

\subsection{Calibrer une doublure du gold}\label{sec:calibration}
Pour approcher le comportement d'un gold, on interroge l'original là où il n'est pas mémorisé :
sur des \emph{vues fortes} de chaque image (recadrage d'une aire tirée dans $[a,a+0{,}1]$, puis
variation de couleur). Un gold traite une image oubliée comme une image inédite ; sa précision y
est sa précision de test, inconnue mais supérieure à \SI{70}{\percent} puisque la porte l'exige.
La \cref{fig:calibration} montre qu'on règle cette précision avec $a$ : l'accord avec l'argmax de
l'image exacte passe de \SI{33}{\percent} ($a=0{,}15$) à \SI{89}{\percent} ($a=0{,}55$). La
confiance reste élevée même quand la réponse change, signe de la sur-confiance bien connue des
réseaux modernes~\cite{guo2017}. Nous retenons trois forces, $a\in\{0{,}40;0{,}50;0{,}60\}$, qui
encadrent un gold à environ 78, 86 et \SI{92}{\percent} de précision (les deux premières par
interpolation, la dernière par extrapolation).

\begin{figure}[ht]
  \centering
  \begin{tikzpicture}
    \begin{axis}[width=0.86\linewidth,height=6.2cm,
        xlabel={aire minimale $a$ du recadrage},ylabel={\%},
        xmin=0.12,xmax=0.64,ymin=25,ymax=100,
        xticklabel style={/pgf/number format/.cd,fixed,precision=2,use comma},
        legend pos=south east,legend style={font=\small,draw=none,fill=white,fill opacity=0.85},
        grid=major,grid style={gray!18},axis line style={gray!60},tick style={gray!60},
        label style={font=\small},tick label style={font=\small}]
      \fill[sarcelle!12] (axis cs:0.12,78) rectangle (axis cs:0.64,92);
      \node[font=\scriptsize\sffamily,text=sarcelle!80!black,anchor=west] at (axis cs:0.125,90)
        {précision plausible d'un gold};
      \foreach \x in {0.40,0.50,0.60}{\edef\temp{\noexpand\draw[bleu,dashed,thin]
        (axis cs:\x,25) -- (axis cs:\x,100);}\temp}
      \addplot[bleu,very thick,mark=*,mark size=1.8pt] coordinates
        {(0.15,32.6) (0.20,43.4) (0.25,53.1) (0.30,61.5) (0.35,73.0) (0.45,83.4) (0.55,89.5)};
      \addlegendentry{accord de l'argmax avec l'image exacte}
      \addplot[ambre,very thick,mark=square*,mark size=1.6pt] coordinates
        {(0.15,78.0) (0.20,80.6) (0.25,84.7) (0.30,85.7) (0.35,88.4) (0.45,92.3) (0.55,94.6)};
      \addlegendentry{confiance moyenne}
    \end{axis}
  \end{tikzpicture}
  \caption{Original interrogé sur des vues fortes de 256 images oubliées (2 vues par image, Apple
  M2). En pointillés, les trois forces retenues pour les doublures du gold du banc local.}
  \label{fig:calibration}
\end{figure}

\subsection{L'oubli est propre à chaque image}\label{sec:local}
Une mesure pas à pas (\cref{tab:local}) a révélé un fait structurant. Après une mise à jour sur un
lot de 8 images, les images qui n'y figuraient pas restent à une divergence de l'ordre de
\num{e-6} avec l'original, alors que la moyenne sur les 32 images grimpe. Les images déjà
entraînées continuent même de dériver quand on en entraîne d'autres (\num{2.85} au quatrième pas,
quand le premier lot revient).

\begin{table}[ht]
  \centering\small
  \caption{Divergence $\KL(p_\orig\|p_\theta)$ à $T=1$ pendant un entraînement \code{trust\_region}
  (32 images, lots de 8, taux \num{5e-5}, Apple M2). \enquote{Lot} : divergence du lot courant
  \emph{avant} sa mise à jour ; \enquote{toutes} : moyenne exacte sur les 32 images \emph{après}.}
  \label{tab:local}
  \begin{tabular}{@{}cccS[table-format=1.6]S[table-format=1.4]@{}}
    \toprule
    Pas & Lot & Déjà vu ? & {Lot (avant)} & {Toutes (après)}\\
    \midrule
    0 & 1 & non & 0.000000 & 0.0515\\
    1 & 2 & non & 0.000001 & 0.2589\\
    2 & 3 & non & 0.000025 & 0.4884\\
    3 & 4 & non & 0.000022 & 0.7286\\
    4 & 1 & oui & 2.848580 & 0.7493\\
    5 & 2 & oui & 0.138682 & 0.9860\\
    \bottomrule
  \end{tabular}
\end{table}

\begin{retenir}
Deux conséquences. \textbf{(1)} En essai réduit, la casse collatérale sur des images de la
distribution non ciblées est faible ; à pleine échelle, ce n'est plus le cas
(\cref{sec:pleine}). \textbf{(2)} La divergence du lot courant ne
renseigne pas sur celle des images déjà poussées : un critère d'arrêt doit mesurer un échantillon
\emph{fixe} d'images (\cref{sec:trust}).
\end{retenir}

\section{Méthodes}\label{sec:methodes}

Toutes les méthodes partagent une même interface, \code{unlearn(model, forget\_images, config)},
dans \chemin{src/hackathon/unlearning/methods.py}, et sont écrites avec PyTorch~\cite{paszke2019}.

\subsection{Références}
\begin{description}[leftmargin=1.2em,font=\sffamily\bfseries\color{bleu}]
  \item[identity] copie intacte ; vaut 0 par construction et sert à valider la chaîne.
  \item[scrub] SCRUB~\cite{kurmanji2023} adapté à l'absence d'étiquettes et de données de
    rétention : des pas qui \emph{maximisent} $\KL(p_\orig\|p_\theta)$ sur le forget set, plafonnée
    pour ne pas dégénérer en montée de gradient naïve, alternent avec des pas qui la
    \emph{minimisent} sur des images naturelles hors distribution
    (COCO~\cite{lin2014coco}, DIV2K~\cite{agustsson2017}, Kodak~\cite{kodak}).
  \item[gold\_soft] la méthode de départ de l'équipe : la cible d'une image est la moyenne des
    sorties de l'original sur plusieurs vues fortes, adoucie par une température ; l'élève
    l'apprend sur l'image exacte et sur une copie altérée, avec une ancre de rétention sur les
    images hors distribution.
\end{description}

\subsection{Détail des références et réglages par défaut}\label{sec:refsdetail}
\paragraph{\code{scrub}.} Chaque époque $e\le m$ (\code{scrub\_msteps}, 1 par défaut) commence par
une passe de \emph{max-steps} sur le forget set : pour un lot d'images exactes et une copie altérée
du même lot (une des quatre altérations simulées, tirée au hasard), on \emph{maximise}
$T^2\,\KL(p^{(T)}_\orig\|p^{(T)}_\theta)$ à $T=2$, écrêtée à $c$ (\code{scrub\_kl\_cap}, 2) : au-delà
du plafond le gradient est nul et la montée s'arrête. Chaque époque se termine par autant de
\emph{min-steps} qui minimisent la même divergence sur des lots d'images proxy (600 images COCO,
DIV2K et Kodak), pondérés par \code{retain\_weight}. Avec les défauts, une seule époque sur trois
contient des max-steps.

\paragraph{\code{gold\_soft}.} À chaque pas, la cible du lot est la moyenne de
$\operatorname{softmax}(z_\orig(\text{vue})/\tau)$ sur $V$ vues fortes (aire minimale
\code{target\_view\_min\_area}), recalculée à chaque pas (la V2 de \code{trust\_region} la calcule une seule fois,
\cref{sec:trustv2}). L'élève l'apprend sur l'image exacte et sur une copie altérée avec la perte
$\ell$ de \code{trust\_region} (moyenne des divergences à $T=1$ et $T=2$), plus l'ancre de rétention
de \code{scrub} sur les images proxy. Aucun arrêt : la quantité d'oubli dépend du taux et du nombre
d'époques.

\paragraph{Commun.} AdamW fusionné, sans décroissance des poids, gradient écrêté à norme 1, montée
linéaire du taux sur \SI{10}{\percent} des pas puis cosinus, bfloat16 sur GPU avec poids float32 ;
évaluation en float16 comme le serveur. Le \cref{tab:defauts} donne les réglages par défaut
(\chemin{src/hackathon/unlearning/configs.py}) ; les réglages effectivement utilisés par chaque run
sont dans le registre.

\begin{table}[ht]
  \centering\small
  \caption{Réglages par défaut (\code{UnlearnConfig}) et méthodes qui les lisent.}\label{tab:defauts}
  \begin{tabularx}{\linewidth}{@{}>{\raggedright}p{4.3cm}>{\raggedright}p{2.4cm}X@{}}
    \toprule
    Réglage & Défaut & Méthodes\\
    \midrule
    \code{batch\_size}, \code{epochs}, \code{lr} & 32, 3, \num{e-5} & toutes\\
    \code{warmup\_fraction}, \code{weight\_decay} & \num{0.1}, 0 & toutes\\
    \code{trainable\_blocks} & $-1$ (tout le réseau) & toutes\\
    \code{kd\_temperature} & 2 & ancre de rétention, perte $\ell$, \code{scrub}, \code{dumb}\\
    \code{retain\_count}, \code{retain\_weight} & 600, 1 & \code{scrub}, \code{gold\_soft} ;
      \code{retain\_weight} : \code{dumb}\\
    \code{dumb\_target\_loss} & \num{0.45} & \code{dumb} (V\_Dumb)\\
    \code{altered\_weight} & 1 & \code{scrub}, \code{gold\_soft}\\
    \code{scrub\_msteps}, \code{scrub\_kl\_cap} & 1, 2 & \code{scrub}\\
    \code{target\_temperature} & 2 & \code{gold\_soft}, \code{trust\_region} (specs : 6 en V1, 4 en V2,
      6 pour la couverture de la V3)\\
    \code{target\_views}, \code{target\_view\_min\_area} & 3, \num{0.35} & \code{gold\_soft}, V2, V3\\
    \code{kl\_radius} & \num{0.5} & \code{trust\_region}\\
    \code{trust\_target} & \code{tempered} (V1) ; \code{views} (V2), \code{calibrated} (V3)
      & \code{trust\_region}\\
    \code{radius\_check\_every} & 10 (campagne, V2 et V3 : 1) & \code{trust\_region}\\
    \code{trust\_lambda}, \code{trust\_hedge} & \num{0.25}, \num{0.25} & V3\\
    \code{radius\_fraction} & aucun : \code{kl\_radius} s'applique (V3 : \num{0.9}) & \code{trust\_region}\\
    \code{altered\_fraction} & 0 (V3 : 0 ou \num{0.5}) & \code{trust\_region}\\
    \code{anchor\_folds}, \code{anchor\_weight}, \code{anchor\_views} & 0 (sans ancre), 1, 0
      & \code{trust\_region} (V4)\\
    \code{radius\_dual\_lr} & 0 (arrêt au premier franchissement) & \code{trust\_region} (V4)\\
    \code{retain\_source} & \code{proxy} (\code{forget\_views} : conforme au sujet)
      & \code{scrub}, \code{gold\_soft}\\
    \code{holdout\_fraction} & 0 (runs de réglage : \num{0.2}) & toutes, via \code{run}\\
    \bottomrule
  \end{tabularx}
\end{table}

\subsection{Région de confiance en divergence (\code{trust\_region})}\label{sec:trust}
\paragraph{Idée.} SCRUB+R~\cite{kurmanji2023} revient en arrière jusqu'au point où l'erreur sur le
forget set égale celle d'un jeu de validation : il n'oublie pas \enquote{trop}. Sans jeu de
validation, ce point devient un \emph{rayon} $\varepsilon$, un budget d'oubli, comme une région de
confiance en optimisation~\cite{schulman2015} :
\begin{equation}\label{eq:trust}
  \min_\theta\ \Exp_{x\in F}\ \ell(\theta;x)
  \quad\text{sous}\quad
  K(\theta)=\Exp_{x\in F}\KL\bigl(p^{(1)}_\orig(x)\,\big\|\,p^{(1)}_\theta(x)\bigr)\le\varepsilon,
\end{equation}
où $p^{(1)}$ désigne le softmax à température 1. Comme l'original est quasi certain, $K(\theta)$
est à peu près la perte moyenne de l'élève sur les pseudo-étiquettes : $\varepsilon$ s'interprète
comme la perte moyenne visée, celle qu'aurait un gold sur des images inédites.

\paragraph{Perte.} L'élève est tiré vers la sortie de l'original adoucie,
$q_\tau(x)=\operatorname{softmax}(z_\orig(x)/\tau)$ avec $\tau=6$, une distillation à
température~\cite{hinton2015} : $\ell(\theta;x)=\tfrac12\bigl[\KL(q_\tau\|p^{(1)}_\theta)+
4\,\KL(q_\tau^{(2)}\|p^{(2)}_\theta)\bigr]$, où $q_\tau^{(2)}$ est la cible portée à température 2 :
on vise à la fois la température 1 des pertes du score et la température 2 de ses distributions. Cette perte est convexe en les logits de l'élève, bornée, et
préserve le classement des classes ; elle ne contient aucune montée de gradient, que le sujet
pénalise et dont l'effondrement est documenté~\cite{zhang2024npo}. Elle n'utilise que les images
exactes du forget set : aucune augmentation, aucune image externe.

\paragraph{Arrêt.} Tous les $k$ pas, $K(\theta)$ est mesurée exactement sur un échantillon fixe de
256 images entraînées, et l'entraînement s'arrête dès qu'elle atteint $\varepsilon$ ; une mesure
sur toutes les images est faite à la fin. Le premier choix, $k=10$, s'est révélé trop grossier à
pleine échelle ; la campagne utilise $k=1$ (réglage \code{radius\_check\_every},
\cref{sec:campagne}). La divergence du lot courant ne peut pas servir
(\cref{sec:local}). Avant d'entraîner, on vérifie que la cible est au-delà du rayon :
$\KL(p_\orig\|q_\tau)=\num{1.20}$ pour $\tau=6$ sur l'échantillon d'essai, ce qui rend atteignables
tous les $\varepsilon\le\num{0.8}$ balayés.

\paragraph{Lots stratifiés.} Les images sont regroupées par pseudo-étiquette (l'argmax de
l'original), puis le lot $i$ prend une image sur $n$ dans cet ordre : chaque lot couvre les classes
dans les proportions du forget set, et aucun pas ne déplace la tête du classifieur vers quelques
classes.

\paragraph{Réglages.} AdamW~\cite{loshchilov2019} (taux \num{e-5}, \SI{10}{\percent} de montée puis
cosinus, gradient écrêté à 1), lots de 32, bfloat16 sur GPU avec poids float32, 10 époques au plus.
Le rayon a été balayé de \num{0.01} à \num{0.8} ; $\varepsilon=\num{0.15}$ est retenu
(\cref{sec:campagne}).

\subsection{V2 : une cible propre à chaque image (\code{trust\_target=views})}\label{sec:trustv2}
\paragraph{Pourquoi.} La V1 adoucit toutes les images de la même façon. Or le gold est hétérogène :
sur des images inédites, il se trompe sur une partie d'entre elles (sa précision de test, de
l'ordre de 78 à \SI{92}{\percent} selon la \cref{sec:calibration}), et vers une classe précise. Une
perte moyenne juste ne suffit pas : $S_{\text{ms}}$ compare aussi l'écart-type des pertes, et
$S_{\text{pred}}$ la distribution complète, image par image (\cref{eq:ms,eq:pred}).

\paragraph{Cible.} Pour chaque image $x$, on tire $V$ vues fortes $t_1(x),\dots,t_V(x)$
(recadrage d'aire $a$, puis couleur, \cref{sec:calibration}) et l'on pose
\begin{equation}\label{eq:qviews}
  q_V(x)=\frac1V\sum_{v=1}^{V}\operatorname{softmax}\bigl(z_\orig(t_v(x))/\tau\bigr),
\end{equation}
calculée une seule fois avant l'entraînement, sans gradient. C'est la cible de \code{gold\_soft},
employée dans la région de confiance : même perte $\ell$ avec $q_V$ à la place de $q_\tau$, même
contrainte \eqref{eq:trust}, même arrêt au rayon, mêmes lots stratifiés, toujours sur les seules
images exactes. La moyenne sur les vues tient lieu de moyenne de plusieurs gold (le sujet moyenne
leurs distributions) ; la température éloigne la cible pour que le rayon reste atteignable. La
cible donne la direction et l'hétérogénéité, $\varepsilon$ la quantité. L'arrêt anticipé ne gomme
pas cette hétérogénéité : le gradient de $\KL(q\|p^{(1)}_\theta)$ par rapport aux logits vaut
$p^{(1)}_\theta-q$, donc les images dont la cible s'écarte le plus de l'original sont poussées le
plus fort à chaque pas.

\paragraph{Faisabilité et choix de $\tau$.} La \cref{tab:cible} donne la divergence de la cible
mesurée sur le vrai modèle. À $\tau=1$, la cible par vues reste très proche de l'original :
moyenner trois vues remonte la probabilité de la pseudo-étiquette, car $-\log$ de la moyenne est
inférieur à la moyenne des $-\log$. Seul $\varepsilon=0{,}2$ serait atteignable, et le contrôle de
faisabilité refuserait les autres rayons. Monter $\tau$ éloigne la cible, mais trop l'adoucir
efface ce qu'elle apporte : à $\tau=6$, l'hétérogénéité relative de ses pertes (écart-type sur
moyenne) tombe à \num{0.23}, sous celle de la V1 (\num{0.28}). Nous retenons $\tau=4$ : la cible
est à une distance comparable à celle de la V1 (\num{1.08} contre \num{1.22}), au-delà des quatre
rayons de la V1, que nous reprenons pour comparer à budget égal, et son hétérogénéité relative
vaut \num{0.42}. Aire $a=0{,}45$ : à $0{,}55$, la cible est juste sur \SI{98}{\percent} des images,
bien plus qu'un gold plausible. Même à $0{,}45$, elle l'est sur \SI{96.5}{\percent} : trois vues
moyennées se trompent moins qu'une seule, une limite de cette cible. La campagne à pleine échelle
(\cref{sec:campagnev2}) a depuis montré que le rayon utile est petit ($\varepsilon\approx\num{0.15}$),
ce qui rend $\tau=1$ atteignable, et que $\tau=1$ fait nettement mieux que $\tau=4$.

\begin{table}[ht]
  \centering\small
  \caption{Cibles de \code{trust\_region} sur le vrai modèle : 256 images du forget set, Apple M2,
  3 vues d'aire $[0{,}45;0{,}55]$ (sauf mention). KL : $\KL(p_\orig\|q)$ à $T=1$ ; accord : part
  des images où l'argmax de la cible est celui de l'original ; perte : $-\log q_y$ sur cette
  pseudo-étiquette $y$, moyenne et écart-type sur les images.}\label{tab:cible}
  \sisetup{table-format=1.3}
  \begin{tabular}{@{}l SSSS@{}}
    \toprule
    Cible & {KL} & {Accord} & {Perte moy.} & {Perte é.-t.}\\
    \midrule
    V1, $q_\tau$, $\tau=6$ & 1.218 & 1.000 & 1.222 & 0.345\\
    V2, $\tau=1$ & 0.257 & 0.945 & 0.256 & 0.319\\
    V2, $\tau=2$ & 0.375 & 0.957 & 0.377 & 0.332\\
    V2, $\tau=3$ & 0.661 & 0.965 & 0.663 & 0.399\\
    \textbf{V2, $\tau=4$ (retenue)} & 1.079 & 0.965 & 1.082 & 0.455\\
    V2, $\tau=6$ & 1.911 & 0.965 & 1.914 & 0.436\\
    V2, $\tau=1$, aire $[0{,}55;0{,}65]$ & 0.147 & 0.980 & 0.146 & 0.234\\
    \bottomrule
  \end{tabular}
\end{table}

\paragraph{Un biais à garder en tête.} La doublure du gold du banc est \emph{aussi} l'original sur
des vues fortes (\cref{sec:banc}). Pour limiter ce recouvrement, les vues de la V2 sont tirées
d'une graine issue de celle du run, jamais de celle du banc, et à une aire différente ($[0{,}45;
0{,}55]$ contre $[0{,}40;0{,}50]$, $[0{,}50;0{,}60]$ et $[0{,}60;0{,}70]$). Le générateur de vues
reste le même : un avantage local de la V2 sur la V1 est en partie mécanique et ne prouve pas un
gain au score réel. Seul compte un écart qui tient sous la pire doublure et dépasse le bruit
d'entraînement (\cref{sec:banc}).

\subsection{V3 : une cible pour chacun des deux objets du score (\code{trust\_target=calibrated})}\label{sec:trustv3}
\paragraph{Pourquoi.} Le score compare l'élève à deux objets différents du gold
(\cref{sec:score}) : $S_{\text{ms}}$ à des statistiques de pertes calculées modèle par modèle puis
moyennées, $S_{\text{pred}}$ au \emph{mélange} des distributions des modèles. Notons $y$ la
pseudo-étiquette (argmax de l'original), $\ell_k(x)$ la perte à $T=1$ du modèle $k$ du gold et
$\bar\ell(x)=\Exp_k\ell_k(x)$. Par linéarité, la perte moyenne visée est $\Exp_x\bar\ell(x)$ ;
mais l'écart-type visé, $\Exp_k\operatorname{std}_x\ell_k$, dépasse $\operatorname{std}_x\bar\ell$
(inégalité triangulaire), et un élève qui copie le mélange $\bar p$ a une perte plus basse encore
que $\bar\ell$ (inégalité de Jensen). Une note de recherche interne l'a chiffré (256 images,
original seul, Apple M2, doublure par recadrage d'aire $0{,}50$) : copier le mélange de 4 vues donne
une perte de \num{0.081} $\pm$ \num{0.260}, contre \num{0.496} $\pm$ \num{1.556} pour les
statistiques par vue, six fois trop bas sur les deux moments. Le second moment pèse lourd : comme
l'original a une perte quasi nulle, $d_\orig$ vaut à peu près la moyenne plus l'écart-type visés,
et un élève dont la moyenne est juste mais dont l'écart-type vaut $c$ fois sa moyenne vérifie
\begin{equation}\label{eq:msbound}
  S_{\text{ms}}\le 1-\frac{\mathrm{CV}_\gold-c}{1+\mathrm{CV}_\gold},\qquad
  \mathrm{CV}_\gold=\frac{\Exp_k\operatorname{std}_x\ell_k}{\Exp_x\bar\ell}\quad(c<\mathrm{CV}_\gold).
\end{equation}
Les six doublures du banc ont $\mathrm{CV}_\gold$ entre \num{2.3} et \num{3.8}
(\cref{tab:familles} ; \num{1.7} pour un masquage de \SI{40}{\percent}) : un modèle sur-confiant se trompe avec assurance~\cite{guo2017}, d'où une
queue de pertes lourde. Les cibles V1 et V2 ont $c=\num{0.28}$ et \num{0.42} (\cref{tab:cible}) :
sous la doublure d'aire $0{,}50$ ($\mathrm{CV}_\gold=\num{3.1}$), \eqref{eq:msbound} plafonne
$S_{\text{ms}}$ à \num{0.31} et \num{0.35}. Ce plafond vaut pour un élève qui copierait sa cible ;
le réseau ne le fait pas : à pleine échelle, la V1 atteint $S_{\text{ms}}=\num{0.456}$ sous
$s_{60}$ (\cref{tab:meilleure}), au-delà du plafond de \num{0.27} que \eqref{eq:msbound} donne pour
sa cible sous cette doublure ($\mathrm{CV}_\gold=\num{3.77}$). L'argument de la V3 sur
$S_{\text{ms}}$ en sort affaibli ; son argument sur $S_{\text{pred}}$, la composante qui limite tous
les runs mesurés (\cref{sec:meilleure}), reste entier.

\paragraph{Cible.} Une seule fois, sans gradient, sur $V$ vues fortes $t_v(x)$ (les tirages de la
V2) et avec $T=2$ la température du score :
\begin{gather}
  \hat p(x)=\frac1V\sum_{v=1}^{V}\operatorname{softmax}\bigl(z_\orig(t_v(x))/2\bigr),\qquad
  \bar\ell(x)=\frac1V\sum_{v=1}^{V}\operatorname{CE}\bigl(z_\orig(t_v(x)),y\bigr),\label{eq:phat}\\
  p_{\text{dir}}=(1-b)\,\hat p+b\,\operatorname{softmax}(z_\orig/\tau),\qquad
  z=2\log p_{\text{dir}},\qquad \ell_{\text{mix}}=\operatorname{CE}(z,y),\label{eq:pdir}\\
  \ell^\star=\ell_{\text{mix}}+\lambda\,\bigl(\bar\ell-\ell_{\text{mix}}\bigr).\label{eq:lstar}
\end{gather}
Seul le logit de $y$ change ensuite. Avec $A=\log\sum_{c\ne y}e^{z_c}$, on a
$\operatorname{CE}(z,y)=\log\bigl(1+e^{A-z_y}\bigr)$, d'où une forme close, sans recherche
numérique :
\begin{equation}\label{eq:zy}
  z_y\leftarrow A-\log\bigl(e^{\ell^\star}-1\bigr)\quad\Longrightarrow\quad
  \operatorname{CE}(z,y)=\ell^\star\ \text{exactement}
\end{equation}
($\ell^\star$ est planché à \num{e-6}). Les autres logits restent $2\log p_{\text{dir}}$ : à $T=2$,
les rapports entre classes fausses sont exactement ceux de $p_{\text{dir}}$. Tout le reste est
celui de la V1 : même perte $\ell$, dont le terme à $T=1$ porte désormais la perte $\ell^\star$ visée
par $S_{\text{ms}}$ et le terme à $T=2$ la direction visée par $S_{\text{pred}}$, même arrêt, mêmes
lots, mêmes images exactes. Le rayon est tiré de la cible :
\begin{equation}\label{eq:epsv3}
  \varepsilon=\rho\,\Exp_{x\in F}\KL\bigl(p^{(1)}_\orig(x)\,\big\|\,q(x)\bigr)\approx\rho\,\Exp_x\ell^\star(x),
  \qquad \rho=\text{\code{radius\_fraction}}=0{,}9,
\end{equation}
car l'original est quasi certain : l'élève qui atteint la cible a $K(\theta)\approx\Exp_x\ell^\star$,
et $\rho<1$ garde la cible au-delà du rayon.

\paragraph{Justification.}
\begin{itemize}[leftmargin=1.2em]
  \item \textbf{Direction : le mélange, couvert.} $\KL(p\|q)$ est une divergence de Bregman ; pour
    une loi $P$ incertaine, $\arg\min_q\mathbb{E}\,\KL(P\|q)=\mathbb{E}\,P$~\cite{banerjee2005bregman} :
    face à des gold incertains, la meilleure prédiction est leur mélange, d'où $\hat p$. Quand la
    direction des vues est fausse, $\KL(p_\gold\|\cdot)$ récompense de répartir la masse : la
    couverture $b$ par l'original adouci (la cible de la V1) achète cette robustesse.
  \item \textbf{Le seul logit vrai.} Augmenter le logit de $y$ de $\delta$ (à $T=2$) coûte en
    divergence $\log\bigl(1-\hat p_y+\hat p_ye^\delta\bigr)-\hat p_y\delta$, de pente nulle en 0 et de
    courbure $\hat p_y(1-\hat p_y)$ : un coût du second ordre, alors que la perte à $T=1$ bouge au
    premier ordre. Le produit $S_{\text{ms}}S_{\text{pred}}$ a donc un optimum intérieur en
    $\lambda$, que la note de recherche situe vers $0{,}25$ à $0{,}5$ dans son modèle (ci-dessous).
    La classe fausse préférée de la cible ne bouge pas : pas de lissage d'étiquettes, qui efface la
    ressemblance entre classes~\cite{muller2019} dont $S_{\text{pred}}$ dépend ; c'est une
    distillation à température~\cite{hinton2015} dont la cible est construite pour le score.
  \item \textbf{Les vues comme estimation de \enquote{hors entraînement}.} La précision attendue
    d'un modèle sur $x$ quand $x$ n'est pas dans son entraînement (le C-score) exige de nombreux
    réentraînements~\cite{jiang2021cscore} ; les vues fortes en sont un substitut bon marché, non
    validé par ces travaux. $\bar\ell(x)$ joue le rôle d'une difficulté par image, dont dépend la
    difficulté de l'oubli~\cite{zhao2024hard} : les fortes pertes vont aux images où les vues
    doutent le plus.
  \item \textbf{Force des vues.} Si l'on ignore la précision du gold, le choix minimax de la
    moyenne visée sur $[\mu_{\text{lo}},\mu_{\text{hi}}]$ est la moyenne harmonique
    $2\mu_{\text{lo}}\mu_{\text{hi}}/(\mu_{\text{lo}}+\mu_{\text{hi}})$ ; avec les doublures
    extrêmes par recadrage (\cref{tab:familles}),
    $2\times\num{0.954}\times\num{0.333}/\num{1.287}=\num{0.494}$, la perte de la force centrale :
    d'où l'aire $a=0{,}50$.
\end{itemize}

\paragraph{Copies douces (\code{altered\_fraction}).} Les copies altérées cachées pèsent la moitié
de chaque terme du score. En essai réduit, l'oubli était très local (\cref{sec:local}) et ne passait
qu'à moitié aux copies ; à pleine échelle, la V1 transfère déjà largement (la pire copie bouge
\num{0.70} fois autant que l'image exacte, \cref{sec:campagne}). Les copies douces visent ce reste :
entraîner sur un voisinage de $x$, pas sur le seul point $x$. Les attaques d'appartenance par étiquette seule
exploitent justement la stabilité de la prédiction sous perturbation~\cite{choquettechoo2021}. En
distillation, professeur et élève gagnent à voir la même vue~\cite{beyer2022} ; ici la cible de la
source est invariante, donc cohérente par construction. Dans chaque lot, une part
\code{altered\_fraction} des images est remplacée, moitié par son retournement, moitié par un
recadrage de \SIrange{70}{100}{\percent} de son aire, avec la cible de l'image source. Le jitter
et le JPEG ne sont jamais vus : sur le banc, leurs groupes mesurent le transfert à des familles
inédites, un contrôle non circulaire. Si l'élève ne transfère qu'une fraction $r$ de son mouvement
aux copies, le terme d'oubli vaut grossièrement $\bigl((1+r)/2\bigr)^2S_{\text{ms}}S_{\text{pred}}$.

\paragraph{Réglages et plafonds.} $\lambda\in\{0{,}25;0{,}5\}$ croisé avec
$\code{altered\_fraction}\in\{0;0{,}5\}$ (4 variantes, un nœud), $b=0{,}25$, $\tau=6$, $V=4$,
$a=0{,}50$, $\rho=0{,}9$, rayon testé à chaque pas comme l'impose la campagne
(\cref{sec:campagne}) ; taux, époques et lots de la V1, pour comparer toutes choses égales. La
note de recherche donne des \emph{plafonds}, dans un modèle simplifié à logits libres (chaque image
garde ses propres logits, qui suivent le flot de gradient de $\ell$ jusqu'au rayon ; ce n'est pas
un réseau), sur le seul groupe des images exactes : sous la pire des six doublures du banc
(\cref{sec:banc}), \num{0.086} pour cette cible ($\lambda=b=0{,}25$) contre \numrange{0.030}{0.037}
pour la V1 et la V2 ; pour les copies, passer de $r\approx\num{0.70}$ (V1 à pleine échelle) à
$r\approx1$ multiplierait l'oubli par \num{1.38} au plus. Ce ne sont pas des
attentes : un réseau, un rayon non atteint ou un transfert partiel feront moins, et seul le banc à
pleine échelle tranchera.

\paragraph{Mesure de la cible.} Sur l'essai réduit (52 images, 4 vues, $\lambda=b=0{,}25$) : divergence
$\KL(p_\orig\|q)=\num{0.173}$, accord de l'argmax \num{0.962}, $\ell^\star=\num{0.173}\pm\num{0.482}$,
soit un rapport écart-type sur moyenne de \num{2.8}, contre \num{0.28} et \num{0.42} pour les cibles V1
et V2 (\cref{tab:cible}, 256 images) : la cible porte l'hétérogénéité des doublures.

\paragraph{Confrontation à la campagne.} La note de recherche a été écrite sans connaître la
campagne à pleine échelle (\cref{sec:campagne}). Elle prédisait qu'avec la cible de la V1,
\code{score\_min} chuterait au-delà d'un rayon de \numrange{0.2}{0.3}, par la composante des
prédictions sous la doublure la plus sûre d'elle : la campagne le mesure (\num{0.034} à
$\varepsilon=\num{0.3}$, \num{0.002} à \num{0.5}, 0 à \num{0.8}, $s_{60}$ étant chaque fois la pire
doublure, \cref{tab:rayon}). Le rayon que la V3 tire de sa cible (\num{0.9} fois sa divergence, soit
\numrange{0.16}{0.18} selon l'échantillon mesuré) tombe dans le plateau de la V1
(\numrange{0.1}{0.2}). La campagne affaiblit en revanche l'argument sur $S_{\text{ms}}$ (ci-dessus),
et montre qu'à divergence égale la direction d'une cible \emph{tempérée} ne change rien de
mesurable ($\tau$ de 4 à 12) ; la V3 parie sur autre chose, la forme de la distribution des pertes
image par image, que $\tau$ ne change pas. Elle ne traite pas non plus le défaut mesuré de la V1, la
localité (\cref{sec:pistes}). Seuls ses runs à pleine échelle trancheront, notés en
\code{bench-v4} et comparés à des références recalculées sous ce protocole. Ils ont tranché : à son
meilleur réglage, la V3 fait jeu égal avec la V2, et le mécanisme par $S_{\text{ms}}$ n'est pas
observé (\cref{sec:campagnev3}).

\paragraph{Un pronostic réfuté.} Écrite sans connaître la campagne V3, la relecture des 54 runs de
la V1 (\cref{sec:diagnostic}) pronostiquait pour la V3 le plateau de la V1 (\numrange{0.060}{0.067}
en équivalent \code{bench-v3}) : au rayon, l'élève ne suit pas la structure image par image de sa
cible. La campagne l'a réfuté sur le score : la V3 retenue fait \num{0.0798 +- 0.0012}, jeu égal avec
la V2 et nettement au-dessus de la V1, par la direction de sa cible. Elle l'a confirmé sur le
mécanisme : la dispersion des pertes est fixée par le réseau, pas par la cible, et le fold d'audit
bouge autant qu'avec la V1 (\cref{sec:diagnostic}).

\subsection{V4 : une ancre tirée du forget set}\label{sec:trustv4}
\paragraph{Pourquoi.} Environ \SI{88}{\percent} du mouvement de la V1 est commun à toutes les images
mémorisées, entraînées ou non (\cref{sec:diagnostic}). La perte de V1--V3 ne contient aucun terme
qui le freine : rien ne lui coûte de déplacer une image mémorisée qu'elle ne vise pas. La
littérature le freine par un terme sur des données de rétention : la KL vers l'original de
SCRUB~\cite{kurmanji2023}, une contrainte sur la perte de rétention résolue en
primal-dual~\cite{entesari2025constrained} ou une géométrie tirée du retain~\cite{huang2024unified},
la phase de réparation de toutes les méthodes analysées du concours NeurIPS~\cite{triantafillou2024}.
SCAR remplace le retain par des images hors distribution~\cite{bonato2024scar}, comme notre
\code{scrub}, et le sujet l'exclut pour une méthode soumise~\cite{sujet2026}. Cha et al. tirent du
forget set des exemples adverses comme proxys~\cite{cha2024learning}, mais ce ne sont pas des images
mémorisées, celles que le mouvement commun déplace. La revue de l'équipe (22 novembre) n'a trouvé
aucun travail qui contraigne la rétention à partir du seul forget set.

\paragraph{Ancre croisée (\code{anchor\_folds}).} Les images entraînées sont coupées en $K$ plis
$F_1,\dots,F_K$ stratifiés par pseudo-étiquette, oubliés l'un après l'autre. Pendant la phase $k$,
le pli $F_k$ est tiré vers sa cible (celle de la V3 ; les spécifications prennent la V3 retenue) jusqu'à son propre
rayon, et chaque pas ajoute un lot d'ancre tiré des autres plis, tenus à des sorties fixées :
\begin{equation}\label{eq:v4}
  \min_\theta\ \Exp_{x\in F_k}\ell(\theta;x)\;+\;w\,\Exp_{x\in F\setminus F_k}
  4\,\KL\bigl(p^{(2)}_{\mathrm{ref}(x)}(x)\,\big\|\,p^{(2)}_\theta(x)\bigr)
  \quad\text{jusqu'à}\quad \Exp_{x\in P_k}\KL\bigl(p^{(1)}_\orig(x)\,\|\,p^{(1)}_\theta(x)\bigr)\ge\varepsilon_k,
\end{equation}
où la référence $\mathrm{ref}(x)$ est l'original pour les images jamais poussées ($F_{k+1},\dots,F_K$)
et une copie figée de l'élève prise au début de la phase pour les images déjà oubliées
($F_1,\dots,F_{k-1}$) ; $w$ = \code{anchor\_weight}, $P_k$ la sonde de la phase, $\varepsilon_k$ son
rayon. Tenir les plis jamais poussés à l'original, et non à leur état du début de phase, rappelle
aussi la dérive héritée des phases précédentes. La divergence à $T=2$ est celle du retain du score
\eqref{eq:retain} ; les sorties de référence sont calculées sous la même précision que l'élève, de
sorte que l'ancre vaut exactement 0 avant le premier pas. Une part \code{anchor\_views} du lot d'ancre
est montrée en vue forte (recadrage de \SIrange{25}{35}{\percent} de l'aire, puis couleur), pour
tenir aussi la réponse sur des entrées que l'original n'a pas mémorisées ; cette aire n'en recouvre
aucune du banc (doublures de \num{0.40} à \num{0.70}, vue d'audit \num{0.50}), mais la couleur est
proche du jitter du banc : pour ces variantes, la famille des masquages tranche. Chaque phase a son
optimiseur, son calendrier de taux et ses flux de tirages : la longueur d'une phase ne décale pas la
suivante.

\paragraph{Pourquoi elle devrait localiser.} Les plis jamais poussés sont des images mémorisées de
la même distribution que le control set : le meilleur proxy de rétention disponible. Un mouvement
commun à toutes les images mémorisées les déplace, et l'ancre le fait payer ; seul le mouvement
propre aux images poussées reste gratuit. Qu'un tel mouvement existe est plausible : la mémorisation
d'un exemple tient à peu de neurones~\cite{maini2023localized} ; en essai réduit, les mises à jour de
la V1 sur 8 images ne bougeaient pas les autres (\cref{sec:local}), mais à pleine échelle la campagne
mesure l'inverse, et seuls les runs diront si l'ancre le retrouve. Le run mesure lui-même sa fuite :
à la fin de chaque phase $k<K$, la divergence à $T=2$ des plis jamais poussés vers l'original
(\code{pending\_kl\_p}$k$, la dernière dans \code{final.pending\_kl}), l'analogue du fold d'audit,
disponible même dans un run final qui n'en a pas.

\paragraph{Portée et limites.} \textbf{Le gain visé est plafonné.} La V4 s'applique à la V3 retenue
(\cref{sec:campagnev3}) : à la graine 0, oubli \num{0.103} $\times$ retain \num{0.758}, avec un accord de
\num{0.943} sur le fold d'audit. Une dérive d'audit nulle donnerait au mieux \num{0.097}
($+\SI{24}{\percent}$ ; \num{0.102} à la graine 1), au-dessus du seuil d'amélioration (\num{0.086}) et
au niveau de \code{scrub} dosé ; un gain au-delà devra venir du terme d'oubli. Sur la V1 le même
plafond n'était que de \num{0.079}. \textbf{Il est en partie
mécanique} : l'ancre tient exactement la population que mesure le retain du banc (images
mémorisées non entraînées) ; il ne vaut pour le score officiel que si le control set est fait
d'images d'entraînement mémorisées, ce que le sujet ne précise pas. \textbf{La dernière phase} n'a
dans son ancre que des images déjà oubliées, qui ne freinent le mouvement commun que s'il les
déplace aussi ; $K=4$ garde des plis jamais poussés dans l'ancre pendant trois phases sur quatre.
\textbf{Coût} : par pas, une passe avant-arrière de plus et une passe avant sans gradient ; en
mémoire, une copie de l'original et, dès la phase 2, une copie figée (\SI{344}{\mega\byte} chacune).

\paragraph{Contrainte active (\code{radius\_dual\_lr}).} Les arrêts de la V1 tombent tous autour de
la fin du warmup, au taux maximal (\cref{sec:diagnostic}) : le modèle est figé en plein vol, au
premier franchissement du rayon. Avec \code{radius\_dual\_lr}~$=\eta>0$, la phase parcourt tout son
calendrier et le rayon devient une contrainte active, à la manière d'un
primal-dual~\cite{entesari2025constrained} : la perte du lot reçoit
$\mu\,\KL(p^{(1)}_\orig\|p^{(1)}_\theta)$, et à chaque test de la sonde
$\mu\leftarrow\max\bigl(0,\mu+\eta\log(\hat K/\varepsilon)\bigr)$, un pas symétrique : un
dépassement pendant le warmup se rattrape aussi vite qu'il s'est payé. L'élève peut continuer
d'apprendre la forme de sa cible sans dépasser le budget d'oubli. Prédiction, non mesurée : avec
une cible peu dispersée (V1), converger vers elle égaliserait les pertes et ferait baisser
$S_{\text{ms}}$ ; le mode n'a de chance qu'avec la cible de la V3, et sans doute à $\lambda=\num{0.5}$
(l'élève générique de la V1 a déjà des pertes plus dispersées que la cible de la V3 à $\lambda=\num{0.25}$).
Chez Entesari et al., la contrainte porte sur la rétention ; ici, faute de retain, elle porte sur le
budget d'oubli.

\paragraph{SCRUB conforme (\code{retain\_source=forget\_views}).} À perte moyenne égale, \code{scrub}
bat la V1 (\cref{tab:diagnostic}, avec la réserve sur les copies). Mais son ancre utilise des images
externes. L'option \code{forget\_views} remplace ces images par des vues fortes du forget set
lui-même (même aire que l'ancre de la V4), tenues aux sorties de l'original : l'ascension plafonnée
efface la pointe de mémorisation sur l'image exacte, l'ancre garde la fonction autour d'elle.
C'est le test le moins cher et le plus décisif de la campagne proposée : il dit si la méthode en
tête le reste une fois conforme au sujet. \code{gold\_soft} accepte la même option.

\paragraph{Relecture.} Le code de la V4 a été relu par un agent indépendant, sans entraînement :
douze défauts relevés, tous corrigés avant la campagne (fuite mesurée contre le mauvais modèle,
montée duale asymétrique, indicateurs de rayon mal définis par phase, réglages texte fusionnés dans
le tableau comparatif, plis plus nombreux que les images, précision de l'ancre, flux de tirages
partagés entre phases, aire des vues qui recouvrait la doublure $s_{40}$). Son verdict, que nous
reprenons : l'ancre croisée est un test de mécanisme légitime, pas le pari principal ; c'est le
terme d'oubli qui bloque, d'où l'ordre de la campagne (\cref{sec:plan}).

\paragraph{Rayon et sonde.} Le rayon \code{radius\_fraction} est calculé sur toutes les images,
alors que l'arrêt est testé sur la sonde de 256 : avec des pertes de cible aussi dispersées que
celles de la V3, un élève qui suivrait exactement sa cible pourrait ne jamais s'arrêter (piste,
non vérifiée, pour le \enquote{rayon jamais atteint} de la V3 à 1 ou 2 vues). Les options de la V4
prennent le rayon sur la sonde ; la V1, la V2 et la V3 gardent leur calcul, pour que tous les runs
mesurés se rejouent exactement (\code{ref4-*}, vérifié sur le petit modèle de test).

\paragraph{État.} Codée et couverte par les tests unitaires ; un essai intermédiaire s'était arrêté
sur une erreur propre au GPU d'Apple (pas de float64, corrigée). Mesurée depuis à pleine échelle,
aux seuils posés avant les runs (\cref{sec:campagnev4}) : ancre croisée et contrainte active
réfutées, SCRUB conforme confirmé une fois redosé.

\subsection{V5 : oubli en bande, dosé hors pli (\code{band\_ascent})}\label{sec:v5}
\paragraph{Pourquoi.} Les V1 à V4 de \code{trust\_region} descendent vers une cible, et environ
\SI{88}{\percent} du mouvement de la V1 est commun à toutes les images mémorisées
(\cref{sec:diagnostic}). Un calcul dans le régime linéarisé~\cite{jacot2018ntk,ortizjimenez2023tangent}
dit pourquoi, et ce qu'il faudrait changer. Notons $g_i=\partial\mathcal L/\partial z(x_i)$ la force que
la perte applique aux logits de l'image entraînée $x_i$, et $k(x,x')$ le noyau tangent (préconditionné
par Adam). Sur une passe de taux $\eta$, une image $x$ bouge de $-\eta\sum_ik(x,x_i)\,g_i$. Dans un
modèle où les images d'une même pseudo-classe partagent un mode commun,
$k(x_i,x_j)\approx c\,\mathbf 1[\hat y_i=\hat y_j]+d\,\mathbf 1[i=j]$,
\begin{equation}\label{eq:v5lin}
  \Delta z(x)\approx-\eta\,c\,n_k\,\bar g_k,\qquad
  \Delta z(x_i)\approx-\eta\,\bigl(c\,n_k\,\bar g_k+d\,g_i\bigr),\qquad
  \bar g_k=\frac1{n_k}\sum_{j:\,\hat y_j=k}g_j,
\end{equation}
pour une image $x$ non vue et une image entraînée $x_i$ de la classe $k$, qui en compte $n_k$ parmi
les images entraînées. La fuite suit la force \emph{moyenne} de la classe, la part propre la force
de l'image. Une descente vers une cible a $g_i=p^{(1)}_\theta(x_i)-q(x_i)$ d'ordre 1 et de même sens
(\enquote{moins sûr}) pour toutes les images : $\bar g_k\approx g_i$, et la part commune du mouvement
vaut $c\,n_k/d$ fois sa part propre, ce qui est grand dès que $n_k\approx28$ (2800 images, 100 classes).
Le mode moyen, de plus grande valeur propre, est en outre ajusté le premier ; un arrêt précoce au rayon
le capture avant tout le reste, ce qui rend compte de l'insensibilité mesurée à la cible ($\tau$ de 4 à
12, \cref{sec:diagnostic}). Pour un changement atteint $r_i=\Delta z(x_i)$, la divergence étant
localement quadratique,
\begin{equation}\label{eq:v5ratio}
  \frac{\Exp_{x\notin F}\KL\bigl(p^{(2)}_\orig(x)\,\|\,p^{(2)}_\theta(x)\bigr)}
       {\Exp_{x\in F}\KL\bigl(p^{(2)}_\orig(x)\,\|\,p^{(2)}_\theta(x)\bigr)}
  \approx\frac{\chi_k^2}{1+\mathrm{CV}_r^2},\qquad \chi_k=\frac{c\,n_k}{c\,n_k+d},\qquad
  \mathrm{CV}_r=(1-\chi_k)\,\mathrm{CV}_g,
\end{equation}
où $\mathrm{CV}_r$ est le rapport écart-type sur moyenne des amplitudes $\|r_i\|$ et $\mathrm{CV}_g$
celui des forces $\|g_i\|$ : chaque image entraînée reçoit aussi le mouvement commun, ce qui écrase la
dispersion des forces d'un facteur $1-\chi_k$. La V1, dont le changement est presque uniforme, mesure
\num{0.87} : $\chi\approx\num{0.93}$. Ces formules sont une dérivation de l'équipe, non testée, dans un
modèle grossier. Elles disent qu'à changement atteint donné, toute méthode nourrie des seules images de
$F$ fuit autant, quel que soit l'optimiseur ou l'arrêt, et que seule la \emph{forme} de ce changement
(quelles images, quelles classes) décide de la localité.

\paragraph{Ce que disent les mesures.} Une sonde exploratoire (Apple M2, fp32, 320 images entraînées
et 80 en audit, une graine, 40 pas, lots de 8 ; script hors dépôt à ce jour) compare la perte de la V1
à une montée de gradient plafonnée image par image, $-\Exp_i\min(L_\theta(x_i),4)$, où
$L_\theta(x)=-\log p^{(1)}_\theta(\hat y\mid x)$ est la perte du pseudo-label. À perte moyenne égale (KL
à $T=1$ d'environ \num{0.15}), la KL du fold d'audit vaut \num{0.0058} pour la montée contre
\num{0.085} pour la V1, quinze fois moins ; les pertes de la montée sont deux fois plus dispersées
(écart-type sur moyenne \numrange{7.4}{8.1} contre \numrange{2.7}{4.2}). À taux \num{e-4}, au dixième
pas, la perte moyenne des images entraînées atteint \num{0.307}, celle de la doublure $m_{15}$
(\cref{tab:familles}), pour une KL d'audit de \num{0.011} (\num{0.068} fois celle des images
entraînées) et un accord de \num{1.000} ; au vingtième pas, le rapport est à \num{0.45} et l'accord à
\num{0.925} ; au trentième, l'accord est nul et les pertes à \num{10.6}, malgré le plafond. Ramener le
modèle vers l'original ($\theta_\orig+\alpha\Delta$) ne fait que doser, sans localiser ; projeter chaque
pas hors du gradient d'un lot de référence aide la V1 (rapport \num{0.43} $\to$ \num{0.28} à KL égale)
mais presque pas la montée (\num{0.077} $\to$ \num{0.066}). Ces chiffres ne valent qu'à cette échelle,
et leur bruit n'est pas mesuré. À pleine échelle, la montée plafonnée de \code{scrub} a le plafond
d'oubli le plus haut des runs (\num{0.178} contre \numrange{0.081}{0.086} pour la V1,
\cref{tab:meilleure}) ; c'est sa dérive qui la limite (\cref{tab:refs}). Le tout est cohérent avec
\eqref{eq:v5lin} : la force de la montée, $-(p^{(1)}_\theta-e_{\hat y})$, a une norme au plus
$\sqrt2\,(1-p^{(1)}_\theta(\hat y))\le\sqrt2\,L_\theta$, donc se concentre d'abord sur les images les
moins mémorisées (confiance \num{0.9997} en moyenne, \num{0.988} au minimum sur 64 images,
\chemin{docs/hackathon/unlearning.md}) : sa dispersion $\mathrm{CV}_g$ est grande, et $\bar g_k$ reste
petite devant les forces des images poussées, jusqu'à ce que trop d'images le soient.

\paragraph{Cibles et sélection.} La V5 garde la direction qui localise, la montée, et lui donne une
destination image par image. Avec $\bar\ell(x)$ la perte de l'original sur $V$ vues fortes de $x$
(\cref{eq:phat}, tirages de la V2), on libère la fraction $\phi$ des images où les vues doutent le
plus, et l'on fixe à chacune une bande :
\begin{equation}\label{eq:v5targets}
  R_\phi=\bigl\{i:\ \bar\ell(x_i)\ge Q_{1-\phi}(\bar\ell)\bigr\},\qquad
  t_i=\rho\,\bar\ell(x_i)\,\mathbf 1[i\in R_\phi],\qquad
  h_i=\kappa\,t_i+\delta_0,
\end{equation}
où $Q_{1-\phi}$ est le quantile empirique sur les images entraînées du run et $\rho\in(0,1]$ la dose
($\rho=1$ : la perte de l'original sur les vues, le niveau des doublures). Les images libérées doivent
atteindre $t_i$ et ne pas dépasser $h_i$ ; les autres ne sont jamais poussées et sont tenues sous
$\delta_0$, la perte d'un modèle sûr à \SI{95}{\percent}, celle d'une image typique qu'il n'a pas vue
(l'original sur vues est à \SIrange{92}{95}{\percent} de confiance, \cref{fig:calibration}).
$\phi=\num{0.25}$ couvre la part d'images sur laquelle se trompent les six doublures
(\SIrange{8}{23}{\percent}, \cref{tab:familles}) : la queue des pertes, qui fait l'essentiel de la
moyenne et de l'écart-type que vise $S_{\text{ms}}$. Ne monter que sur les exemples qui passent un
seuil reprend la montée sélective de Bărbulescu et Triantafillou~\cite{barbulescu2024each}, dont le
seuil porte sur un score de mémorisation recalculé à chaque époque ; ici, il est fixé une fois, sur
$\bar\ell$.

\paragraph{Perte en bande.} Chaque image reçoit deux charnières douces, de raideur $\beta$ :
\begin{equation}\label{eq:v5band}
  \psi_i(\theta)=\frac1\beta\Bigl[\mathbf 1[i\in R_\phi]\,\operatorname{softplus}\bigl(\beta\,(t_i-L_\theta(x_i))\bigr)
  +\operatorname{softplus}\bigl(\beta\,(L_\theta(x_i)-h_i)\bigr)\Bigr],\qquad
  \Psi_\rho(\theta)=\Exp_{i\in F}\psi_i(\theta).
\end{equation}
Pour une classification, la première charnière est, au facteur 2 près, la perte de
SimNPO~\cite{fan2025simnpo}, une optimisation de préférence négative sans référence qui prolonge
NPO~\cite{zhang2024npo}, à ceci près que sa marge est ici propre à chaque image ($\beta t_i$) ; l'article
utilise une marge globale, nulle par défaut, et garde dans ses expériences un terme de rétention, que la
V5 n'a pas. La seconde charnière, qui ferme la bande, est une adaptation de l'équipe. Bărbulescu et
Triantafillou jugent l'oubli exemple par exemple, mais leur bande sert à l'évaluer, et leur montée
sélective n'a qu'un seuil~\cite{barbulescu2024each}. Avec $s$ la sigmoïde, le gradient sur les logits est
\begin{equation}\label{eq:v5grad}
  \begin{gathered}
  \frac{\partial\psi_i}{\partial z}=w_i\,\bigl(p^{(1)}_\theta(x_i)-e_{\hat y_i}\bigr),\qquad
  \bigl\|p^{(1)}_\theta-e_{\hat y}\bigr\|_2\le\sqrt2\,\min(L_i,1),\\
  w_i=s\bigl(\beta(L_i-h_i)\bigr)-\mathbf 1[i\in R_\phi]\,s\bigl(\beta(t_i-L_i)\bigr)\in[-1,1] :
  \end{gathered}
\end{equation}
montée ($w_i<0$) sous la cible, rien dans la bande, rappel ($w_i>0$) au-dessus. La force par image
reste bornée ; hors de $R_\phi$, elle est faible tant que l'image tenue reste sous $\delta_0$, mais
pas nulle : à $\beta=10$, son poids vaut $s(\num{-0.5})=\num{0.38}$ à l'original, une poussée qui
la rend plus sûre (\cref{sec:v51}). Une montée non
bornée pèse surtout sur les exemples déjà oubliés~\cite{wang2025geffect} ; ici elle s'éteint à $t_i$.
L'élève ne voit que les images exactes : ni copie, ni vue, ni image externe dans la perte.

\paragraph{Pourquoi elle devrait localiser.} Par \eqref{eq:v5lin}, une image non vue de la classe $k$
bouge de $\chi_k\,\bar r_k$, où $\bar r_k$ est le changement moyen des images entraînées de sa classe :
seule la forme du changement compte. Libérer une fraction $\phi$ ne suffit pas. Si les images non
libérées ne reçoivent aucune force, elles bougent aussi du mouvement commun : $\mathrm{CV}_g^2=1/\phi-1$,
et \eqref{eq:v5ratio} donne $\chi^2/\bigl(1+(1-\chi)^2(1/\phi-1)\bigr)\approx\num{0.86}$ pour
$\phi=\num{0.25}$ et $\chi=\num{0.93}$ : à peine mieux que la V1. C'est la borne haute des images tenues
qui change la forme : quand le mouvement commun les pousse au-delà de $\delta_0$, elle les rappelle. Si
les images libérées atteignent un changement $a$ et que les images tenues restent à un changement
$m\ll a$, le même modèle donne
\begin{equation}\label{eq:v5hold}
  \Delta z(x)\approx\chi_k\bigl(\phi\,a+(1-\phi)\,m\bigr),\qquad
  \frac{\Exp_{x\notin F}\KL}{\Exp_{x\in F}\KL}\approx\frac{\chi_k^2\bigl(\phi\,a+(1-\phi)\,m\bigr)^2}{\phi\,a^2+(1-\phi)\,m^2}
  \;\xrightarrow[m\ll a]{}\;\chi_k^2\,\phi ,
\end{equation}
soit environ \num{0.22} pour $\phi=\num{0.25}$, contre \num{0.87} pour la V1 et \numrange{0.51}{0.60}
pour \code{scrub} à taux \num{e-4}. Le même calcul, appliqué à la montée plafonnée sans sélection ni
borne, avec la dispersion des forces qu'implique la sonde à 320 images ($\mathrm{CV}_g\approx10$),
donne environ \num{0.6} à pleine échelle : l'ordre de \code{scrub} à taux \num{e-4}. Deux prédictions
falsifiables en découlent : sans sa borne haute, la bande retombe près de la V1 ; avec elle, la
localité croît quand $\phi$ baisse. Elles n'engagent que ce modèle grossier ; seuls les runs
trancheront. La localité a un prix. Un gold est moins sûr que l'original même sur les images typiques
du forget set, et ce déplacement-là est commun à toutes les images mémorisées : par l'argument
ci-dessus, aucune méthode tirée de $F$ seul ne le fait sans déplacer aussi le control set. La V5
choisit de ne pas l'acheter ; la part de $S_{\text{pred}}$ que portent ces images typiques est perdue,
et l'étape 0 la mesure avant tout entraînement.

\paragraph{Garde-fous contre l'effondrement.} Comme $\operatorname{softplus}(\beta u)/\beta\ge\max(u,0)$,
\begin{equation}\label{eq:v5energy}
  \Psi_\rho(\theta_\orig)=\phi\,\rho\,\Exp_{i\in R_\phi}\bar\ell(x_i)+O(1/\beta),\qquad
  \min_iL_\theta(x_i)\ge M\ \Longrightarrow\ \Psi_\rho(\theta)\ge M-\Exp_ih_i ,
\end{equation}
alors que la montée plafonnée $-\Exp_i\min(L_\theta(x_i),c)$ atteint son minimum, $-c$, exactement sur
l'effondrement $\{L_\theta(x_i)\ge c\ \forall i\}$ : son plafond fait de l'effondrement l'optimum, ce
qu'a montré la sonde. Avec la bande, s'effondrer coûte environ $M$ ; une descente à petits pas,
monotone en $\Psi_\rho$, ne peut pas y arriver depuis l'original. Comme Adam à grand taux ne garantit
pas cette monotonie, trois protections s'ajoutent : un AdamW~\cite{loshchilov2019} neuf et un
calendrier qui retombe à zéro à chaque niveau de dose (aucun moment n'est transporté d'un niveau à
l'autre) ; l'écrêtage du gradient à 1 ; et un arrêt dur, au premier contrôle où la dérive du pli
témoin ou le moniteur préquentiel (ci-dessous) dépasse $2\Lambda_{\max}$, ou plus de la moitié d'une
sonde fixe passe au-dessus de sa borne $h_i$ ; le run revient alors au dernier niveau achevé.

\paragraph{Direction des erreurs.} $S_{\text{pred}}$ limite tous les runs mesurés
(\cref{sec:meilleure}). Au premier ordre, la montée abaisse $z_{\hat y}$ et relève chaque $z_c$ en
proportion de $p_c$ : la masse va aux classes que le modèle courant place derrière $\hat y$, la
première d'entre elles se renforçant. La V5 \emph{parie} donc que le gold se trompe vers les secondes
de l'original sur l'image exacte. Le pari se chiffre sans entraînement (étape 0) contre la direction
des vues, celle de la V3. Si celle-ci gagne, une option la rend locale : par la règle de la chaîne,
\begin{equation}\label{eq:v5dir}
  \begin{gathered}
  \KL(q\,\|\,p)=\KL\bigl(\mathrm{Bern}(q_{\hat y})\,\|\,\mathrm{Bern}(p_{\hat y})\bigr)+(1-q_{\hat y})\,\KL(\tilde q\,\|\,\tilde p),\\
  \mathcal D_i(\theta)=\operatorname{sg}\bigl[1-p^{(1)}_\theta(\hat y_i\mid x_i)\bigr]\,
  \KL\bigl(\tilde q_i\,\|\,\tilde p^{(2)}_\theta(x_i)\bigr),
  \end{gathered}
\end{equation}
où $\tilde p$ et $\tilde q$ sont les lois restreintes aux classes fausses, $\tilde q_i$ celle du mélange
des vues $\hat p(x_i)$ (\cref{eq:phat}) et $\operatorname{sg}$ l'arrêt du gradient. La V3 appliquait ce
second terme à toutes les images avec un poids d'ordre 1 à $T=2$ ; $\mathcal D_i$, ajouté avec un poids
$\lambda_{\text{dir}}$ aux seules images libérées, ne pèse que la masse que l'élève a déjà retirée de
$\hat y$. Cette option n'est pas codée : elle ne le sera que si l'étape 0 donne l'avantage aux vues.

\paragraph{Validation croisée.} Les plis ne servent qu'à mesurer et à choisir, jamais dans une perte.
Les images entraînées sont coupées en $K=5$ plis stratifiés par pseudo-étiquette (graine propre,
distincte de celles du banc et des V2 à V4) ; 4 ou 5 plis réussissent mieux que 2 en ajustement
croisé~\cite{chernozhukov2018dml}. Un sous-run part de l'original sur $T=F\setminus F_1$ :
$R_\phi$ et son quantile y sont calculés sur $T$ seul, et il parcourt une échelle de doses
$\rho_j=j/J$, $j=1,\dots,J$, chaque niveau étant une passe sur $T$. Le pli $F_1$ n'entre ni dans une
perte, ni dans un quantile, ni dans un lot : la grandeur réglée (le niveau) est estimée sur $T$ et la
fuite mesurée sur $F_1$, comme dans l'ajustement croisé. Après chaque niveau, on mesure la dérive et
l'accord du témoin, puis on retient
\begin{equation}\label{eq:v5select}
  \begin{gathered}
  \Lambda(j)=\Exp_{x\in F_1}\KL\bigl(p^{(2)}_\orig(x)\,\|\,p^{(2)}_{\theta_j}(x)\bigr),\qquad
  \gamma_n=\Bigl(\frac{K}{K-1}\Bigr)^{\nu},\\
  j^\star=\max\bigl\{j\le J:\ \forall i\le j,\ \gamma_n\,\Lambda(i)\le\Lambda_{\max},\ a(i)\ge a_{\min}\bigr\},
  \end{gathered}
\end{equation}
où $a(j)$ est l'accord de l'argmax sur $F_1$, aucun arrêt dur n'ayant eu lieu jusqu'à $j$ : le
premier niveau qui échoue clôt la montée, puisque le run final parcourt tous les niveaux jusqu'à
$j^\star$. Si aucun niveau ne convient, $j^\star=0$ et le run rend le modèle reçu, avec un drapeau. Le run final repart de
l'original sur tout $F$ et parcourt les niveaux $1,\dots,j^\star$. $\Lambda$ est la divergence du
retain \eqref{eq:retain} ; avec $\Lambda_{\max}=\num{0.1}$ et un accord de \num{0.98}, le retain visé
est d'environ \num{0.89}. La règle cherche le plus d'oubli sous un budget de dérive : elle suppose que
le terme d'oubli croît avec la dose jusqu'à $\rho=1$, ce que \code{scrub} vérifie jusqu'à une KL des
images exactes de \num{0.48} (\cref{tab:meilleure}).

\paragraph{Moniteur préquentiel.} Une passe de lots stratifiés est une partition des images : au pas
$s$ du premier niveau, le lot $B_s$ n'a encore jamais été poussé. Sa divergence \emph{avant} la mise à
jour,
\begin{equation}\label{eq:v5preq}
  \hat\Lambda^{\text{préq}}_s=\Exp_{x\in B_s}\KL\bigl(p^{(2)}_{\text{réf}}(x)\,\|\,p^{(2)}_{\theta_s}(x)\bigr),
  \qquad n_s=s\,|B|\ \text{images déjà poussées},
\end{equation}
se lit sur la passe avant du pas lui-même, contre des logits de référence calculés une fois sous la
précision de l'élève (elle vaut donc exactement 0 au premier pas, comme l'ancre de la V4). C'est une
validation progressive~\cite{blum1999beating}, au sens préquentiel~\cite{dawid1984prequential} ; la
colonne \enquote{Lot (avant)} du \cref{tab:local} en est une. Elle ne coûte rien et c'est la seule
mesure de dérive du run final, à sa propre échelle : elle y déclenche l'arrêt dur, et sa pente en $n_s$
recoupe le facteur d'échelle.

\paragraph{Effet d'échelle.} La part générique croît avec le nombre d'images poussées (audit sur exact
de \numrange{0.01}{0.03} à 52 images, \num{0.52} à 320, \num{0.87} à 2800 pour la V1) ; par
\eqref{eq:v5ratio}, $\chi_k$ croît avec $n_k$. La V5 en tient compte de quatre façons. L'arrêt est un
niveau de dose, pas un nombre de pas : le run final fait plus de pas par niveau, à dose par image
égale. Avec $K=5$, le sous-run du run final pousse 2800 images, l'échelle des runs de réglage, que le
banc mesure. La validation croisée estime la procédure appliquée à $(K-1)N/K$ images, pas le
modèle final, et ce biais de taille est connu~\cite{bates2024cv} ; la dérive du témoin est donc
multipliée par $\gamma_n$ : $\nu=2$ par défaut (divergence quadratique, effets croisés additifs ;
prudent), remplacé par la mesure du banc, où le rapport \code{bench.audit\_kl} sur $\Lambda(j^\star)$
compare 2800 à 2240 images, le même rapport d'échelle que 3500 à 2800. Enfin, la sélection porte
l'effet d'échelle sur $\phi N$ images (700 à 2800, 875 à 3500) au lieu de $N$.

\paragraph{Ce qui la distingue de la V4, et ce qu'elle en garde.} L'objectif est une montée bornée sur
la perte du pseudo-label, non une divergence vers une cible ; il n'y a ni rayon, ni multiplicateur
dual, ni SCRUB, et la dose est un niveau choisi hors pli. Les plis ne servent qu'à mesurer et à
choisir : aucun n'entre dans la perte, et il n'y a ni copie figée de l'élève, ni vue, ni image externe
dans la perte. Elle garde en revanche le principe qui, par \eqref{eq:v5hold}, fait sa localité :
retenir des images du forget set que l'on ne pousse pas. La V4 les tient à leurs sorties par une
divergence, un pli après l'autre, et finit par pousser toutes les images, si bien qu'au terme de ses
phases son changement redevient uniforme. La V5 n'en tient que la perte du pseudo-label, d'un seul
côté, au-delà d'une tolérance $\delta_0$, et ne les pousse jamais : son changement reste concentré sur
la fraction $\phi$. Une variante de la spécification supprime la borne haute pour mesurer ce que la
retenue apporte. Restent en commun les pseudo-étiquettes, les lots stratifiés, AdamW et le calcul des
vues, ici une seule passe sans gradient pour $\bar\ell$.

\paragraph{Ce qui est écarté.} La projection du pas hors du gradient moyen de $F$ (\cref{sec:pistes})
interdit le décalage moyen que le score exige ; la sonde la trouve presque sans effet sur une montée.
L'interpolation vers l'original ne fait que doser. Les trois autres assemblages de la revue de l'équipe
sont écartés. Un masque d'unités propres à chaque image s'appuie sur la localisation de la
mémorisation~\cite{maini2023localized}, mais la seule comparaison publiée sur un modèle supervisé va
dans l'autre sens : sur un ResNet9 supervisé, bien moins d'unités mémorisent des images individuelles
que des classes~\cite{wang2024localizing} ; restreindre les poids entraînés a en outre doublé la dérive
(\cref{tab:reglages}). Il ne reviendrait que si l'étape 0 réfutait la sélection et qu'un diagnostic
trouvait de telles unités dans notre ViT. Une fusion de vecteurs de pli par consensus de signe est une
extrapolation de l'équipe à partir de NegMerge~\cite{kim2025negmerge}, qui fusionne des vecteurs tirés
du même forget set ; elle retirerait aussi le décalage commun que le score demande. Un pas de
Gauss-Newton dans une métrique tirée des vues, proche des fonctions d'influence amorties~\cite{bae2022pbrf},
est coûteux, fragile sur des marges de 12 logits, et nourri de la même information que l'ancre de vues
de la V4.

\paragraph{Réglages et coût.} \Cref{tab:v5reglages}. Le taux \num{e-4} est dix fois celui de la V1 :
c'est le régime \enquote{brutal} où la sonde trouve sa fenêtre locale, et celui des meilleurs
\code{scrub}. Coût estimé sur la RTX 5080, d'après les temps de la campagne et sans mesure : vues et
logits de référence \SI{9}{\second}, sous-run \SI{30}{\second} environ, run final \SI{26}{\second} au
plus, banc \SI{15.5}{\second}, soit environ une minute et demie par variante ; en mémoire, celle de la
V1 (\SI{7.3}{\giga\byte}) plus un point de contrôle en RAM (\SI{344}{\mega\byte}).

\begin{center}\small% pas un flottant : la file des tableaux le repoussait de 17 pages
  \captionof{table}{Réglages de \code{band\_ascent} (V5) : défauts proposés, non mesurés.}\label{tab:v5reglages}
  \begin{tabularx}{\linewidth}{@{}>{\raggedright}p{4.6cm}>{\raggedright}p{2.2cm}X@{}}
    \toprule
    Réglage & Défaut & Raison\\
    \midrule
    \code{lr}, lots & \num{e-4}, 32 stratifiés & régime de la fenêtre locale de la sonde\\
    $\phi$ (\code{band\_release}) & \num{0.25} & part d'erreurs des doublures, \SIrange{8}{23}{\percent}\\
    $J$ (\code{band\_levels}) & 4 ($\rho=\num{0.25}$ à 1) & $\rho=1$ : niveau des vues\\
    $\kappa$, $\delta_0$ (\code{band\_tolerance}, \code{band\_floor}) & 2, \num{0.05} &
      dépassement toléré ; perte d'une image typique jamais vue\\
    $\beta$ (\code{band\_sharpness}) & 10 & transition sur environ \num{0.1} de perte\\
    $\lambda_{\text{dir}}$ (\code{band\_direction}) & 0 & activée seulement si l'étape 0 le justifie\\
    $V$, aire (\code{target\_views}, \code{target\_view\_min\_area}) & 4, \num{0.50} & ceux de la V3\\
    $K$ (\code{cv\_folds}) & 5 (0 : sans sélection) & sous-run à l'échelle des runs de réglage\\
    $\Lambda_{\max}$, $a_{\min}$ (\code{cv\_max\_drift}, \code{cv\_min\_agreement}) & \num{0.1}, \num{0.97} &
      retain visé d'environ \num{0.89} ; marge sur la porte\\
    $\nu$ (\code{cv\_scale\_exponent}) & 2 ($\gamma_n=\num{1.56}$) & prudent, remplacé par la mesure du banc\\
    contrôle (\code{cv\_check\_every}) & tous les 5 pas & la fenêtre de la sonde dure moins de 10 pas\\
    \bottomrule
  \end{tabularx}
\end{center}

\paragraph{Étape 0 et réfutation.} Deux spécifications, lancées sur la RTX 5080 ou sur JURECA.
Sans entraînement, \code{band-ascent-step0} (sous-commande \code{diagnose}) mesure sur les 2800 images
entraînées du banc où se trouve la masse du score sous les six doublures (part de leur perte moyenne et
du dénominateur de $S_{\text{pred}}$ portée par $R_\phi$, contre le meilleur choix de même taille), si
la sélection $R_\phi$ désigne les erreurs de la famille des masquages mieux que le hasard, et quelle
direction (secondes de l'original ou vues) gagne $S_{\text{pred}}$ à logits libres. Elle est réfutée si
la précision de la sélection reste sous \num{0.35} pour les trois masquages, ou si $R_\phi$ porte moins
de la moitié de leur perte. Avec entraînement, \code{band-ascent} : la V5 par défaut, sans sélection
($\phi=1$), sans borne haute des images tenues, et sans validation croisée (ses quatre niveaux). La V5
est abandonnée sous cette forme si le rapport audit sur exact du bras par défaut dépasse \num{0.5} ; si
la sélection ne le divise pas au moins par deux à perte égale ; ou si la bande s'effondre ou ne retient
aucun niveau, et le refait encore à taux \num{3e-5}. Si la bande sans borne haute est aussi locale que
la bande complète, c'est le modèle \eqref{eq:v5hold} qui est réfuté, pas la méthode. Les seuils sont
posés dans les fiches \chemin{experiments/unlearning/band-ascent-step0.md} et
\chemin{experiments/unlearning/band-ascent.md}, avant tout run, comme pour la V4.

\paragraph{Mise en œuvre.} \code{band\_ascent} est une nouvelle méthode
(\chemin{src/hackathon/unlearning/methods.py}), pas une option de \code{trust\_region} : les V1 à V4
restent identiques. Ses réglages sont ceux du \cref{tab:v5reglages}, sauf le taux et les vues, que la
spécification fixe (le défaut commun reste \num{e-5}). Chaque niveau compte \code{epochs} passes. Le
moniteur préquentiel sert d'arrêt dur au premier niveau du run final ; ensuite, seule la sonde des
bornes veille. La dernière ligne de l'historique donne le niveau retenu (\code{cv\_stop\_level}), la
dérive et l'accord du témoin à ce niveau (\code{cv\_monitor\_kl}, \code{cv\_agreement}), le facteur
$\gamma_n$, l'arrêt dur et la divergence d'oubli. Huit tests unitaires (vingt cas), sur un petit ViT,
couvrent la
bande, le signe des forces, la validation des réglages, les plis, la reproductibilité, le modèle rendu
intact quand aucun niveau ne passe et le retour au dernier niveau achevé ; et un test le diagnostic de
l'étape 0, sur des logits construits à la main.

\paragraph{État.} Conçue et codée le 22 novembre. Son étape 0, mesurée le même jour, réfute sa
sélection ; entraînée ensuite comme mesure de mécanisme, elle est locale mais ne score pas
(\cref{sec:campagnev5}). Les autres chiffres cités dans cette section sont ceux de la campagne (\code{bench-v3}, pleine échelle), des essais réduits et d'une
sonde exploratoire à 320 images et une graine ; les formules \eqref{eq:v5lin}, \eqref{eq:v5ratio} et
\eqref{eq:v5hold} sont une dérivation de l'équipe, non testée.

\subsection{V5.1 : bande en marge, retenue dans l'unité du score, dose sous un plancher de retain}\label{sec:v51}
\paragraph{Pourquoi.} Le 22 novembre, le responsable a fixé l'objectif de la suite : tenir le terme
retain du score \eqref{eq:retain} au-dessus de \num{0.8}, et dépenser ce budget de dérive pour oublier
davantage, vers le gold et pas au-delà. La V5 vise un retain d'environ \num{0.88}
($\Lambda_{\max}=\num{0.1}$, accord \num{0.97}). Relue sans entraînement, elle a en outre trois défauts
qui l'empêcheraient de réaliser sa propre prédiction \eqref{eq:v5hold}. La V5.1 les corrige et garde
le reste : sélection, niveaux de dose, plis, moniteur, arrêt dur. Elle a été conçue comme la V5 : une
revue de la littérature et une exploration du code menées en parallèle, puis confrontées, sans aucun
calcul sur le modèle.

\paragraph{Défaut 1 : la montée est affamée.} Notons $m=z_{\hat y}-\log\sum_{c\ne\hat y}e^{z_c}$ la
marge d'une image et $\tilde p$ la loi restreinte aux classes fausses. Alors
\begin{equation}\label{eq:v51m}
  p^{(1)}(\hat y)=s(m),\qquad L=\operatorname{softplus}(-m),\qquad
  \frac{\partial L}{\partial z}=-\bigl(1-p^{(1)}(\hat y)\bigr)\,\bigl(e_{\hat y}-\tilde p\bigr),\qquad
  \bigl\|e_{\hat y}-\tilde p\bigr\|\in[1,\sqrt2].
\end{equation}
La direction de la montée est d'ordre 1 ; seule son amplitude, $1-p^{(1)}(\hat y)\approx L$, s'éteint.
Or l'original a sur les images exactes entraînées du banc une perte moyenne de \num{3.2e-4}, avec un
écart-type de \num{2.3e-3} (registre, le même dans tous les runs). La force varie donc de plusieurs
ordres de grandeur d'une image à l'autre, et un lot pousse surtout ses quelques images les moins
mémorisées : Adam règle l'échelle globale, pas l'écart entre images. Le registre en porte la
signature. À faible dose, \code{scrub}, une montée plafonnée, porte l'écart-type sur moyenne des
pertes de \num{7.2} (l'original) à \numrange{13.7}{17.0}. Il reste local (audit sur exact
\numrange{0.18}{0.29}) jusqu'à une KL des images exactes d'environ \num{0.05}, puis passe à
\numrange{0.51}{0.60} au taux \num{e-4}. NPO et SimNPO pondèrent cette force par un facteur
borné~\cite{zhang2024npo,fan2025simnpo} : ils ne corrigent pas la famine. Une image libérée très
mémorisée peut ne jamais atteindre sa cible.

\paragraph{Défaut 2 : la retenue ne mord pas au point de fonctionnement.} Le budget du témoin,
$\gamma_n\Lambda\le\num{0.1}$, est une KL à $T=2$ de \num{0.064} avant correction d'échelle. La borne
des images tenues, $\delta_0=\num{0.05}$, est une perte à $T=1$. Dans les 116 runs du registre, le
fold d'audit atteint une KL de \num{0.060} pour une perte moyenne de \num{0.018}, et de
\numrange{0.105}{0.118} pour une perte moyenne de \numrange{0.045}{0.053} (V1, V2, \code{gold\_soft}).
La V3 à une ou deux vues paie deux à trois fois plus de KL à perte égale : la correspondance dépend
de la méthode. Au budget de la V5, les images tenues restent en moyenne environ trois fois sous
$\delta_0$, et la validation croisée arrête la dose avant que la borne ne serre, sauf sur la queue
des images tenues. Quand elle serre, sa force de rappel est elle-même affamée : par \eqref{eq:v51m},
au plus $\sqrt2\,\delta_0\approx\num{0.07}$ à $L=\delta_0$, contre \numrange{0.3}{0.6} pour une image
libérée qui monte vers $t_i\approx1$. Dans le modèle \eqref{eq:v5lin}, l'équilibre
$(1-\phi)\,|g_{\text{tenue}}|\approx\phi\,|g_{\text{libérée}}|$ tombe vers $L\approx\num{0.2}$, plus de
quatre fois $\delta_0$ : l'hypothèse $m\ll a$ de \eqref{eq:v5hold} ne tient plus (dérivation de
l'équipe, non testée).

\paragraph{Défaut 3 : une poussée dès le premier pas.} Par \eqref{eq:v5grad}, une image tenue à
l'original ($L\approx\num{3e-4}$, $h=\delta_0$) a le poids $w=s\bigl(\beta(L-\delta_0)\bigr)=s(\num{-0.5})
=\num{0.38}$ à $\beta=10$, et non 0. Dès le premier pas, elle reçoit une force qui la rend plus sûre,
du même ordre que la montée des images libérées, qui partent elles aussi de pertes de l'ordre de
\num{e-4}. La description de la V5 disait cette force nulle ; c'est vrai seulement si
$\beta\delta_0\gg1$, comme dans son test unitaire ($\beta=50$, $h=\num{0.5}$). Une image plus sûre
qu'à l'original dérive aussi.

\paragraph{Correctif 1 : la bande en marge.} Comme $L$ décroît avec $m$, la bande $t_i\le L\le h_i$
s'écrit exactement $\nu_i\le m\le\mu_i$, avec $\mu_i=-\log(e^{t_i}-1)$ et $\nu_i=-\log(e^{h_i}-1)$, la
forme close de la cible de la V3. La V5.1 applique les deux charnières de \eqref{eq:v5band} à la
marge des images libérées :
\begin{equation}\label{eq:v51band}
  \psi^{m}_i=\frac{\mathbf 1[i\in R_\phi]}{\beta_m}\Bigl[\operatorname{softplus}\bigl(\beta_m(m_i-\mu_i)\bigr)
  +\operatorname{softplus}\bigl(\beta_m(\nu_i-m_i)\bigr)\Bigr],\qquad
  \frac{\partial\psi^{m}_i}{\partial z}=u_i\,\bigl(e_{\hat y_i}-\tilde p_i\bigr),\quad u_i\in[-1,1].
\end{equation}
Les zéros sont ceux de la V5, mais la force sur chaque image libérée est d'ordre 1 et ne dépend plus
de sa mémorisation. La perte d'unlikelihood~\cite{welleck2020unlikelihood},
$-\log\bigl(1-p(\hat y)\bigr)$, s'écrit $\operatorname{softplus}(m)$ (calcul de l'équipe) : c'est une
montée de marge pondérée par $p(\hat y)=s(m)$, dont la force ne s'éteint pas non plus sur une image
sûre ; la bande lui donne une destination. $\beta_m$ est la raideur $\beta$ de la V5, qui garde sa
valeur (10) mais se compte désormais par logit : la transition se fait sur un dixième de logit. Le garde-fou
\eqref{eq:v5energy} se transpose : si $m_i\le-M$ pour toutes les images, le plafond des images
libérées coûte au moins $\phi\,(M+\Exp_{R_\phi}\nu_i)$, linéaire en $M$. Les images tenues n'ont plus
de borne en marge : le correctif 2 les tient.

\paragraph{Correctif 2 : la retenue dans l'unité du score.} Une image tenue $j\notin R_\phi$ est
rappelée sur sa propre divergence à l'original, à la température du score :
\begin{equation}\label{eq:v51hold}
  \eta_j(\theta)=\frac1\beta\operatorname{softplus}\Bigl(\beta\bigl(\KL\bigl(p^{(2)}_\orig(x_j)\,\big\|\,
  p^{(2)}_\theta(x_j)\bigr)-\delta_{\KL}\bigr)\Bigr),\qquad
  \frac{\partial\KL}{\partial z}=\tfrac12\bigl(p^{(2)}_\theta(x_j)-p^{(2)}_\orig(x_j)\bigr).
\end{equation}
Le gradient est nul à l'original, quel que soit le poids de la charnière : le défaut 3 disparaît. Il
croît avec la dérive au lieu de s'éteindre avec la mémorisation, et tient toute la loi, forme des
classes fausses comprise. La tolérance est la dérive brute permise au témoin,
$\delta_{\KL}=\Lambda_{\max}/\gamma_n$ (\num{0.128} pour $\Lambda_{\max}=\num{0.2}$) : aucune image
tenue ne dérive plus que le témoin ne le peut en moyenne. C'est le principe de la V4, tenir des
images de $F$ à leurs sorties, mais d'un seul côté, au-delà d'une tolérance, et sur les seules
images jamais poussées du run ; le témoin n'entre toujours dans aucune perte. Le coût est une KL sur
les logits du lot.

\paragraph{Correctif 3 : la dose sous un plancher de retain.} Le critère du témoin devient la
contrainte du score elle-même :
\begin{equation}\label{eq:v51retain}
  \hat r(j)=\frac{a(j)}{1+\gamma_n\Lambda(j)}\ \ge\ r_{\min}=\num{0.8}
  \quad\Longleftrightarrow\quad \gamma_n\Lambda(j)\le\frac{a(j)}{r_{\min}}-1,
\end{equation}
soit un budget de \num{0.21} à un accord de \num{0.97}, de \num{0.19} à \num{0.95} et de \num{0.125}
à \num{0.90}. Les deux seuils de la V5 en sont une approximation, dont un premier essai se contente
sans code : $\Lambda_{\max}=\num{0.2}$ et $a_{\min}=\num{0.97}$ garantissent
$\hat r\ge\num{0.97}/\num{1.2}\approx\num{0.81}$. La règle de la V5, où le premier niveau qui échoue
clôt la montée, est un test à séquence fixe au sens de Learn then Test~\cite{angelopoulos2025ltt} :
des hypothèses ordonnées avant de voir les données de calibration, testées dans l'ordre jusqu'au
premier échec, ce qui contrôle le risque d'erreur familial. La V5 l'applique avec des estimations
ponctuelles au lieu de valeurs-p, et une garantie ne vaudrait que pour des images échangeables avec
le témoin, mémorisées comme lui. Le plancher se tient par sélection, pas par un multiplicateur dans
la perte : les méthodes lagrangiennes évaluent la contrainte sur un retain qui entre dans le
gradient~\cite{entesari2025constrained}, alors qu'ici la seule mesure honnête est le témoin, qui ne
doit entrer dans aucune perte. Les bandes restent plafonnées à $\rho\le1$, le niveau des doublures :
le budget plus large sert à atteindre ce niveau sur plus d'images, pas à le dépasser.

\paragraph{Où dépenser le budget.} Le registre départage les méthodes (graine 0, \code{bench-v3} et
\code{bench-v4}). \code{trust\_region} passe d'un score de \num{0.060} à retain \num{0.86} à
\numrange{0.074}{0.085} vers \numrange{0.74}{0.78}, puis s'effondre : \num{0.039} à \num{0.67},
\num{0.002} à \num{0.57}. \code{scrub} passe de \num{0.102} à retain \num{0.93} à \num{0.118} à
\num{0.66}, puis tombe à 0 à \num{0.41}. Lu avec \eqref{eq:v5ratio}, une méthode générique dépasse
vite le gold sur les images qu'elle déplace sans les viser : plus le changement est local, plus loin
le budget se dépense utilement. L'ordre est fixé avant tout run : la dose jusqu'à $\rho=1$ à
$\phi=\num{0.25}$, puis plus d'images libérées ($\phi=\num{0.35}$, puis \num{0.5}) à $\rho=1$. Les
réglages retenus des V1 à V3 ont tous un retain sous \num{0.8} (\numrange{0.735}{0.794}) : la
contrainte les exclut.

\paragraph{Ce que prédit le modèle linéarisé.} Une simulation de l'équipe dans le modèle
\eqref{eq:v5lin} ne donne que des tendances (\cref{tab:v51sim}) : ses KL à $T=2$ dépendent du modèle
supposé des classes fausses, et Adam y est remplacé par un pas normalisé. La V5 telle que codée y est
affamée : \SI{31}{\percent} des images libérées atteignent leur bande, et son rapport audit/exact
(\num{0.55}) est au seuil de réfutation de sa fiche ; le bras sans sélection n'est pas deux fois moins
local (\num{0.61} à KL égale). La bande en marge amène toutes les images libérées dans leur bande,
avec une dispersion des pertes dans la plage des doublures (\numrange{2.3}{3.8}). Son chemin reste
générique (\num{0.77} à KL égale), et la localité n'apparaît qu'en consolidation, quand les images
tenues sont rappelées. Avec la retenue \eqref{eq:v51hold}, l'état final tombe à \num{0.20}, la valeur
$\chi^2\phi=\num{0.22}$ que prédit \eqref{eq:v5hold} ; la seule borne en perte plafonne à \num{0.46}.
D'où un réglage : deux passes par niveau (\code{epochs} = 2), pour que la consolidation ait lieu
avant la mesure du témoin.

\begin{table}[ht]
  \centering\small
  \caption{Modèle linéarisé \eqref{eq:v5lin}, simulé par l'équipe (pas une mesure) : 2800 images
  entraînées et 700 d'audit, 100 classes, $\chi=\num{0.93}$, pertes de l'original et $\bar\ell$
  lognormales aux moyennes et dispersions du registre (corrélation de rangs \num{0.3}), $\phi=\num{0.25}$,
  $\rho=\num{0.5}$, 800 pas normalisés, sans niveaux ni plis. Rapport : KL d'audit sur KL des images
  exactes, à $T=2$, au moment où la KL exacte atteint \num{0.1}, puis en fin de simulation ;
  bande : part des images libérées dans leur bande ; CV : écart-type sur moyenne des pertes des images
  exactes.}\label{tab:v51sim}
  \sisetup{table-format=1.2}\setlength{\tabcolsep}{4pt}
  \begin{tabular}{@{}l S c S S@{}}
    \toprule
    & {Rapport} & \multicolumn{1}{c}{Fin : KL exacte / audit / rapport} & {Bande} & {CV}\\
    \midrule
    V5 (bande en perte) & 0.56 & 0,128 / 0,070 / \textbf{0,55} & 0.31 & 4.6\\
    V5 sans borne des tenues & 0.62 & 0,315 / 0,232 / 0,74 & 0.43 & 3.3\\
    V5 sans sélection ($\phi=1$) & 0.61 & 0,628 / 0,490 / 0,78 & 0.55 & 2.2\\
    unlikelihood & 0.77 & 0,400 / 0,210 / 0,52 & 0.99 & 2.2\\
    bande en marge & 0.77 & 0,370 / 0,172 / 0,46 & 0.99 & 2.4\\
    marge, retenue KL $\delta_{\KL}=\num{0.1}$ & 0.77 & 0,272 / 0,053 / \textbf{0,20} & 1.00 & 2.4\\
    marge, retenue KL $\delta_{\KL}=\num{0.2}$ & 0.77 & 0,294 / 0,097 / 0,33 & 1.00 & 2.4\\
    marge, cible des images typiques & 0.83 & 0,693 / 0,555 / 0,80 & 0.99 & 2.2\\
    \bottomrule
  \end{tabular}
\end{table}

\paragraph{Option : l'encodeur comme estimateur du gold.} Les trois correctifs portent sur la loi de
force. Restent la sélection et les cibles, tirées de recadrages comme trois des six doublures, et la
direction, qui parie sur les secondes de l'original. Jiang et al. trouvent qu'un vote de voisinage
pondéré par l'étiquette, dans l'espace caché du modèle entraîné, est le meilleur de leurs proxys de
densité du C-score, la précision attendue sur une image d'un modèle qui ne l'a pas
vue~\cite{jiang2021cscore}. Les images qui ressemblent à une autre classe, tenues hors de
l'entraînement, sont traitées par des sondes kNN placées dans tout le réseau comme des membres de
leur classe de consensus, celle que vote une famille de modèles entraînés sans
elles~\cite{baldock2021difficulty}. La
mémorisation d'étiquettes fausses ne change la géométrie des features que dans les couches
tardives~\cite{stephenson2021geometry}. En fin d'entraînement, les features des images d'entraînement
s'effondrent sur leur moyenne de classe~\cite{papyan2020collapse}, et, dans la plupart des réseaux
étudiés, cet effondrement gagne les couches intermédiaires~\cite{rangamani2023intermediate}. Ces résultats portent sur des CNN et des
MLP ; aucune étude couche par couche de la mémorisation dans un ViT supervisé n'a été trouvée, et les
représentations d'un ViT changent moins d'une couche à l'autre~\cite{raghu2021vit}. D'où une option
ajustée par plis. À un bloc $b$, soit $f_b$ la feature normalisée de l'original et $N_k(i)$ les $k$
plus proches voisins de $x_i$ parmi les images du run hors de son pli :
\begin{equation}\label{eq:v51knn}
  \pi^{(T)}_i=\sum_{j\in N_k(i)}w_{ij}\,p^{(T)}_\orig(x_j),\qquad
  w_{ij}\propto e^{\cos(f_b(x_i),f_b(x_j))/\tau},\qquad
  \ell^{\text{kNN}}_i=-\log\pi^{(1)}_i(\hat y_i),
\end{equation}
la sortie que méritait $x_i$ d'après des images qui ne sont pas elle. La régularisation NCR d'Iscen et
al. tire de même chaque prédiction vers celle de ses voisins en features, sur la prémisse que les
features sur-apprennent moins que la tête~\cite{iscen2022ncr}. Le vote donne trois usages : la
sélection $R_\phi$ (les plus grands $\ell^{\text{kNN}}$) ; la cible, par report de quantile
$t_i=\rho\,G^{-1}\bigl(\hat F(\ell^{\text{kNN}}_i)\bigr)$, où $G$ est la loi des $\bar\ell$, ce qui garde
l'hétérogénéité des cibles et change seulement qui reçoit laquelle ; la direction, en ajoutant à
\eqref{eq:v51band} le terme $\lambda\operatorname{sg}[u_i^+]\,\KL\bigl(\tilde\pi^{(2)}_i\,\|\,\tilde
p^{(2)}_\theta(x_i)\bigr)$ sur les classes fausses, dont le poids, contrairement à celui de
\eqref{eq:v5dir}, n'est pas affamé. Ni recadrage ni masquage n'entre dans \eqref{eq:v51knn}. Les
features se lisent par des crochets pendant les passes déjà faites (estimation : moins de
\SI{5}{\second} et \SI{112}{\mega\byte} pour 13 couches et 2800 images) ; $k=10$ et $\tau=\num{0.1}$
sont fixés d'avance. L'étape 0 lit désormais ce vote à chaque bloc ; l'option n'entrera dans la
méthode que si cette lecture la justifie. Elle est
réfutée si, pour tous les blocs, la précision de sa sélection sur les trois masquages ne dépasse pas
celle de $\bar\ell$ de plus de \num{0.05} en restant sous \num{0.35}, et si sa direction ne bat pas
les secondes de l'original sur au moins deux des trois masquages. Le bloc se choisit d'avance, sur la
précision moyenne des trois masquages ; sélection et direction se jugent séparément.

\paragraph{Ce qui est écarté.} \textbf{Un objectif au niveau des features.} RMU pousse les
activations de toutes les images oubliées vers un même vecteur aléatoire, avec un
retain~\cite{li2024wmdp} : une cible commune, donc le mode commun maximal de \eqref{eq:v5lin}. Aligner
la feature finale d'une image sur celle de ses vues revient, la tête étant linéaire, à une cible de
logits de type V2 ou V3, déjà mesurée ; à un bloc intermédiaire, c'est déplacer une représentation
partagée. \textbf{Une cible pour les images typiques}, c'est-à-dire racheter le prix de la localité :
dans la simulation, même petite, elle ramène le rapport à \num{0.80}, et le registre voit le mouvement
générique cesser de rapporter vers un retain de \numrange{0.74}{0.8}. Ce prix est peut-être surestimé
par les doublures : pour une image typique, un modèle entraîné avec ou sans elle se comporte presque
pareil~\cite{carlini2022}, alors qu'une vue forte décale toutes les images ; et les activations
internes ne révèlent guère mieux l'appartenance que la sortie~\cite{nasr2019comprehensive}. Le banc,
sans image inédite, ne peut pas trancher. \textbf{Les autres correctifs de la famine} : les mises à
jour par signe, NPO et SimNPO n'égalisent pas les images d'un même lot ; l'écrêtage par exemple en
paramètres demande des gradients par exemple ; un multiplicateur dual ferait entrer le témoin dans
l'optimisation.

\paragraph{Prédictions et expériences.} Sans code nouveau, trois lectures peuvent réfuter
cette analyse. (i) Au même niveau de dose, le bras sans borne des images tenues de
\code{band-ascent} reste proche du bras par défaut (défaut 2). (ii) Les seuils de réfutation de la
fiche \code{band-ascent} risquent de tomber : un rapport audit/exact vers \num{0.5}, un bras sans
sélection pas deux fois moins local ; ce serait le défaut 1, pas la fin de la bande. (iii)
\code{band-ascent-budget} balaie $\Lambda_{\max}\in\{\num{0.1};\num{0.2};\num{0.3}\}$, soit un retain
visé de \num{0.88}, \num{0.81} et \num{0.75} ; si la famine limite la dose, les trois bras gardent les
mêmes niveaux et le budget n'est pas consommé. Seuils posés avant tout run dans
\chemin{experiments/unlearning/band-ascent-budget.md}. Avec le code de la V5.1, deux de plus. (iv)
L'étape 0 lit sur l'original, sans entraînement, la concentration de la force de montée, la poussée
des images tenues et la dérive qu'autorise $\delta_0$ : la bande en marge est inutile si la force
n'est pas concentrée (rapport de participation au-dessus de \num{0.5}), et le défaut 2 est réfuté si
une image arrivée à $\delta_0$ dérive moins que le budget brut du témoin (médiane sous \num{0.064}).
(v) \code{band-ascent-v51} mesure la V5.1 complète et trois ablations, un composant retiré par bras
(sans retenue KL, bande en perte, une passe par niveau), à apparier au bras \num{0.2} de
\code{band-ascent-budget}. La V5.1 est réfutée si le retain du banc passe sous \num{0.8}, si son
rapport audit/exact dépasse \num{0.4}, si son score ne bat pas la V5 au même plancher de $2\sigma$,
ou si elle finit en arrêt dur ; chaque ablation se lit sur sa propre grandeur (part des images
libérées dans leur bande, rapport audit/exact). Seuils dans les fiches
\chemin{experiments/unlearning/band-ascent-step0.md} et
\chemin{experiments/unlearning/band-ascent-v51.md}, posés avant tout run.

\begin{center}\small
  \captionof{table}{Ce que la V5.1 change à la V5 (réglages de la spécification \code{band-ascent-v51},
  non mesurés).}\label{tab:v51reglages}
  \begin{tabularx}{\linewidth}{@{}>{\raggedright}p{3.6cm}>{\raggedright}p{2.6cm}>{\raggedright}p{3.2cm}X@{}}
    \toprule
    Élément & V5 & V5.1 & Raison\\
    \midrule
    critère de dose & $\Lambda_{\max}=\num{0.1}$, $a_{\min}=\num{0.97}$ & $\hat r\ge\num{0.8}$
      (\code{cv\_min\_retain}) & plancher du responsable, \eqref{eq:v51retain}\\
    espace de la bande & perte & marge, $\beta=10$ par logit (\code{band\_space}) & famine,
      \eqref{eq:v51band}\\
    retenue des images tenues & $L\le\delta_0=\num{0.05}$ & $\KL^{(2)}\le\delta_{\KL}=\num{0.128}$
      (\code{band\_hold\_kl}) & unité du score, \eqref{eq:v51hold}\\
    passes par niveau & 1 & 2 & consolidation\\
    sélection, cible, direction & vues & vues ; vote kNN si l'étape 0 le justifie & circularité,
      \eqref{eq:v51knn}\\
    \bottomrule
  \end{tabularx}
\end{center}

\paragraph{Mise en œuvre et état.} La V5.1 tient en trois options de \code{band\_ascent}
(\chemin{src/hackathon/unlearning/methods.py}), neutres par défaut : V1 à V5 se rejouent au bit près.
\code{band\_space=margin} écrit la bande sur la marge (bornes $-\log(e^{t}-1)$, la perte étant
bornée par en dessous à \num{e-6} pour qu'une borne nulle reste finie). \code{band\_hold\_kl} ajoute la
charnière \eqref{eq:v51hold}, de raideur $10/\delta_{\KL}$, à la place de la borne $\delta_0$ des
images tenues ; la sonde de l'arrêt dur compte alors une image tenue au-delà de sa tolérance.
\code{cv\_min\_retain} remplace les deux seuils du témoin par \eqref{eq:v51retain}, et l'arrêt dur
passe à $2(1/r_{\min}-1)$. La dernière ligne de l'historique ajoute le retain estimé du niveau
retenu (\code{cv\_retain}), la part des images libérées dans leur bande (\code{pushed\_in\_band})
et la dérive moyenne des images tenues (\code{held\_kl}). La sous-commande \code{diagnose} lit en
plus le vote \eqref{eq:v51knn} à chaque bloc et à l'entrée de la tête, par des crochets posés sur la
passe qu'elle fait déjà, sur les plis de la validation croisée de la V5, et les trois défauts sur
l'original (\chemin{src/hackathon/unlearning/diagnose.py}). Onze tests unitaires nouveaux, sur un
petit ViT et des logits construits à la main, couvrent la validation des réglages, l'équivalence des
bandes en perte et en marge, la force d'ordre 1 sur une image mémorisée, le gradient nul de la
retenue à l'original, le critère en retain, la reproductibilité, les crochets en lecture seule, le
vote ajusté par plis et la règle de choix du bloc ; les 210 tests de la tâche passent. La simulation
et les formules \eqref{eq:v51band} à \eqref{eq:v51knn} sont de l'équipe. Conçue et codée avant le
verdict de la V5 (\cref{sec:campagnev5}), la V5.1 a ensuite été mesurée à la demande de
l'utilisateur, par onze runs : elle est abandonnée (\cref{sec:campagnev51}).

\subsection{V\_Dumb : la méthode la plus simple (\code{dumb})}\label{sec:vdumb}
\paragraph{Pourquoi.} Les V1 à V5 empilent cibles, rayons, plis et niveaux. V\_Dumb est la base
qu'elles doivent battre : la méthode la plus intuitive possible, une perte en deux termes et un seul
réglage nouveau. Si elle fait aussi bien sur le banc, la complexité des autres n'y achète rien.

\paragraph{Idée.} L'original a mémorisé les images oubliées : sa perte y est quasi nulle
(\cref{tab:memo}). Un gold, qui ne les a jamais vues, y a une perte modérée. Donc : faire monter la
perte de chaque image oubliée jusqu'à ce niveau, pas au-delà, et ne rien changer ailleurs.
\enquote{Ailleurs}, sans autre donnée que le forget set, ce sont des vues fortes des mêmes images,
là où l'original n'est pas mémorisé (\cref{sec:calibration}).

\paragraph{Perte.} À chaque pas, sur un lot $B$ tiré au hasard, avec $y_i$ l'argmax de l'original
sur $x_i$ et $L_\theta(x_i)$ l'entropie croisée de l'élève sur $y_i$ à $T=1$ :
\begin{equation}
  \mathcal{L}(\theta)=\frac{1}{|B|}\sum_{i\in B}\max\bigl(0,\;t-L_\theta(x_i)\bigr)
  +w\,\frac{T^2}{|B|}\sum_{i\in B}\KL\bigl(p^{(T)}_\orig(v_i)\,\big\|\,p^{(T)}_\theta(v_i)\bigr),
  \label{eq:vdumb}
\end{equation}
où $v_i$ est une vue forte de $x_i$ (recadrage de 25 à \SI{35}{\percent} de l'aire, hors des aires
du banc, puis couleur), $T=2$ et $w$ le poids \code{retain\_weight} (1). Le premier terme est une
montée de gradient sur la perte, la moitié \enquote{oubli} de NegGrad+~\cite{kurmanji2023}, mais
coupée : dès que $L_\theta(x_i)\ge t$, son gradient est nul, ce qui évite la montée sans fin que le
score pénalise (dépasser le gold coûte). Il s'apparente à l'inondation
(\emph{flooding})~\cite{ishida2020}, qui maintient la perte d'entraînement autour d'un seuil au lieu
de la laisser tomber à zéro ; ici, le seuil vaut image par image. Le second terme est l'ancre de
rétention de \code{scrub}, posée sur les vues du forget set : rien n'entre dans la perte hors du
forget set, comme le sujet l'exige.

\paragraph{Cible $t$.} Un seul nombre pour toutes les images : la perte moyenne de l'original sur
une vue forte d'aire \num{0.50}, la doublure du milieu du banc, mesurée sur les 2800 images
entraînées : \num{0.449} (\chemin{experiments/unlearning/measurements/v3_targets.json}). Défaut
\code{dumb\_target\_loss} = \num{0.45} ; la spécification teste aussi \num{0.31} et \num{0.63}
(aires \num{0.60} et \num{0.45}).

\paragraph{Ce qu'elle ne fait pas, exprès.} Pas de distribution cible : la masse quitte la classe
prédite vers les classes que le gradient choisit. Pas d'arrêt : trois passes, le défaut commun
(taux \num{e-5}, \SI{10}{\percent} de montée puis cosinus). Ni pli ni sélection. Pronostic, non
mesuré : toutes les images finissent près de la même perte, alors que celles d'un gold sont
hétérogènes ; l'écart-type des pertes restera bas, et le terme d'écart-type de $S_{\text{ms}}$
\eqref{eq:ms} devrait borner le score.

\paragraph{Mise en œuvre et état.} \code{dumb} est une nouvelle méthode
(\chemin{src/hackathon/unlearning/methods.py}) avec un réglage nouveau,
\code{dumb\_target\_loss} ; les autres méthodes sont inchangées. La dernière ligne de l'historique
donne la divergence d'oubli exacte (\code{forget\_kl}), la moyenne et l'écart-type des pertes des
images entraînées (\code{forget\_loss\_mean}, \code{forget\_loss\_std}) et la part des images à la
cible (\code{reached}). Deux tests unitaires, sur un petit ViT, couvrent la reproductibilité, le
modèle d'entrée intact, la perte qui monte et la charnière plate au-dessus de la cible. La
spécification \code{v-dumb} a quatre bras ($t\in\{\num{0.31};\num{0.45};\num{0.63}\}$, puis
$t=\num{0.45}$ sans ancre) ; ses seuils de réfutation sont posés avant tout run dans
\chemin{experiments/unlearning/v-dumb.md} : réfutée si le meilleur bras reste sous la V1 moins
$2\sigma$ (\num{0.0577}) ou sous un retain de \num{0.8}, l'ancre jugée inutile si le bras sans elle
garde le même retain à $2\sigma$ près. Codée le 22 novembre ; mesurée le même jour, abandonnée
(\cref{sec:campagnevdumb}).

\subsection{V\_Fusion : SCRUB conforme ancré sur des images synthétiques}\label{sec:vfusion}
\paragraph{Pourquoi.} SCRUB conforme est la meilleure méthode conforme mesurée
(\num{0.1085 +- 0.0068}, retain \num{0.86}, \cref{tab:v4scrub}). Un cran de dose plus haut, il oublie
davantage (\num{0.121 +- 0.007}), mais son retain tombe à \num{0.73}, sous le plancher de \num{0.8}
fixé par le responsable. La campagne V4 dit d'où vient la limite. Son ancre, faite de vues fortes des
images oubliées, protège bien (à lr \num{2e-4}, retain \num{0.96} et KL d'audit \num{0.04}), mais
elle tient le contenu même de ces images et défait l'oubli aux faibles doses : à lr \num{e-4}, la KL
des images exactes ne vaut que \numrange{0.27}{0.31} fois celle de l'ancre externe. L'ancre externe
ne défait pas l'oubli, mais protège mal (retain \num{0.66} et KL d'audit \num{0.42} à lr \num{2e-4}).
La V5 a montré que la localité seule ne rapporte rien sur ce banc (\cref{sec:campagnev5}) : le but
n'est donc pas un oubli plus local, mais un retain tenu à une dose plus forte. Le pari, non mesuré :
une ancre faite d'images des bonnes classes, qui ne sont pas les images oubliées, garde la fonction
sans retenir ces images. Des images synthétiques tirées du seul modèle original restent dans le cadre
du sujet, comme le désapprentissage sans données qui tire du modèle lui-même ses images de
rétention~\cite{chundawat2023}.

\paragraph{Méthode.} SCRUB conforme est inchangé (\cref{sec:refsdetail}), sauf la source de ses
min-steps : une distillation $T^2\KL\bigl(p^{(2)}_\orig\,\|\,p^{(2)}_\theta\bigr)$ sur des images
synthétiques augmentées, au lieu de vues fortes du forget set. Une divergence, et non une entropie
croisée sur l'argmax : c'est l'unité du retain \eqref{eq:retain}, son gradient est nul à l'original,
et une entropie croisée ne donnerait presque aucune force sur des images que leur génération a rendues
sûres, sinon pour les rendre plus sûres encore, ce qui est aussi une dérive. L'alternance de SCRUB,
des pas qui effacent puis des pas qui réparent, est le schéma des méthodes analysées du concours
NeurIPS~\cite{triantafillou2024}, qui réparent sur les vraies données à garder ; ici, sur des images
synthétiques.

\paragraph{Les images synthétiques.} La génération suit Plug-In Inversion, qui inverse des ViT sans
régulariseur à régler pour chaque réseau, par des augmentations~\cite{ghiasi2022pii}. Chaque image part
du bruit à \SI{56}{px}, puis \num{112} et \num{224} (le zoom de l'article), et elle est optimisée
(Adam, taux \num{0.05}, 100 pas par étage) pour que l'original prédise sa classe sur quatre copies
augmentées (décalage circulaire d'au plus un patch, décalage de couleur par canal) ; le centrage de
l'article n'est pas repris. Au dernier étage, une image s'arrête dès que l'original donne
\num{0.9} à sa classe, le niveau des vues fortes (\cref{fig:calibration}) : à $T=2$, sa sortie doit
encore dessiner les classes fausses. Ni image externe, ni générateur, ni image du forget set ;
retournement, recadrage et flou ne servent jamais à la génération, pour que la porte a les teste sur
du neuf. Les images se rangent en trois lots, par classe : A, les classes des pseudo-étiquettes de
l'original sur le forget set, moins neuf classes mises à part ; B, les classes absentes du forget
set ; C, ces neuf classes (0, 6, 7, 8, 29, 30, 31, 33, 38), qu'une note de l'équipe suppose jamais
vues par le gold, hypothèse non mesurée : les tenir à l'original garderait un savoir que le gold n'a
pas, et elles n'entrent dans aucune ancre. Vingt images par classe ; les cinq dernières servent de
thermomètre et n'entrent dans aucune perte. Deux lots, parce que leurs risques sont opposés. B ne
partage aucune classe avec le forget set et ne peut pas en retenir les images, mais il protège là où
la fuite n'est pas : le fold d'audit, tiré des classes du forget set, bouge \num{0.85} fois autant que
les images oubliées à la dose retenue. A protège là où ça fuit, au risque de défaire une part de
l'oubli par le mode commun de la classe. Le banc tranche.

\paragraph{Portes.} Fixées avant toute lecture des résultats.
\begin{description}[leftmargin=1.2em,font=\sffamily\bfseries\color{bleu}]
  \item[a. Robustesse] Sous retournement, recadrage de 60 à \SI{90}{\percent} de l'aire et flou
    gaussien ($\sigma=1$), l'original garde la classe sur au moins \SI{80}{\percent} des images, pour
    chaque transformation.
  \item[b. Appartenance] Le jeton de classe à l'entrée de la tête est plus proche du centre de sa
    classe, calculé sur les vraies images du forget set, que de tout autre centre, pour au moins
    \SI{70}{\percent} des images. Pas mesurable pour le lot B, faute de vraies images de ses classes ;
    un prototype tiré des poids de la tête serait circulaire, la génération maximisant ce score.
  \item[c. Variété] Au plus \SI{5}{\percent} des paires d'une même classe sont aussi semblables
    qu'une image et son retournement (au-delà du cinquième centile de ces similarités).
  \item[c'. Distance au forget set] Au bloc du milieu, une image aussi proche d'une image du forget set
    que l'est, en médiane, une image du forget set de sa vue forte (aire \num{0.5}) est retirée. La
    porte échoue si plus de \SI{10}{\percent} du lot A part.
  \item[d. Réparation] Un bruit sur les poids, proportionnel à l'écart-type de chaque tenseur et
    calibré par bisection, fait changer d'avis l'original sur \SI{20}{\percent} des vues fortes
    (aire \num{0.5}) d'une moitié du forget set. Une distillation à $T=2$ sur les seules images
    d'ancre d'un lot (300 pas, lr \num{3e-5}) doit ramener l'accord à au moins \num{0.95}. Un témoin
    répare avec le même budget sur des vues fortes de l'autre moitié : s'il échoue, le budget est
    court et la porte ne dit rien du lot. Elle dit aussi si le lot B répare des images des classes du
    forget set, ce qu'il doit faire pour protéger le control set.
\end{description}

\paragraph{Protections des poids, en bras.} Le gel de couches, une pénalité L2 vers les poids
d'origine (L2-SP~\cite{li2018l2sp}), une pénalité pondérée par l'importance de chaque poids
(EWC~\cite{kirkpatrick2017ewc}, MAS~\cite{aljundi2018mas}) et l'ancre en KL forment une même famille,
de la plus grossière à la plus fine. À l'ordre deux autour de l'original, la KL de l'ancre vaut
$\tfrac12\,\Delta\theta^\top\mathcal I\,\Delta\theta$, où $\mathcal I$ est l'information de Fisher à
$T=2$ sur les images d'ancre : EWC en garde la diagonale, sans passe avant par pas, et agit aussi
pendant les pas de montée, où l'ancre de SCRUB est absente. Une L2 uniforme ne sait rien des directions
générales ; la sonde de la V5 a vu que ramener le modèle vers l'original ne fait que doser
(\cref{sec:v5}). Le gel des premiers blocs est déjà mesuré en V1, dans le mauvais sens : geler les 6
ou 10 premiers blocs, c'est-à-dire n'entraîner que les 6 ou 2 derniers, double la dérive d'audit
(\cref{tab:reglages}). Il reste un bras à faible dose de gel, l'embedding et les deux premiers blocs
(\code{trainable\_blocks=10}, sans code), à côté de deux bras à coder, EWC et L2. Pronostic, non
mesuré : le gel et la L2 ne déplacent que la dose.

\paragraph{Effacer puis réparer, dans l'autre ordre.} Une note de l'équipe propose d'effacer d'abord
par atténuation sélective des poids (SSD~\cite{foster2024ssd}, dans sa version sans étiquettes, qui
remplace la diagonale de Fisher par le gradient de la norme de la sortie~\cite{foster2024lfssd}),
puis de réparer sur les images synthétiques. Elle entre dans le même cadre comme un autre opérateur
d'oubli, après les bras d'ancre, avec deux réserves. Contraster l'importance sur le forget set et sur
des images synthétiques mesure aussi l'écart entre photos et images inversées : le contraste avec les
vues fortes du forget set, même contenu et même domaine, isole mieux la mémorisation. Et ses deux
réglages dosent l'oubli sans lien avec le score : la dose se choisira sous le plancher.

\paragraph{Mise en œuvre et état.} \chemin{src/hackathon/unlearning/synthetic.py} : génération par
tranches de classes (\code{synth-generate --shard k/n}), étiquetage par lot et par rôle, portes
(\code{synth-gates}, qui fusionne les tranches, écrit le masque \code{kept} et une grille d'aperçu).
\chemin{jureca/synthetic.sbatch} lance les quatre tranches sur les quatre A100 d'un même nœud, puis les
portes. Neuf tests unitaires, sur un petit ViT et des features construites à la main, couvrent les
lots, les tranches et leurs graines, la montée de la confiance, le plafond, la reproductibilité, le
retrait des copies du forget set, le comptage des quasi-copies, les centres de classe, le rapport
complet des portes, la lecture des fichiers et la ligne de commande. Les réglages sont proposés, non
mesurés. L'ancre synthétique de \code{scrub} est l'option \code{retain\_source=synthetic} : ses
min-steps tirent leurs lots des images gardées (\code{kept}) d'un fichier écrit par
\code{synth-gates} (\code{synthetic\_path}), hors thermomètre, dans les lots
\code{synthetic\_pools} (A, B ou AB), augmentées comme le proxy. Deux tests unitaires couvrent le
choix des images d'ancre et un run \code{scrub} sur cette ancre. Mesures :
\cref{sec:campagnevfusion}, où V\_Fusion est abandonnée sous sa première génération.

\subsection{V\_Dumb2 : la charnière de V\_Dumb dans le squelette de SCRUB conforme}\label{sec:vdumb2}
\paragraph{Pourquoi.} V\_Dumb (\cref{sec:vdumb}) visait un niveau de perte pour chaque image, sans
autre réglage ; sa campagne l'a abandonnée pour sa dynamique (\cref{sec:campagnevdumb}) : montée et
ancre sommées dans le même pas, d'où une falaise entre l'immobilité et l'effondrement ; une charnière
d'un seul côté, que le mouvement commun dépasse de 13 à 27 fois ; une force en perte qui s'éteint sur
une image mémorisée (\cref{sec:v51}, défaut 1) ; et un taux de \num{e-5}, trente fois sous la dose de
SCRUB conforme. SCRUB conforme a la dynamique qui marche (\cref{sec:campagnev4}) : une passe de
montée plafonnée, puis des passes d'ancre sur des vues fortes du forget set. Mais son plafond porte sur
la KL \emph{moyenne du lot} (2), loin de ce qu'il atteint (KL exacte \num{0.17} à lr \num{3e-4}) : il ne
lie jamais, et l'oubli se concentre sur peu d'images. Au registre, graine 0, le coefficient de variation
des pertes des images exactes vaut \num{7.5} à lr \num{3e-4} et \num{5.2} à \num{4e-4}.

\paragraph{Un seul niveau pour toutes les images.} Sans savoir sur quelles images le gold doute, et la
V5 a montré qu'aucune sélection tirée de vues de l'original ne le sait (\cref{sec:campagnev5}), un même
niveau pour toutes les images est le meilleur choix pour la composante des prédictions. Si $p$ varie
d'une image à l'autre et que l'élève ne peut pas les distinguer, la prédiction $q$ qui minimise
$\Exp\,\KL(p\|q)$ est la moyenne $\Exp p$ : c'est une propriété de la KL directe, et notre dérivation
l'applique à la masse de la pseudo-étiquette. La V5 a aussi mesuré que les doublures sont moins sûres
que l'original sur \emph{toutes} les images, typiques comprises. L'argument ignore le terme d'écart-type
de $S_{\text{ms}}$, qu'un niveau uniforme abaisse (pronostic de V\_Dumb). L'étape 0 mesure l'effet net
sur le produit.

\paragraph{Perte.} Les min-steps et le calendrier sont ceux de SCRUB conforme ; seul le max-step change
(\code{scrub\_target\_loss} $=t$). Avec $m_i=z_{\hat y}-\log\sum_{c\ne\hat y}e^{z_c}$ la marge de la
pseudo-étiquette (\cref{eq:v51m}) et $\nu(\ell)=-\log(e^{\ell}-1)$ la marge à laquelle l'entropie
croisée vaut $\ell$, chaque image $x_i$ du lot $B$ et sa copie altérée $a(x_i)$ (une des quatre
altérations du banc, tirée au hasard, avec la pseudo-étiquette de $x_i$) reçoivent
\begin{gather}
  \psi(x)=\tfrac{1}{\beta}\Bigl[\operatorname{softplus}\bigl(\beta(m(x)-\nu(t))\bigr)
  +\operatorname{softplus}\bigl(\beta(\nu(\kappa t)-m(x))\bigr)\Bigr],\label{eq:vdumb2}\\
  \mathcal{L}_{\max}=\frac{1}{1+w_a}\,\frac{1}{|B|}\sum_{i\in B}\bigl[\psi(x_i)+w_a\,\psi(a(x_i))\bigr],
  \notag
\end{gather}
avec $t=\num{0.45}$ (la cible de V\_Dumb), $\kappa=2$ (\code{band\_tolerance}), $\beta=10$ par logit
(\code{band\_sharpness}) et $w_a=1$ (\code{altered\_weight}). C'est la bande de la V5.1 en marge
(\cref{eq:v51band}), sans sélection. Une image mémorisée est poussée avec une force d'ordre 1, jusqu'à
une perte $t$. Entre $t$ et $\kappa t$, plus rien ne la pousse. Au-delà, elle est rappelée vers
l'original. Chaque pièce vient d'une mesure ou d'une analyse déjà faite. Le niveau vient de V\_Dumb.
La borne haute vient de la V5, où elle s'est montrée indispensable contre l'effondrement. L'écriture
en marge vient de la V5.1. L'alternance, les copies et l'ancre viennent de SCRUB conforme. La campagne
V5.1 a depuis mesuré la marge à pleine échelle (\cref{sec:campagnev51}) : sa force d'ordre 1 fait
dépasser la bande en 15 à 40 pas dès lr \num{3e-5}, et la V5.1 s'arrête alors au premier niveau sans en
garder. V\_Dumb2 l'emploie à \num{3e-4} et \num{4e-4}, sans arrêt dur ni niveaux : seuls le rappel et les
passes d'ancre la retiennent. C'est un risque mesuré, que les seuils de la spécification jugeront. Pari, non mesuré : à dose égale, l'arrêt par image dépense la dérive sur les images en retard
plutôt que sur celles qui ont déjà dépassé le gold. L'oubli serait plus uniforme, la composante des
prédictions plus haute, et le retain tenu à lr \num{4e-4}, où SCRUB conforme tombe à \num{0.73}.

\paragraph{Étape 0, sans entraînement.} \code{diagnose} rend désormais \code{level\_ceiling} :
l'oubli du banc ($S_{\text{ms}}\times S_{\text{pred}}$) sur les seules images exactes, quand chaque
marge est ramenée exactement dans la bande $[\nu(\kappa t),\nu(t)]$ par sa seule pseudo-étiquette. Les
classes fausses gardent la forme de l'original, rien d'autre ne bouge, et le retain vaut 1. C'est un
plafond, calculé pour $t\in\{\num{0.10};\num{0.20};\num{0.31};\num{0.45};\num{0.63};\num{1.00}\}$. La
référence de même portée est SCRUB conforme sur les images exactes (\code{bench.exact.*}, graine 0, pire
des six doublures) : \num{0.155} à lr \num{3e-4}, \num{0.198} à \num{4e-4}. Si le plafond à $t=\num{0.45}$
reste sous \num{0.155} sous le pire masquage, la spécification n'est pas lancée : même atteinte, la
bande oublierait moins que SCRUB conforme à sa dose retenue.

\paragraph{Spécification.} \code{v-dumb2} croise la dose (lr \num{3e-4} et \num{4e-4}) et le rappel
($\kappa=2$ ou $\kappa=10^6$, qui revient à ne rien rappeler), graine 0, $t=\num{0.45}$. Elle se compare,
appariée, à SCRUB conforme de même graine, déjà mesuré : \num{0.1151} (retain \num{0.868}) à
\num{3e-4}, \num{0.1279} (retain \num{0.734}) à \num{4e-4}. Seuils posés avant tout run, dans
\chemin{experiments/unlearning/v-dumb2.md}. La méthode est réfutée si aucun bras de retain supérieur ou
égal à \num{0.8} n'atteint \num{0.129} (SCRUB conforme apparié plus $2\sigma$, $\sigma=\num{0.0068}$), ou
si le bras à \num{4e-4} avec rappel passe sous un retain de \num{0.8}. Le mécanisme est réfuté si le
coefficient de variation des pertes du meilleur bras dépasse la moitié de celui de SCRUB conforme au
même taux. Le rappel est jugé inutile si les bras sans lui restent à $2\sigma$ du score et à \num{0.02}
du retain. Un positif se confirme sur les graines 1 et 2 (moyenne au-delà de \num{0.122}).

\paragraph{Mise en œuvre et état.} Une option de \code{scrub}, \code{scrub\_target\_loss}
(\chemin{src/hackathon/unlearning/methods.py}, \code{level\_band\_loss}), qui réutilise la bande et la
marge de la V5.1. Elle est neutre par défaut : sans elle, SCRUB rejoue la V4 au bit près (poids et
historique identiques sur le petit ViT des tests, sous trois réglages). Les deux rétropropagations
(image exacte, copie) restent séparées, comme dans SCRUB, pour garder le pic mémoire à une passe avant.
La dernière ligne de l'historique donne \code{forget\_kl}, la moyenne et l'écart-type des pertes, et
\code{in\_band}, la part des images dans $[t,\kappa t]$ à la fin. Quatre tests unitaires couvrent la
validation, la loi de force (ordre 1 sur une image mémorisée, nulle dans la bande, rappel au-delà, pas
de rappel sans plafond), la reproductibilité et le résumé de l'entraînement, et le plafond de l'étape 0.
Codée le 22 novembre ; aucun run.

\subsection{Pistes ouvertes}\label{sec:pistes}
La V2 (\cref{sec:trustv2}) puis la V3 (\cref{sec:trustv3}) traitent l'hétérogénéité des pertes par
la cible ; la V4 (\cref{sec:trustv4}) traite la localité, le défaut mesuré, par une ancre tirée du
forget set. Deux pistes de la revue restent à coder, si la V4 ne suffit pas. \textbf{Un masque
propre à l'instance} : ne mettre à jour que les poids dont l'importance pour l'image oubliée dépasse
nettement leur importance moyenne, à la manière de SSD~\cite{foster2024ssd}, la moyenne sur le
forget set tenant lieu de celle sur tout l'entraînement (ce qui suppose un forget set tiré au
hasard). Sa variante la plus simple, n'entraîner que les derniers blocs, a été réfutée par la
campagne (dérive doublée) ; un masque à la SalUn~\cite{fan2024salun}, sans contraste, risque le même
sort. \textbf{Le retrait de la composante commune des gradients}, projeter chaque mise à jour hors
du gradient moyen du forget set : une adaptation de l'équipe, sans appui publié dans ce cadre, et
écartée par l'étude de la V5, car elle interdit le décalage moyen que le score exige
(\cref{sec:v5}).

\section{Le banc d'essai local}\label{sec:banc}

Sans gold ni control set, le classement public est la seule mesure officielle ; il ne porte que sur
\SI{30}{\percent} des données et l'équipe a décidé de ne pas soumettre tant qu'une mesure locale ne
le justifie pas. Le banc (\chemin{src/hackathon/unlearning/bench.py}, protocole \code{bench-v4}
depuis la V3 ; la campagne du \cref{sec:campagne} est notée en \code{bench-v3}, qui n'a que la
première famille de doublures) applique les formules exactes \eqref{eq:ms}--\eqref{eq:retain} et remplace chaque pièce manquante
par une doublure déclarée (\cref{fig:banc,tab:doublures}).

\begin{figure}[ht]
  \centering
  \begin{tikzpicture}[font=\small\sffamily,
      boite/.style={draw=bleu,fill=brume,rounded corners=2pt,align=center,inner sep=4pt,
        minimum height=10mm,minimum width=24mm},
      cle/.style={boite,draw=sarcelle,fill=sarcelle!10},
      fleche/.style={-{Stealth[length=2.2mm]},thick,bleu!70}]
    \node[boite,minimum width=20mm] (forget) at (0,0) {Forget set\\\num{3500} images};
    \node[boite] (train) at (3.2,1.35) {Fold entraîné\\\SI{80}{\percent}};
    \node[boite] (audit) at (3.2,-1.35) {Fold d'audit\\\SI{20}{\percent}, graine 31};
    \node[boite] (methode) at (6.5,1.35) {Méthode\\\code{trust\_region}, \dots};
    \node[boite] (gold) at (6.5,0) {Doublures du gold\\recadrages, masquages};
    \node[boite,minimum width=20mm] (cand) at (9.6,1.35) {Candidat\\\code{.pt}};
    \node[cle,minimum width=20mm] (banc) at (9.6,0) {Banc\\\code{bench-v4}};
    \node[cle,minimum width=20mm] (score) at (12.6,0) {\code{score\_min}\\registre};
    \draw[fleche] (forget.north) |- (train.west);
    \draw[fleche] (forget.south) |- (audit.west);
    \draw[fleche] (train) -- (methode);
    \draw[fleche] (methode) -- (cand);
    \draw[fleche] (train.south) |- (gold.west);
    \draw[fleche] (cand) -- (banc);
    \draw[fleche] (gold) -- (banc);
    \draw[fleche] (audit.east) -| (banc.south);
    \draw[fleche] (banc) -- (score);
  \end{tikzpicture}
  \caption{Chaîne d'un run de réglage : le fold d'audit n'est jamais entraîné et joue le control
  set ; les doublures du gold sont calculées une fois par protocole et mises en cache.}
  \label{fig:banc}
\end{figure}

\begin{table}[ht]
  \centering\small
  \caption{Transposition locale des pièces du score officiel.}\label{tab:doublures}
  \begin{tabularx}{\linewidth}{@{}>{\raggedright}p{3.1cm}X>{\raggedright\arraybackslash}p{2.6cm}@{}}
    \toprule
    Pièce officielle & Doublure locale & Fiabilité\\
    \midrule
    Formules & exactes, \cref{eq:ms,eq:pred,eq:final,eq:retain} & exacte\\
    Control set & fold d'audit : \SI{20}{\percent} du forget set, stratifié par pseudo-étiquette,
      graine fixe, jamais entraîné & bonne ; mais ces images sont mémorisées\\
    Porte de \SI{70}{\percent} & accord avec l'original sur le fold d'audit & indirecte\\
    Vraies étiquettes & argmax de l'original (mémorisé, donc presque toujours juste) & bonne\\
    Gold (plusieurs modèles) & l'original sur 4 vues par doublure, deux familles : recadrages à 3
      forces (\cref{sec:calibration}) et masquage de patchs à 3 taux (\cref{tab:familles} ;
      \code{bench-v3} : recadrages seuls) & hypothèse ; les recadrages favorisent \code{gold\_soft},
      la V2 et la V3\\
    Copies altérées cachées & retournement, recadrage \SIrange{60}{90}{\percent}
      (\code{roi\_align}~\cite{he2017maskrcnn}), couleur, JPEG q50 & simulées\\
    \bottomrule
  \end{tabularx}
\end{table}

\paragraph{Deux familles de doublures.} Les recadrages et les cibles de \code{gold\_soft}, de la V2
et de la V3 sont tous \enquote{l'original sur des vues fortes} : le banc les favorise par
construction. Dans le modèle à logits libres de la note de recherche, une cible tirée des vues
gagnait un facteur 5 au sein de la seule famille des recadrages, mais seulement 2 à 3 face à une
autre famille. Le protocole \code{bench-v4} ajoute donc une seconde famille, que rien n'entraîne :
4 vues où une part des patchs $16\times16$ (un jeton du ViT) est remplacée par du gris. Les trois
taux sont choisis, comme les forces des recadrages, pour encadrer un gold à environ 78, 86 et
\SI{92}{\percent} de précision (\cref{tab:familles}) ; \SI{10}{\percent} de masquage
(accord \num{0.944}) et \SI{40}{\percent} (\num{0.661}) sortent de cette plage.

\begin{table}[ht]
  \centering\small
  \caption{Les six doublures du gold de \code{bench-v4} : l'original sur 4 vues par doublure,
  256 images du forget set, Apple M2, passes avant seules. Accord : part des vues dont l'argmax est
  celui de l'image exacte ; perte : entropie croisée à $T=1$ sur cette pseudo-étiquette, moyenne par
  vue ; CV : écart-type des pertes sur leur moyenne, moyennés vue par vue (le $\mathrm{CV}_\gold$
  de \eqref{eq:msbound}). Source : note de recherche du 21 novembre (masque tiré patch par patch),
  sauf (*) : mesuré avec le code du banc (\code{mask\_patches}, nombre exact de patchs par image) ;
  à \SI{20}{\percent} et \SI{30}{\percent}, la note donnait un accord de \num{0.868} et \num{0.775}.}\label{tab:familles}
  \sisetup{table-format=1.3}
  \begin{tabular}{@{}llSSS[table-format=1.2]@{}}
    \toprule
    Doublure & Paramètre & {Accord} & {Perte moy.} & {CV}\\
    \midrule
    $s_{40}$ & recadrage, aire $[0{,}40;0{,}50]$ & 0.782 & 0.954 & 2.30\\
    $s_{50}$ & recadrage, aire $[0{,}50;0{,}60]$ & 0.880 & 0.496 & 3.14\\
    $s_{60}$ & recadrage, aire $[0{,}60;0{,}70]$ & 0.916 & 0.333 & 3.77\\
    $m_{30}$ & \SI{30}{\percent} des patchs masqués (*) & 0.766 & 1.034 & 2.34\\
    $m_{20}$ & \SI{20}{\percent} des patchs masqués (*) & 0.856 & 0.472 & 2.63\\
    $m_{15}$ & \SI{15}{\percent} des patchs masqués (*) & 0.917 & 0.309 & 3.28\\
    \bottomrule
  \end{tabular}
\end{table}

\paragraph{Chiffre suivi.} \code{score\_min} est le score \eqref{eq:final} sous la \emph{pire} des
six doublures (trois en \code{bench-v3}) : il ne monte que si un candidat progresse quelle que soit l'hypothèse sur le gold,
dans les deux familles ; \code{score\_min\_family} donne le pire cas de chaque famille. L'identité
vaut 0 ; un candidat qui atteindrait la doublure vaudrait 1.

\paragraph{Dérive hors score.} Le fold d'audit est mémorisé par l'original : il peut sous-estimer la
dérive sur des images jamais vues. Le banc rapporte donc aussi, hors score, $\KL(p_\orig\|p_\sub)$ à
$T=2$ sur une vue forte (aire $0{,}50$, graine fixe) de chaque image du fold d'audit, des entrées
\enquote{comme inédites} (\code{audit\_view\_kl}). Prévue comme seuil de décision avant la campagne,
elle sert désormais à juger une ancre de rétention : la V1 déplace déjà le fold d'audit exact de
\num{0.23} à $\varepsilon=\num{0.15}$ (\cref{sec:campagne}). Coût : une passe du candidat sur le fold d'audit.


\paragraph{Reproductibilité.} Chaque rapport porte l'identifiant complet du protocole (versions,
tailles, graines, forces, altérations) ; deux chiffres ne se comparent qu'à protocole égal. Les
sorties de l'original sont mises en cache par protocole (écriture atomique, partagée par les
variantes parallèles), et la commande \code{bench --candidate <fichier>} recalcule n'importe quel
fichier : sur un essai, le recalcul a redonné exactement les chiffres du run. Seuls les runs sur le
forget set complet sont classés (\code{bench.*}) ; un essai réduit écrit \code{bench\_smoke.*}.
Depuis \code{bench-v4}, chaque rapport porte aussi le numéro du protocole
(\code{bench.protocol\_version}), et le tableau comparatif fait une section par protocole : un
chiffre \code{bench-v3} n'est jamais classé avec un chiffre \code{bench-v4}.

\paragraph{Bases de comparaison.} Deux expériences encadrent toute idée : les références
\code{gold\_soft} et \code{scrub} notées par le même protocole, et le bruit d'entraînement mesuré
sur trois graines. Une différence ne comptera que si elle dépasse deux écarts-types de ce bruit.
Pour juger la V2 et la V3, ces références (la V1 à $\varepsilon=\num{0.15}$, \code{scrub} et
\code{gold\_soft} dosés, le bruit de graine) doivent d'abord être recalculées en \code{bench-v4}.

\section{Protocole expérimental et outillage}\label{sec:protocole}
\paragraph{Expériences déclarées.} Une expérience est une spécification TOML
(\chemin{experiments/unlearning/<id>.toml}) : une commande, un balayage cartésien de quatre
variantes au plus, des arguments propres au local et à JURECA, la métrique qui la classe.
\code{make hk-run SPEC=<id>} l'exécute, en local variante après variante, ou sur JURECA à raison
d'une variante par A100 et d'un seul nœud pour toute l'équipe, en file derrière ses autres jobs.
Aucun run n'envoie rien au serveur.

\paragraph{Hypothèses posées avant.} Chaque spécification a sa fiche (\code{<id>.md}), qui fixe
l'hypothèse et les critères de réfutation \emph{avant} le lancement ; la campagne a sa fiche propre
(\chemin{experiments/unlearning/trust-region-campaign.md}), complétée étape par étape, chaque étape
étant posée avant d'être lancée.

\paragraph{Traçabilité.} Chaque run écrit son dossier (copie figée de la spécification, commandes
exactes, journaux, fichier produit) et une ligne au registre versionné
(\chemin{experiments/registry.jsonl}) : identifiant, commit, drapeau \code{dirty} (code modifié par
rapport au commit, registre et documentation exclus), paramètres et toutes les métriques
numériques du run. Les 84 runs des deux campagnes ont tourné sur des commits publiés, aucun avec du
code modifié.
\code{make hk-unl-scores} en tire \chemin{experiments/unlearning/SCOREBOARD.md} : chaque run, et
une synthèse par réglage avec moyenne et écart-type des graines, une section par protocole du
banc.

\paragraph{Déterminisme.} À graine égale, un run se reproduit : le run de calibration de
\code{gold\_soft} et sa reprise par le lab ont donné les mêmes pertes à chaque pas et le même banc ;
après la fusion avec la V2, le run de référence de la V1 redonne exactement \num{0.0647}
(\cref{sec:fusion}). Tout l'écart entre deux runs d'un même réglage vient donc de la graine.

\paragraph{Machine.} Les runs locaux passent par une garde (\chemin{scripts/gpu_guard.sh}) : un
seul job à la fois sur la carte, refus de démarrer sous \SI{9}{\giga\byte} de mémoire libre,
suspension à \SI{83}{\celsius}, télémétrie de chaque job (\cref{sec:couts}).

\paragraph{Règles de décision.} On classe sur \code{score\_min}, porte comprise ; une différence ne
compte que si elle dépasse deux écarts-types du bruit entre graines du même réglage, et l'on
compare une idée à des références réglées, pas à leurs défauts.

\section{Résultats à ce jour}\label{sec:resultats}

\begin{avertissement}
\textbf{Portée.} Deux familles de chiffres. Les essais réduits (Apple M2 : 52 images entraînées,
12 en audit, quelques dizaines de pas) valident la chaîne. Les campagnes à pleine échelle
(RTX 5080 : 2800 images entraînées, 700 en audit ; 174 runs d'entraînement, plus l'étape 0 de la V5,
sans entraînement) mesurent les méthodes sur le banc :
\code{bench-v3} (trois doublures par recadrage) pour la V1, les références et la V2,
\code{bench-v4} (six doublures, recadrage et masquage) pour la V3, la V4, la V5, V\_Dumb et pour
les points retenus de toutes les versions, recalculés. Une graine par réglage, trois pour les points clés (notés $\pm$
un écart-type). Deux protocoles ne se comparent jamais entre eux. Ce sont des distances à une
doublure du gold, pas le score du serveur, et rien n'a été soumis.
\end{avertissement}

\subsection{Vue d'ensemble : de la V1 à V\_Fusion}\label{sec:vue}
Cette section donne le chemin et la comparaison finale ; les sections suivantes détaillent chaque
étape. Chaque campagne a posé ses hypothèses et ses seuils avant de lancer ses runs
(\cref{sec:protocole}).

\begin{enumerate}[leftmargin=1.6em,itemsep=2pt]
  \item \textbf{La chaîne tient} (essais réduits) : à petite échelle, l'oubli reste propre à chaque
    image.
  \item \textbf{V1 : l'arrêt décide} (\cref{sec:campagne}). Testé tous les 10 pas, le rayon ne
    contrôle rien ; testé à chaque pas, il donne un score positif, avec un plateau pour
    $\varepsilon\in[\num{0.1};\num{0.2}]$. Le taux, le warmup, la température et le lot sont sans
    effet ; à pleine échelle, l'oubli n'est plus local (ratio \num{1.4}).
  \item \textbf{Les références dosées encadrent} (\cref{tab:refs}) : \code{scrub} dosé fait mieux
    en moyenne, mais il est bruité et non conforme au sujet (images proxy, \cref{sec:pistes}).
  \item \textbf{V2 : la direction de la cible compte} (\cref{sec:campagnev2}). Une cible tirée des
    vues, à basse température ($\tau=1$), bat la V1 de \SI{30}{\percent} sous recadrage, à quantité
    d'oubli égale ; trop basse ($\tau\le\num{0.5}$), le rayon n'est plus atteint.
  \item \textbf{\code{bench-v4} corrige la mesure} (\cref{tab:v4}). Sous le masquage, famille de
    doublures étrangère à toute cible, le gain de la V2 tombe à \SI{12}{\percent} : une partie du gain
    venait de la famille de vues commune à la cible et aux doublures.
  \item \textbf{V3 : la cible sature} (\cref{sec:campagnev3}). La cible calibrée fait jeu égal avec la
    V2 ; comme son rayon est tiré de la cible, tous ses réglages agissent par le rayon, dont
    l'optimum retombe vers \numrange{0.15}{0.18}. Le mécanisme attendu (dispersion des pertes) n'est
    pas observé : le réseau la fixe.
  \item \textbf{Le levier suivant est la localité} : les trois versions abîment
    autant le fold d'audit (ratio \numrange{1.3}{1.4} ; \numrange{0.86}{0.91} fois l'image exacte,
    \cref{sec:diagnostic}) ; c'est là que \code{scrub} se distingue, et seule une ancre conforme est
    admise. La V4 (\cref{sec:trustv4}) en propose une tirée du forget set lui-même.
  \item \textbf{V4 : SCRUB conforme passe en tête, la localité reste à obtenir}
    (\cref{sec:campagnev4}). Ancré sur des vues fortes du forget set et redosé (lr \num{3e-4}),
    SCRUB bat la V3 de \SI{36}{\percent} sous les deux familles, apparié sur trois graines : c'est la
    meilleure méthode conforme mesurée. Il ne gagne pas par la localité (\num{0.85} fois l'image
    exacte, comme la V3). L'ancre croisée freine l'oubli partout sans le localiser ; la contrainte
    active ne lie pas, et poursuivre la V3 sur tout son calendrier divise son score par deux.
  \item \textbf{V5 : la localité est obtenue, et elle ne paie pas} (\cref{sec:campagnev5}). Sa
    sélection échoue avant tout entraînement : les images où l'original doute sur des vues par
    recadrage ne sont pas celles où se trompent les doublures par masquage, et les deux familles de
    doublures ne s'accordent pas sur les images difficiles. Entraînée comme mesure de mécanisme, la
    bande est la plus locale des méthodes mesurées à KL exacte égale (\cref{fig:v5localite}), mais son
    score est nul : la composante des prédictions demande le mouvement commun que la localité refuse.
  \item \textbf{V\_Dumb : la méthode la plus simple n'a pas de dose} (\cref{sec:campagnevdumb}). Une
    montée de la perte coupée à une cible, plus l'ancre de SCRUB conforme sur les vues, sommées dans
    chaque pas : avec l'ancre, le modèle ne bouge pas ; un peu moins d'ancre, et il s'effondre. La
    charnière ne borne rien. Les cibles des V1 à V3 et l'alternance de SCRUB valent donc quelque chose
    sur ce banc.
  \item \textbf{V5.1 : corriger la loi de force ne suffit pas} (\cref{sec:campagnev51}). Les trois
    défauts de la V5 sont réels, mais la force d'ordre 1 de la marge fait dépasser la bande en
    quelques pas. Au seul taux où la V5.1 garde une dose, elle score comme la V3, sous le plancher de
    retain et sans oubli local.
  \item \textbf{V\_Fusion : une ancre synthétique ne voit pas les dégâts} (\cref{sec:campagnevfusion}).
    Les images tirées de l'original sont saturées : le modèle cassé par la montée de SCRUB reste
    d'accord avec l'original sur elles, et les quatre bras d'ancre s'effondrent. L'ancre de vues du
    forget set reste la seule qui tienne.
\end{enumerate}

\begin{table}[ht]
  \centering\small
  \caption{Comparaison finale, même protocole (\code{bench-v4}, trois graines appariées, RTX 5080 ;
  SCRUB conforme : \code{retain\_source=forget\_views}, \cref{sec:campagnev4}).
  Pire recadrage et pire masquage : pire doublure de chaque famille ; \code{score\_min} : pire
  des six. KL audit : $\KL(p_\orig\|p_\sub)$ à $T=2$ sur le fold d'audit ; ratio : mouvement des
  images oubliées sur celui du fold d'audit. La V5, abandonnée (\code{score\_min} \numrange{0.000}{0.005}
  à son réglage par défaut, deux graines), V\_Dumb (au mieux \num{0.019}, une graine), la V5.1
  (\num{0.080} à retain \num{0.725}, sous le plancher, une graine) et V\_Fusion (effondrée) n'ont pas
  de ligne (\cref{sec:campagnev5,sec:campagnevdumb,sec:campagnev51,sec:campagnevfusion}).}\label{tab:synthese}
  \setlength{\tabcolsep}{3.2pt}
  \begin{tabular}{@{}>{\raggedright}p{4.2cm} ccc S[table-format=1.3] S[table-format=1.3] S[table-format=1.2] c@{}}
    \toprule
    Méthode (réglage retenu) & \code{score\_min} & Pire recadr. & Pire masq. & {Retain} & {KL audit}
      & {Ratio} & Conforme\\
    \midrule
    V1, $\varepsilon=\num{0.15}$ & \num{0.0637 +- 0.0025} & \num{0.0637} & \num{0.0702} & 0.768 & 0.228
      & 1.36 & oui\\
    V2, $\tau=1$, aire \num{0.45}, 3 vues, $\varepsilon=\num{0.15}$ & \num{0.0786 +- 0.0030} &
      \num{0.0831} & \num{0.0786} & 0.778 & 0.225 & 1.36 & oui\\
    V3, $b=0$, $\lambda=\num{0.25}$, $\rho=\num{0.9}$, aire \num{0.5}, 4 vues & \num{0.0798 +- 0.0012}
      & \num{0.0875} & \num{0.0798} & 0.744 & 0.261 & 1.29 & oui\\
    \textbf{SCRUB conforme} (V4, vues du forget set), lr \num{3e-4} & \num{0.1085 +- 0.0068}
      & \num{0.1155} & \num{0.1085} & 0.859 & 0.142 & 1.24 & oui\\
    \midrule
    \code{scrub}, lr \num{e-4}, images proxy (référence) & \num{0.098 +- 0.022} & \num{0.098} & \num{0.101} & 0.931
      & 0.067 & 1.96 & non\\
    \bottomrule
  \end{tabular}
\end{table}

\begin{figure}[ht]
  \centering
  \begin{tikzpicture}
    \begin{axis}[width=0.86\linewidth,height=5.6cm,ybar=1pt,bar width=9pt,
        symbolic x coords={V1,V2,V3,scrubfv,scrub},xtick=data,
        xticklabels={V1,V2,V3,{SCRUB\\conforme},{\code{scrub}\\(non conforme)}},
        xticklabel style={align=center},
        ymin=0,ymax=0.13,ylabel={score},enlarge x limits=0.18,
        yticklabel style={/pgf/number format/.cd,fixed,precision=2,use comma},
        legend columns=3,legend style={at={(0.5,1.03)},anchor=south,font=\small,draw=none},
        grid=major,grid style={gray!18},axis line style={gray!60},tick style={gray!60},
        label style={font=\small},tick label style={font=\small},
        error bars/y dir=both,error bars/y explicit]
      \addplot[fill=bleu!70,draw=bleu] coordinates
        {(V1,0.0637) +- (0,0.0025) (V2,0.0831) +- (0,0.0042) (V3,0.0875) +- (0,0.0022) (scrubfv,0.1155) +- (0,0.0067) (scrub,0.0978) +- (0,0.022)};
      \addlegendentry{pire recadrage}
      \addplot[fill=sarcelle!70,draw=sarcelle] coordinates
        {(V1,0.0702) +- (0,0.0013) (V2,0.0786) +- (0,0.0030) (V3,0.0798) +- (0,0.0012) (scrubfv,0.1085) +- (0,0.0068) (scrub,0.1008) +- (0,0.0234)};
      \addlegendentry{pire masquage}
      \addplot[fill=ambre!80,draw=ambre!80!black] coordinates
        {(V1,0.0637) +- (0,0.0025) (V2,0.0786) +- (0,0.0030) (V3,0.0798) +- (0,0.0012) (scrubfv,0.1085) +- (0,0.0068) (scrub,0.0978) +- (0,0.022)};
      \addlegendentry{\code{score\_min} (pire des six)}
    \end{axis}
  \end{tikzpicture}
  \caption{Les versions retenues sous les deux familles de doublures de \code{bench-v4} (moyenne et
  écart-type sur trois graines). La V2 et la V3 dépassent la V1 sous les deux familles ; leur
  avance est plus nette sous le recadrage, famille que leurs cibles partagent avec les doublures.
  SCRUB conforme (V4) les dépasse sous les deux familles, sans écart propre au recadrage. La V5,
  V\_Dumb, la V5.1 et V\_Fusion, abandonnées, n'y figurent pas.}
  \label{fig:synthese}
\end{figure}

\subsection{Essais réduits}
Le \cref{tab:smoke} note quatre candidats avec \code{bench-v3}. Les scores sont faibles (\numrange{0.03}{0.05}
sur 1), ce qui est attendu : quelques dizaines de pas sur 52 images déplacent peu le modèle. Cinq
constats sont plus instructifs.
\begin{itemize}[leftmargin=1.2em]
  \item \textbf{Le retain est intact} (\numrange{0.996}{0.998}) : la divergence sur le fold d'audit
    reste 40 à 84 fois plus faible que sur les images oubliées, ce que prédisait la
    \cref{sec:local}.
  \item \textbf{\code{gold\_soft} devance légèrement \code{trust\_region}} à pas égaux, comme on
    pouvait l'attendre d'un banc dont la doublure ressemble à sa cible. Sur une seule graine et à
    cette échelle, l'écart n'est pas interprétable.
  \item \textbf{Le transfert aux copies altérées est partiel} : sur les retournements et
    recadrages, la divergence vaut entre la moitié et les deux tiers de celle des images exactes,
    tout près du seuil de réfutation fixé pour \code{trust\_region} (la moitié). C'est le premier
    point à surveiller à pleine échelle.
  \item \textbf{La V2 fait jeu égal avec la V1 en 40 pas} (\code{score\_min} \num{0.0289}
    contre \num{0.0287}, V1 recalculée par le même protocole ; un peu moins de divergence sur le
    fold d'audit). Son léger avantage sous $s_{60}$ vient de la composante de prédiction (\num{0.29}
    contre \num{0.26}), pas de celle des pertes, où elle est un peu en dessous sous les trois
    doublures (\num{0.18} contre \num{0.19} sous $s_{60}$) : l'inverse du mécanisme attendu, à
    vérifier à pleine échelle. À ce taux, aucune des deux n'approche son rayon (divergence
    \num{0.036} contre \num{0.041}, pour $\varepsilon\ge\num{0.2}$) : les quatre variantes de
    \code{trust-region-views} ont fait les mêmes 40 pas et donné le même chiffre. L'essai valide la
    chaîne de la V2 (cible, contrôle de faisabilité, arrêt, banc), pas son intérêt. Sur ses 52
    images, la cible est à une divergence de \num{0.94}, juste sur \SI{92}{\percent} d'entre elles
    (2 vues en essai réduit, contre 3 à pleine échelle).
  \item \textbf{La V3 passe la chaîne en \code{bench-v4}} (\cref{tab:smokev4}), sur une seule
    configuration ($\lambda=b=0{,}25$, copies douces sur la moitié de chaque lot, 40 pas,
    \SI{38}{\second} d'entraînement). Le rayon tiré de la cible vaut \num{0.156} ; il n'est pas
    atteint (divergence \num{0.033}). La dérive sur des vues du fold d'audit (\num{0.0086}) est
    15 fois celle du fold exact (\num{0.0006}). Le premier essai a
    échoué sur le GPU d'Apple, sans float64 : la cible est désormais calculée sur le processeur.
    L'identité vaut exactement 0 sous ce protocole. Rien n'est interprétable à cette échelle.
\end{itemize}

\begin{table}[ht]
  \centering\small
  \caption{Banc \code{bench-v3} sur les essais réduits (Apple M2). $s_{40}$, $s_{50}$, $s_{60}$ :
  score sous chaque doublure du gold ; KL : $\KL(p_\orig\|p_\sub)$ à $T=2$ ; \enquote{ratio} :
  divergence moyenne sur les images oubliées et leurs copies, divisée par celle du fold d'audit.}
  \label{tab:smoke}
  \sisetup{table-format=1.3}\setlength{\tabcolsep}{4.2pt}
  \begin{tabular}{@{}lc SSSS S[table-format=1.4] S[table-format=2.0] S[table-format=1.3]S[table-format=1.3]@{}}
    \toprule
    Candidat & Pas, taux & {\code{score\_min}} & {$s_{40}$} & {$s_{50}$} & {$s_{60}$} & {KL audit}
      & {Ratio} & {KL exact} & {KL retourn.}\\
    \midrule
    identity & -- & 0 & 0 & 0 & 0 & 0 & {--} & 0 & 0\\
    trust\_region & 7, \num{5e-5} & 0.036 & 0.036 & 0.071 & 0.047 & 0.0019 & 84 & 0.176 & 0.093\\
    gold\_soft & 7, \num{5e-5} & 0.051 & 0.051 & 0.111 & 0.055 & 0.0036 & 48 & 0.183 & 0.091\\
    trust\_region & 40, \num{e-5} & 0.029 & 0.029 & 0.059 & 0.048 & 0.0024 & 41 & 0.093 & 0.062\\
    trust\_region V2 & 40, \num{e-5} & 0.029 & 0.029 & 0.058 & 0.051 & 0.0018 & 52 & 0.089 & 0.057\\
    \bottomrule
  \end{tabular}
\end{table}

\begin{table}[ht]
  \centering\small
  \caption{Banc \code{bench-v4} sur l'essai réduit de la V3 (Apple M2, 52 images entraînées, 12 en
  audit). Pire recadrage et pire masquage : \code{score\_min\_family} ; KL vue d'audit :
  \code{audit\_view\_kl}, hors score.}\label{tab:smokev4}
  \sisetup{table-format=1.4}
  \begin{tabular}{@{}lc SSS S S[table-format=1.4] S[table-format=3.0]@{}}
    \toprule
    Candidat & Pas, taux & {\code{score\_min}} & {Pire recadr.} & {Pire masq.} & {KL audit}
      & {KL vue d'audit} & {Ratio}\\
    \midrule
    identity & -- & 0 & 0 & 0 & 0 & 0 & {--}\\
    trust\_region V3 & 40, \num{e-5} & 0.0195 & 0.0224 & 0.0195 & 0.0006 & 0.0086 & 143\\
    \bottomrule
  \end{tabular}
\end{table}

\subsection{Essai intermédiaire de la V3}\label{sec:intermediaire}
Écrit sans connaître la campagne V3 (\cref{sec:campagnev3}), qui l'a depuis mesurée à pleine
échelle, un essai entre les deux échelles (Apple M2, 22 novembre) : 400 images du forget set, dont
320 entraînées et 80 en audit, lots de 8, taux \num{3e-5}, la spécification d'origine de la V3
($\lambda=b=\num{0.25}$, copies douces sur la moitié des lots), une graine, banc \code{bench-v4}
(\code{bench\_smoke}, non classé). Le rayon, testé tous les 4 pas, est dépassé : la divergence passe
de \num{0.065} à \num{0.244} en quatre pas et finit à \num{0.247} pour un rayon de \num{0.169}.
\code{score\_min} vaut \num{0.088}, retain \num{0.833}. Seul enseignement retenu : à cette échelle
déjà, le fold d'audit bouge à \num{0.52} fois l'image exacte (KL \num{0.171} contre \num{0.327}), et
sur une vue forte des mêmes images la dérive atteint \num{0.408} (la saturation du softmax sur les
images mémorisées en explique une partie ; le sens du changement n'est pas mesuré). Ces chiffres
ne se comparent à aucun autre run.

\subsection{Campagne V1 à pleine échelle}\label{sec:pleine}\label{sec:campagne}
Chaque étape a été posée, hypothèse et critère de réfutation compris, avant d'être lancée
(\chemin{experiments/unlearning/trust-region-campaign.md}) ; les 54 runs sont au registre, et
\chemin{experiments/unlearning/SCOREBOARD.md} (\code{make hk-unl-scores}) les compare tous.

\paragraph{Un arrêt trop grossier.} Testé tous les 10 pas, le rayon ne contrôle rien :
$\varepsilon=\num{0.2}$, \num{0.4} et \num{0.6} s'arrêtent au même pas (90) et donnent le même
modèle, car la divergence passe de \num{0.07} à \num{0.65} en dix pas, à la fin du warmup. Ce
modèle a un \code{score\_min} nul (retain \num{0.52}). Les références à leurs réglages par défaut
encadrent mal : \code{gold\_soft} va beaucoup trop loin (retain \num{0.40}, score nul),
\code{scrub} ne bouge presque pas (\num{0.0025 +- 0.0003} sur trois graines).

\paragraph{Le rayon, testé à chaque pas.} La divergence finale suit alors $\varepsilon$ (1,04 à
1,24 fois $\varepsilon$ dès \num{0.05}) et le score devient positif sous les trois doublures
(\cref{fig:rayon,tab:rayon}). Il forme un plateau entre \num{0.1} et \num{0.2}, confirmé sur trois
graines (\cref{tab:refs}). Tous les arrêts tombent entre les pas 72 et 92, pendant le warmup : la
divergence passe de \num{0.02} à \num{0.94} en vingt pas. Le ratio de mouvement reste entre
\num{1.2} et \num{1.8} : l'oubli n'est pas localisé, le fold d'audit dérive avec les images
entraînées, et le score reste positif parce que le gain sur l'oubli l'emporte.

\begin{figure}[ht]
  \centering
  \begin{tikzpicture}
    \begin{axis}[width=0.86\linewidth,height=6.2cm,xmode=log,
        xlabel={rayon $\varepsilon$ (échelle logarithmique)},ylabel={score},
        xmin=0.008,xmax=1,ymin=0,ymax=0.125,
        yticklabel style={/pgf/number format/.cd,fixed,precision=2,use comma},
        xtick={0.01,0.02,0.05,0.1,0.2,0.5,1},
        xticklabels={{0,01},{0,02},{0,05},{0,1},{0,2},{0,5},{1}},
        legend columns=3,legend style={at={(0.5,1.03)},anchor=south,font=\small,draw=none},
        grid=major,grid style={gray!18},axis line style={gray!60},tick style={gray!60},
        label style={font=\small},tick label style={font=\small}]
      \fill[sarcelle!12] (axis cs:0.008,0.076) rectangle (axis cs:1,0.120);
      \draw[sarcelle!80!black,dashed] (axis cs:0.008,0.098) -- (axis cs:1,0.098);
      \node[font=\scriptsize\sffamily,text=sarcelle!80!black,anchor=south east] at (axis cs:0.98,0.098)
        {\code{scrub} dosé (lr \num{e-4}), moyenne $\pm$ écart-type};
      \addplot[bleu,very thick,mark=*,mark size=1.8pt] coordinates
        {(0.01,0.0236) (0.02,0.0328) (0.05,0.0511) (0.1,0.0594) (0.15,0.0647) (0.2,0.0595) (0.3,0.0338) (0.5,0.0018) (0.8,0)};
      \addlegendentry{\code{score\_min}}
      \addplot[ambre,thick,mark=square*,mark size=1.4pt] coordinates
        {(0.01,0.0236) (0.02,0.0328) (0.05,0.0511) (0.1,0.0700) (0.15,0.0820) (0.2,0.0866) (0.3,0.0919) (0.5,0.0887) (0.8,0.0660)};
      \addlegendentry{$s_{40}$}
      \addplot[brique,thick,mark=triangle*,mark size=1.8pt] coordinates
        {(0.01,0.0257) (0.02,0.0354) (0.05,0.0525) (0.1,0.0643) (0.15,0.0647) (0.2,0.0595) (0.3,0.0338) (0.5,0.0018) (0.8,0)};
      \addlegendentry{$s_{60}$}
    \end{axis}
  \end{tikzpicture}
  \caption{\code{trust\_region} sur le banc à pleine échelle, rayon testé à chaque pas (graine 0,
  RTX 5080). $s_{40}$ et $s_{60}$ : score sous la doublure la plus lâche et la plus stricte ;
  \code{score\_min} : la pire des trois. En bande, \code{scrub} dosé sur trois graines.}
  \label{fig:rayon}
\end{figure}

\begin{table}[ht]
  \centering\small
  \caption{Balayage du rayon, testé à chaque pas (graine 0). \enquote{KL/$\varepsilon$} : divergence
  finale sur toutes les images entraînées, rapportée au rayon ; \enquote{ratio} : mouvement des
  images oubliées et de leurs copies sur celui du fold d'audit.}
  \label{tab:rayon}
  \sisetup{table-format=1.3}\setlength{\tabcolsep}{4.5pt}
  \begin{tabular}{@{}S[table-format=1.2] S[table-format=1.2] S[table-format=2.0] S[table-format=1.4] SSS SS[table-format=1.1]@{}}
    \toprule
    {$\varepsilon$} & {KL/$\varepsilon$} & {Pas} & {\code{score\_min}} & {$s_{40}$} & {$s_{50}$}
      & {$s_{60}$} & {Retain} & {Ratio}\\
    \midrule
    0.01 & 1.92 & 72 & 0.0236 & 0.024 & 0.044 & 0.026 & 0.961 & 1.8\\
    0.02 & 1.53 & 75 & 0.0328 & 0.033 & 0.051 & 0.035 & 0.939 & 1.7\\
    0.05 & 1.24 & 79 & 0.0511 & 0.051 & 0.056 & 0.053 & 0.892 & 1.5\\
    0.10 & 1.12 & 82 & 0.0594 & 0.070 & 0.059 & 0.064 & 0.830 & 1.4\\
    0.20 & 1.07 & 85 & 0.0595 & 0.087 & 0.073 & 0.060 & 0.739 & 1.3\\
    0.30 & 1.11 & 87 & 0.0338 & 0.092 & 0.073 & 0.034 & 0.665 & 1.3\\
    0.50 & 1.04 & 89 & 0.0018 & 0.089 & 0.050 & 0.002 & 0.568 & 1.2\\
    0.80 & 1.18 & 92 & 0 & 0.066 & 0.001 & 0 & 0.421 & 1.2\\
    \bottomrule
  \end{tabular}
\end{table}

\paragraph{Les autres réglages ne comptent pas.} Autour de $\varepsilon=\num{0.15}$, où trois graines
donnent \num{0.0637 +- 0.0025}, chaque réglage a été balayé seul (\cref{tab:reglages}). Le taux,
le warmup, la température de la cible et la taille de lot restent dans le bruit : à divergence
égale, ni la vitesse ni la direction de la cible ne changent le score de façon mesurable. Seul
l'ensemble des poids entraînés compte, et dans le sens inattendu : n'entraîner que les derniers
blocs \emph{double} la dérive sur le fold d'audit, et la tête seule n'atteint jamais le rayon.

\begin{table}[ht]
  \centering\small
  \caption{Réglages balayés un à un à $\varepsilon=\num{0.15}$ (une graine chacun ; la référence,
  en gras, a trois graines). Valeur du réglage $\to$ \code{score\_min}.}
  \label{tab:reglages}
  \begin{tabularx}{\linewidth}{@{}>{\raggedright}p{2.9cm}X@{}}
    \toprule
    Réglage & Valeurs $\to$ \code{score\_min}\\
    \midrule
    Taux & \num{3e-6} $\to$ \num{0.064} (170 pas) ; \textbf{\num{e-5}} $\to$ \num{0.064} (85 pas) ;
      \num{3e-5} $\to$ \num{0.067} (50 pas) ; \num{e-4} $\to$ \num{0.062} (29 pas)\\
    Part de warmup & 0 $\to$ \num{0.065} ; \num{0.03} $\to$ \num{0.066} ; \textbf{\num{0.1}} $\to$
      \num{0.064} ; \num{0.3} $\to$ \num{0.065}\\
    Température $\tau$ & 4 $\to$ \num{0.065} ; 5 $\to$ \num{0.065} ; \textbf{6} $\to$ \num{0.064} ;
      8 $\to$ \num{0.065} ; 12 $\to$ \num{0.065}\\
    Blocs entraînés & 0 (tête) $\to$ \num{0.0006} ; 2 $\to$ \num{0.025} ; 6 $\to$ \num{0.025} ;
      \textbf{tous} $\to$ \num{0.064} (KL audit \num{0.20}, contre \num{0.44} à \num{0.51})\\
    Taille de lot & 8 $\to$ \num{0.063} ; 16 $\to$ \num{0.066} ; \textbf{32} $\to$ \num{0.064} ;
      64 $\to$ \num{0.067}\\
    \bottomrule
  \end{tabularx}
\end{table}

\paragraph{Contre des références dosées.} Comparer une méthode réglée à des références laissées à
leurs défauts serait trompeur. Dosées par un seul balayage de taux (\cref{tab:refs}),
\code{gold\_soft} fait jeu égal avec \code{trust\_region} et \code{scrub} la dépasse : ses trois
graines à lr \num{e-4} (\num{0.074} ; \num{0.102} ; \num{0.117}) sont toutes au-dessus des trois de
\code{trust\_region} (\num{0.061} ; \num{0.065} ; \num{0.066}). Le mécanisme se lit dans les
colonnes : à divergence d'oubli comparable, \code{scrub} déplace deux à trois fois moins le fold
d'audit (KL \num{0.067}, contre \num{0.117} pour \code{gold\_soft} et \num{0.17} à \num{0.23} pour
\code{trust\_region}). Mais \code{scrub} est bruité (écart-type neuf fois plus grand) et tient sur
une crête : \num{0.037} à lr \num{5e-5}, \num{0.118} à \num{2e-4} (une graine), effondrement à
\num{3e-4}.

\begin{table}[ht]
  \centering\small
  \caption{La V1 face aux références dosées (RTX 5080). Moyenne $\pm$ écart-type quand il y a trois
  graines ; KL audit : $\KL(p_\orig\|p_\sub)$ à $T=2$ sur le fold d'audit.}
  \label{tab:refs}
  \sisetup{table-format=1.3}\setlength{\tabcolsep}{3.6pt}
  \begin{tabular}{@{}llc c SSS S S S[table-format=1.1]@{}}
    \toprule
    Méthode & Réglage & Gr. & {\code{score\_min}} & {$s_{40}$} & {$s_{50}$} & {$s_{60}$} & {Retain}
      & {KL audit} & {Ratio}\\
    \midrule
    trust\_region & $\varepsilon=\num{0.10}$ & 3 & \num{0.061 +- 0.003} & 0.071 & 0.061 & 0.064 & 0.825 & 0.169 & 1.4\\
    trust\_region & $\varepsilon=\num{0.15}$ & 3 & \num{0.064 +- 0.003} & 0.082 & 0.069 & 0.064 & 0.768 & 0.228 & 1.4\\
    trust\_region & $\varepsilon=\num{0.20}$ & 3 & \num{0.058 +- 0.006} & 0.087 & 0.073 & 0.058 & 0.739 & 0.266 & 1.3\\
    \midrule
    scrub & défaut (\num{e-5}) & 3 & \num{0.0025 +- 0.0003} & 0.003 & 0.005 & 0.010 & 0.998 & 0.002 & 4.9\\
    scrub & lr \num{e-4} & 3 & \num{0.098 +- 0.022} & 0.098 & 0.147 & 0.136 & 0.931 & 0.067 & 2.0\\
    scrub & lr \num{2e-4} & 1 & \num{0.118} & 0.165 & 0.157 & 0.118 & 0.663 & 0.423 & 1.2\\
    scrub & lr \num{3e-4} & 1 & 0 & 0.106 & 0.033 & 0 & 0.405 & 0.960 & 1.1\\
    \midrule
    gold\_soft & défaut (3 époques) & 1 & 0 & 0.073 & 0.014 & 0 & 0.401 & 1.022 & 1.1\\
    gold\_soft & 1 époque, \num{e-5} & 3 & \num{0.059 +- 0.012} & 0.059 & 0.069 & 0.060 & 0.882 & 0.117 & 1.6\\
    \bottomrule
  \end{tabular}
\end{table}

\begin{retenir}[Verdict sur la V1 de \code{trust\_region}]
\textbf{Bonne, mais pas la meilleure.} Contre les références à leurs défauts, elle gagne d'un
facteur 25 ; contre des références dosées, \code{scrub} la dépasse d'environ \SI{50}{\percent} en
moyenne et \code{gold\_soft} fait jeu égal (deux références de banc, non conformes au sujet pour
une méthode soumise, \cref{sec:pistes}). La V2 la dépasse ensuite (\cref{sec:campagnev2}). Ses atouts mesurés : un seul réglage utile ($\varepsilon$,
plateau entre \num{0.1} et \num{0.2}), une reproductibilité neuf fois meilleure que celle de
\code{scrub}, une insensibilité au taux, au warmup, à la température et au lot. Son défaut
mesuré : la localité (ratio \num{1.4} ; \num{1.14} sur les seules images exactes,
\cref{sec:diagnostic}). Réglage à garder : $\varepsilon=\num{0.15}$, $\tau=6$,
rayon testé à chaque pas, tout le réseau entraîné, lr \num{3e-5} pour aller plus vite.
\end{retenir}

\subsection{Campagne V2 : une cible propre à chaque image}\label{sec:campagnev2}
La V2 (\cref{sec:trustv2}) a été écrite en parallèle de la campagne V1 et fusionnée ensuite. Sa
campagne (30 runs, fiche \chemin{experiments/unlearning/trust-region-v2-campaign.md}) part de ce que
la V1 a appris : rayon testé à chaque pas, rayon utile petit, taux et lot sans effet. Toutes les
comparaisons sont \emph{appariées} : à graine égale, V1 et V2 voient les mêmes lots et le même fold
d'audit, seule la cible change. Règle posée avant les runs : la V2 améliore la V1 si son meilleur
réglage dépasse, sur trois graines, \num{0.0637} plus deux écarts-types, soit \num{0.069}.

\paragraph{La cible, mesurée avant d'entraîner.} Sur les 2800 images entraînées, avec les vues que
tire la méthode (\chemin{scripts/unlearning\_v2\_targets.py}), on a mesuré pour chaque cible sa
distance à l'original (le plus grand rayon atteignable), son accord avec la pseudo-étiquette, son
hétérogénéité (écart-type sur moyenne de $-\log q_y$) et un \emph{plafond} : le score d'oubli
qu'obtiendrait un élève dont la sortie serait exactement la cible, sous chaque doublure
(\cref{tab:ciblev2}). Le plafond dit combien de la doublure la cible contient, donc la part
potentiellement mécanique d'un gain ; un run s'arrête au rayon, loin de sa cible, et le plafond
renseigne une direction, pas un score attendu. Deux constats ont orienté la campagne : au rayon
utile, $\tau=1$ devient atteignable (\num{0.26}) ; et seules les cibles à $\tau=1$ ont un plafond
positif sous les trois doublures, la V1 n'en ayant aucun.

\begin{table}[ht]
  \centering\small
  \caption{Cibles mesurées sur les 2800 images entraînées (graine 0, 3 vues sauf mention).
  Atteignable : $\KL(p_\orig\|q)$ à $T=1$ ; accord : part des images où la cible garde la
  pseudo-étiquette ; plafond : score d'oubli, sur les images exactes, d'un élève égal à la
  cible.}\label{tab:ciblev2}
  \sisetup{table-format=1.3}\setlength{\tabcolsep}{4.5pt}
  \begin{tabular}{@{}l S S S[table-format=1.2] SSS@{}}
    \toprule
    Cible & {Atteignable} & {Accord} & {Hétérog.} & {Plafond $s_{40}$} & {Plafond $s_{50}$}
      & {Plafond $s_{60}$}\\
    \midrule
    V1, $\tau=6$ & 1.217 & 1.000 & 0.27 & 0 & 0 & 0\\
    V1, $\tau=3$ & 0.120 & 1.000 & 0.94 & 0.030 & 0 & 0\\
    V2, aire $0{,}45$, $\tau=4$ & 1.043 & 0.950 & 0.46 & 0 & 0 & 0\\
    V2, aire $0{,}45$, $\tau=2$ & 0.363 & 0.951 & 1.06 & 0.124 & 0.040 & 0\\
    \textbf{V2, aire $0{,}45$, $\tau=1$} & 0.264 & 0.944 & 1.76 & 0.147 & 0.185 & 0.145\\
    V2, aire $0{,}45$, $\tau=1$, 1 vue & 0.630 & 0.848 & 2.81 & 0.290 & 0.216 & 0\\
    V2, aire $0{,}45$, $\tau=1$, 10 vues & 0.228 & 0.981 & 1.00 & 0.107 & 0.128 & 0.093\\
    V2, aire $0{,}25$, $\tau=1$ & 1.082 & 0.699 & 1.37 & 0.330 & 0 & 0\\
    V2, aire $0{,}65$, $\tau=1$ & 0.109 & 0.982 & 2.19 & 0.059 & 0.090 & 0.108\\
    \bottomrule
  \end{tabular}
\end{table}

\paragraph{À cible distante, un gain faible mais constant.} Avec la cible de la spécification
d'origine ($\tau=4$), la V2 dépasse la V1 aux six rayons testés, de \num{0.002} à \num{0.006}
(\cref{fig:v1v2}) : un signe constant, mais des écarts d'un à deux écarts-types chacun. Le gain
passe par la composante des prédictions, pas par celle des pertes, l'inverse du mécanisme annoncé.

\begin{figure}[ht]
  \centering
  \begin{tikzpicture}
    \begin{axis}[width=0.86\linewidth,height=6.2cm,xmode=log,
        xlabel={rayon $\varepsilon$ (échelle logarithmique)},ylabel={\code{score\_min}},
        xmin=0.008,xmax=1,ymin=0,ymax=0.1,
        yticklabel style={/pgf/number format/.cd,fixed,precision=2,use comma},
        xtick={0.01,0.02,0.05,0.1,0.2,0.5,1},
        xticklabels={{0,01},{0,02},{0,05},{0,1},{0,2},{0,5},{1}},
        legend columns=2,legend style={at={(0.5,1.03)},anchor=south,font=\small,draw=none},
        grid=major,grid style={gray!18},axis line style={gray!60},tick style={gray!60},
        label style={font=\small},tick label style={font=\small}]
      \addplot[bleu,thick,mark=*,mark size=1.6pt] coordinates
        {(0.01,0.0236) (0.02,0.0328) (0.05,0.0511) (0.1,0.0594) (0.15,0.0647) (0.2,0.0595) (0.3,0.0338) (0.5,0.0018) (0.8,0)};
      \addlegendentry{V1, $\tau=6$ (graine 0)}
      \addplot[ambre,thick,mark=square*,mark size=1.4pt] coordinates
        {(0.05,0.0564) (0.1,0.0610) (0.15,0.0678) (0.2,0.0641) (0.3,0.0394) (0.5,0.0023)};
      \addlegendentry{V2, $\tau=4$ (graine 0)}
      \addplot[brique,very thick,mark=triangle*,mark size=2pt] coordinates
        {(0.1,0.0735) (0.15,0.0811) (0.2,0.0782)};
      \addlegendentry{V2, $\tau=1$ (graine 0)}
      \addplot[only marks,brique,mark=diamond*,mark size=2.4pt,
          error bars/.cd,y dir=both,y explicit] coordinates {(0.15,0.0831) +- (0,0.0042)};
      \addlegendentry{V2, $\tau=1$, trois graines}
    \end{axis}
  \end{tikzpicture}
  \caption{V1 et V2 au même rayon, rayon testé à chaque pas (RTX 5080). Le losange et sa barre :
  moyenne et écart-type de la V2 retenue sur trois graines ; la V1 au même point vaut
  \num{0.0637 +- 0.0025}.}
  \label{fig:v1v2}
\end{figure}

\paragraph{La température décide, l'aire et le nombre de vues non.} À $\varepsilon=\num{0.15}$
(\cref{tab:v2reglages}), le score monte quand $\tau$ baisse, de \num{0.066} à $\tau=6$ à
\num{0.081} à $\tau=1$, alors que la V1 est plate en $\tau$. L'aire des vues ne change presque rien :
à $\tau=2$, cinq aires de \num{0.25} à \num{0.65} donnent le même score à \num{0.0007} près ; à
$\tau=1$, l'aire \num{0.25}, qui ne recoupe aucune doublure et dont la cible se trompe sur
\SI{30}{\percent} des images, garde un gain de \num{0.013} sur la V1. Le gain ne suit donc pas le
recouvrement avec les doublures : l'hypothèse d'un gain surtout mécanique par les aires est
réfutée, dans la forme testable ici. Une seule vue dégrade ; au-delà de trois, rien ne change.
Lecture proposée, cohérente avec ces mesures mais non démontrée : le run s'arrêtant au rayon, seule
compte la \emph{direction} du premier mouvement, $p_\theta-q$, c'est-à-dire la façon dont la cible
répartit la masse retirée à la classe en tête ; un $\tau$ bas la concentre sur les classes que les
vues confondent réellement, l'aire et le nombre de vues changent surtout la quantité retirée, que
le rayon normalise. La V1 ne peut pas descendre en $\tau$ : sa cible y est trop près de l'original.

\begin{table}[ht]
  \centering\small
  \caption{V2, réglages balayés à $\varepsilon=\num{0.15}$ (graine 0 ; V1 au même point :
  \num{0.0647}). Valeur $\to$ \code{score\_min} ; en gras, le réglage retenu.}\label{tab:v2reglages}
  \begin{tabularx}{\linewidth}{@{}>{\raggedright}p{3.4cm}X@{}}
    \toprule
    Réglage & Valeurs $\to$ \code{score\_min}\\
    \midrule
    $\tau$ (aire $0{,}45$, 3 vues) & 6 $\to$ \num{0.066} ; 4 $\to$ \num{0.068} ; 3 $\to$ \num{0.070} ;
      2 $\to$ \num{0.072} ; \textbf{1} $\to$ \num{0.081} ; \num{0.75} $\to$ \num{0.084} ;
      \num{0.5} $\to$ \num{0.085} ; \num{0.25} $\to$ \num{0.077} (rayon jamais atteint)\\
    Aire, à $\tau=1$ & \num{0.25} $\to$ \num{0.078} ; \num{0.35} $\to$ \num{0.081} ;
      \textbf{\num{0.45}} $\to$ \num{0.081}\\
    Aire, à $\tau=2$ & \num{0.25} $\to$ \num{0.072} ; \num{0.35} $\to$ \num{0.072} ; \num{0.45} $\to$
      \num{0.072} ; \num{0.55} $\to$ \num{0.072} ; \num{0.65} $\to$ \num{0.072}\\
    Vues, à $\tau=1$ & 1 $\to$ \num{0.077} ; \textbf{3} $\to$ \num{0.081} ; 6 $\to$ \num{0.080} ;
      10 $\to$ \num{0.080}\\
    Rayon, à $\tau=1$ & \num{0.10} $\to$ \num{0.074} ; \textbf{\num{0.15}} $\to$ \num{0.081} ;
      \num{0.20} $\to$ \num{0.078}\\
    \bottomrule
  \end{tabularx}
\end{table}

\paragraph{Sous $\tau=1$, un mode de défaillance.} Les points $\tau=\num{0.5}$ et $\num{0.75}$
semblaient prolonger la tendance. Sur trois graines, $\tau=\num{0.5}$ a la même moyenne que
$\tau=1$ (\num{0.083}) mais un écart-type deux fois et demie plus grand : trois runs à $\tau\le\num{0.5}$
n'atteignent jamais le rayon. La divergence à l'original y plafonne vers \num{0.085}, l'entraînement
va au bout des 880 pas, et le modèle change de régime (composante des pertes effondrée, vers
\num{0.2}, composante des prédictions vers \num{0.6}). La distance de la cible n'est pas en cause
(\num{0.27} contre \num{0.26} à $\tau=1$) : c'est sa forme, très piquée sur quelques classes. Les
sept runs à $\tau=1$ atteignent tous le rayon, en 89 à 104 pas.

\paragraph{Trois graines, appariées.} Au réglage retenu ($\tau=1$, aire \num{0.45}, 3 vues,
$\varepsilon=\num{0.15}$), la V2 dépasse la V1 à chacune des trois graines, de \num{0.016},
\num{0.027} et \num{0.015} (\cref{tab:v1v2}). Les deux modèles ont bougé exactement autant (KL sur
les images oubliées \num{0.261} contre \num{0.261}), abîmé autant le fold d'audit (KL \num{0.225}
contre \num{0.228}, retain \num{0.78} contre \num{0.77}) et transféré autant aux copies altérées ;
la différence tient entière dans la composante des prédictions, plus haute sous les trois doublures
(par exemple \num{0.245} contre \num{0.180} sous $s_{60}$). Même quantité d'oubli, meilleure
direction.

\begin{table}[ht]
  \centering\small
  \caption{V1 et V2 retenues, trois graines appariées ($\varepsilon=\num{0.15}$, mêmes lots, même
  fold d'audit), et la meilleure référence. $S_{\text{pred}}$ et $S_{\text{ms}}$ sous $s_{60}$.}
  \label{tab:v1v2}
  \sisetup{table-format=1.3}\setlength{\tabcolsep}{3.4pt}
  \begin{tabular}{@{}l c c SSS SS S S@{}}
    \toprule
    Méthode & Graines 0 ; 1 ; 2 & {\code{score\_min}} & {$s_{40}$} & {$s_{50}$} & {$s_{60}$}
      & {$S_{\text{ms}}$} & {$S_{\text{pred}}$} & {Retain} & {KL audit}\\
    \midrule
    V1, $\tau=6$ & \num{0.065} ; \num{0.061} ; \num{0.066} & \num{0.0637 +- 0.0025} & 0.082 & 0.069
      & 0.064 & 0.465 & 0.180 & 0.768 & 0.228\\
    V2, $\tau=1$ & \num{0.081} ; \num{0.088} ; \num{0.080} & \num{0.0831 +- 0.0042} & 0.096 & 0.084
      & 0.086 & 0.452 & 0.245 & 0.778 & 0.225\\
    \midrule
    scrub, lr \num{e-4} & \num{0.102} ; \num{0.074} ; \num{0.117} & \num{0.098 +- 0.022} & 0.098
      & 0.147 & 0.136 & {--} & {--} & 0.931 & 0.067\\
    \bottomrule
  \end{tabular}
\end{table}

\paragraph{Contrôle sous une seconde famille de doublures.} Le banc \code{bench-v4}, fusionné après
la campagne, ajoute aux recadrages le masquage de patchs, une famille de doublures étrangère à la
cible de la V2. Recalculés sans réentraînement (\chemin{experiments/unlearning/measurements/bench\_v4\_rebench.json}),
les fichiers retenus donnent le \cref{tab:v4} : les scores sous recadrage sont identiques à
\code{bench-v3}, et la V2 bat la V1 sous les deux familles, à chaque graine. Mais sous la seule
famille de masquage, son gain n'est plus que de \SI{12}{\percent}, contre \SI{30}{\percent} sous
recadrage : environ la moitié du gain mesuré sous recadrage est probablement liée à la famille de
vues que partagent la cible et les doublures, la part que l'étape sur les aires ne pouvait pas
tester. \code{scrub}, dont la cible n'est pas faite de vues, ne bouge pas d'une famille à l'autre.
Sous la pire des six doublures, la V2 vaut \num{0.0786 +- 0.0030}, soit \SI{23}{\percent} de plus que
la V1, au-dessus du seuil de décision.

\begin{table}[ht]
  \centering\small
  \caption{Contrôle en \code{bench-v4} (trois graines, fichiers recalculés). Pire recadrage : pire de
  $s_{40}$, $s_{50}$, $s_{60}$ ; pire masquage : pire de $m_{30}$, $m_{20}$, $m_{15}$ ;
  \code{score\_min} : pire des six.}\label{tab:v4}
  \setlength{\tabcolsep}{5pt}
  \begin{tabular}{@{}lccc@{}}
    \toprule
    & Pire recadrage & Pire masquage & \code{score\_min} \code{bench-v4}\\
    \midrule
    V1, $\varepsilon=\num{0.15}$ & \num{0.0637 +- 0.0025} & \num{0.0702 +- 0.0013} & \num{0.0637 +- 0.0025}\\
    V2, $\tau=1$ & \num{0.0831 +- 0.0042} & \num{0.0786 +- 0.0030} & \num{0.0786 +- 0.0030}\\
    Écart V2 $-$ V1 & $+\SI{30}{\percent}$ & $+\SI{12}{\percent}$ & $+\SI{23}{\percent}$\\
    \midrule
    \code{scrub}, lr \num{e-4} (référence) & \num{0.098 +- 0.022} & \num{0.101 +- 0.023} & \num{0.098 +- 0.022}\\
    \bottomrule
  \end{tabular}
\end{table}

\begin{retenir}[Verdict sur la V2]
\textbf{La V2 améliore la V1 sur le banc}, graine par graine, à quantité d'oubli et à retain
égaux : de \SI{30}{\percent} sous recadrage (\num{0.0831 +- 0.0042} contre \num{0.0637 +- 0.0025}),
de \SI{23}{\percent} sous la pire des six doublures de \code{bench-v4} (\num{0.0786 +- 0.0030}), de
\SI{12}{\percent} sous la seule famille étrangère à sa cible ; elle dépasse dans tous les cas le
seuil fixé avant les runs (\num{0.069}). Elle rejoint \code{scrub} dosé dans la limite du bruit de
celui-ci, en étant cinq fois plus stable ; \code{scrub} n'étant qu'une référence de banc (non
conforme au sujet), la V2 est la meilleure méthode conforme mesurée. Réglage retenu : \code{trust\_target=views},
$\tau=1$, aire \num{0.45}, 3 vues, $\varepsilon=\num{0.15}$, rayon testé à chaque pas. À éviter : la
spécification d'origine telle quelle, $\tau\ge4$ (gain faible), $\tau\le\num{0.5}$ (rayon parfois
jamais atteint), une seule vue. Réserve : une partie du gain sous recadrage tient à la famille de
vues commune à la cible et aux doublures ; on retient le chiffre \code{bench-v4} et seul le
classement public dira ce qui se transfère au score réel. Défaut commun à la V1 : la localité
(ratio \num{1.36}).
\end{retenir}

\subsection{Campagne V3 : la cible calibrée}\label{sec:campagnev3}
La V3 (\cref{sec:trustv3}) a été testée en \code{bench-v4}, sous les six doublures, par 23 runs
(fiche \chemin{experiments/unlearning/trust-region-v3-campaign.md}), avec pour contrat les critères
de réfutation posés par sa fiche d'origine avant tout run. Les références V1, V2 et \code{scrub}
ont été relancées au lab en \code{bench-v4} (\code{ref4-*}) : elles redonnent exactement les chiffres
recalculés du \cref{tab:v4}. Seuil d'amélioration sur la V2, fixé avant les runs : \num{0.0786} plus
deux écarts-types, soit \num{0.085}, et un gain sous le masquage au-delà de deux écarts-types.

\paragraph{La cible, mesurée avant d'entraîner.} 560 cibles ont été mesurées sur les 2800 images
entraînées avec la fonction de la méthode (\chemin{scripts/unlearning\_v3\_targets.py}). Deux
constats. La cible de la V2 retenue a un plafond nul sous le masquage $m_{15}$, ce qui éclaire le
contrôle du \cref{tab:v4} ; celle de la V3 de la spécification est positive sous les six doublures
(\num{0.041} au pire). Et la V3 diffère de la V2 sur un point de structure : son rayon est une
fraction $\rho$ de la distance de sa cible ($\rho=\num{0.9}$), si bien que l'élève finit près de sa
cible.

\paragraph{La spécification d'origine fait jeu égal avec la V2.} À $\lambda=\num{0.25}$, la V3 de la
spécification vaut \num{0.0761} (V2 à la même graine : \num{0.0757}) ; à $\lambda=\num{0.5}$, sa
cible pousse le rayon à \num{0.25} et le score tombe à \num{0.0566}, avec une dérive d'audit au-delà
du seuil de casse (KL \num{0.343}). Les copies douces ne changent ni le score (\num{0.0762}) ni le
transfert aux altérations jamais vues (jitter et JPEG) : réfutées. Le mécanisme attendu par
$S_{\text{ms}}$ ne s'observe pas : la dispersion des pertes de l'élève vaut environ \num{3.3} fois
leur moyenne pour les trois versions, quelle que soit la cible ; c'est le réseau qui la fixe.

\paragraph{Tous les réglages agissent par le rayon.} Les balayages (\cref{tab:v3reglages})
montrent un motif unique : $\lambda$, la couverture $b$ et l'aire des vues déplacent la distance
de la cible, donc le rayon effectif, et le score suit le rayon avec un optimum vers
\numrange{0.15}{0.18}, comme pour la V1 et la V2. La couverture nuit de façon monotone ($b=0$ est le
meilleur réglage) ; à rayon fixe \num{0.15}, la V3 sans couverture fait jeu égal avec la V2
(\num{0.0762} contre \num{0.0757}). Les cibles à une ou deux vues, les plus prometteuses avant
entraînement (plafond \num{0.12} sous les six doublures), retrouvent le mode de défaillance de la V2
à $\tau\le\num{0.5}$ : le rayon n'est jamais atteint. Un plafond mesuré avant entraînement ne vaut
que si le rayon est atteint.

\begin{table}[ht]
  \centering\small
  \caption{V3, réglages balayés en \code{bench-v4} (graine 0). Base : $\lambda=\num{0.25}$,
  $b=\num{0.25}$, aire \num{0.5}, 4 vues, $\rho=\num{0.9}$, sans copies (\num{0.0761}, rayon
  \num{0.159}). Valeur $\to$ \code{score\_min} (rayon effectif). V2 à la même graine :
  \num{0.0757}.}\label{tab:v3reglages}
  \begin{tabularx}{\linewidth}{@{}>{\raggedright}p{3.2cm}X@{}}
    \toprule
    Réglage & Valeurs $\to$ \code{score\_min} (rayon)\\
    \midrule
    $\lambda$ & 0 $\to$ \num{0.060} (\num{0.07}) ; \num{0.1} $\to$ \num{0.068} (\num{0.10}) ;
      \textbf{\num{0.25}} $\to$ \num{0.076} (\num{0.16}) ; \num{0.5} $\to$ \num{0.057} (\num{0.25})\\
    Couverture $b$ & \num{0.25} $\to$ \num{0.076} ; \num{0.1} $\to$ \num{0.077} ;
      \textbf{0} $\to$ \num{0.078} (\num{0.17})\\
    Rayon & $\rho$ \num{0.6} $\to$ \num{0.072} (\num{0.11}) ; \num{0.75} $\to$ \num{0.074} (\num{0.13}) ;
      \num{0.9} $\to$ \num{0.076} (\num{0.16}) ; à $b=0$ : \num{0.95} $\to$ \num{0.078} (\num{0.18}),
      $\varepsilon$ fixe \num{0.15} $\to$ \num{0.076}\\
    Vues et aire & 8 vues $\to$ \num{0.074} ; à $b=0$ : aire \num{0.4} $\to$ \num{0.042} (\num{0.35},
      casse), \num{0.6} $\to$ \num{0.072} (\num{0.11}) ; 1 ou 2 vues $\to$ \numrange{0.038}{0.054}
      (rayon jamais atteint)\\
    Copies douces & \num{0.5} $\to$ \num{0.076} (sans effet)\\
    \bottomrule
  \end{tabularx}
\end{table}

\paragraph{Trois graines : jeu égal.} Au meilleur réglage ($b=0$, $\lambda=\num{0.25}$,
$\rho=\num{0.9}$, aire \num{0.5}, 4 vues), la V3 vaut \num{0.0798 +- 0.0012} contre
\num{0.0786 +- 0.0030} pour la V2, appariées graine par graine (\num{+0.0027}, \num{-0.0013},
\num{+0.0021}) ; sous le masquage, \num{0.0798} contre \num{0.0786}. Les deux restent sous les seuils
fixés avant les runs : pas d'amélioration. La V3 est deux fois et demie plus stable ; elle abîme un
peu plus le fold d'audit (retain \num{0.74} contre \num{0.78}). Son avance sous recadrage
(\num{0.0876} contre \num{0.0831}) est sur la famille que sa cible partage avec les doublures.

\begin{retenir}[Verdict sur la V3]
\textbf{Jeu égal avec la V2, pas une amélioration.} Meilleur réglage : \code{trust\_target=calibrated},
$b=0$, $\lambda=\num{0.25}$, $\rho=\num{0.9}$, aire \num{0.5}, 4 vues, sans copies douces :
\num{0.0798 +- 0.0012} en \code{bench-v4}. Réfutés : le gain sur V1 et V2, le mécanisme par
$S_{\text{ms}}$, les copies douces. Leçon de méthode : avec un rayon tiré de la cible, tout réglage
de la cible agit d'abord par le rayon ; comparer des cibles demande un rayon fixe. V2 et V3
convergent vers \num{0.08} : la cible n'est plus le levier, la localité (ratio
\numrange{1.25}{1.37}) l'est.
\end{retenir}

\subsection{Diagnostic : un oubli générique}\label{sec:diagnostic}
Relus le 22 novembre, sans nouveau run, les registres des trois campagnes disent \emph{comment}
la région de confiance oublie (\cref{tab:diagnostic} ; tous les chiffres viennent de
\chemin{experiments/registry.jsonl}).
\begin{itemize}[leftmargin=1.2em]
  \item \textbf{Le fold d'audit bouge presque autant que les images entraînées.} Sur les 31 runs
    de la V1 testés à chaque pas, sa divergence vaut de \num{0.80} à \num{0.96} fois celle des images
    exactes entraînées (\num{0.87} en moyenne), quels que soient le rayon, le taux, le lot, le
    warmup, $\tau$ ou la graine ; ses pertes ont la même dispersion (écart-type sur moyenne de
    \num{3.63} contre \num{3.63}, graine 0). Environ \SI{88}{\percent} du mouvement est donc commun à
    toutes les images mémorisées. Le \enquote{ratio 1,4} cité plus haut moyenne les images exactes et
    leurs copies : le jitter, qui bouge deux fois plus que l'image exacte, le gonfle ; sur les images
    exactes seules, il vaut \num{1.14}.
  \item \textbf{Changer la cible n'y a rien changé.} La V2 et la V3 retenues battent la V1 par la
    direction de leur cible (\cref{sec:campagnev2,sec:campagnev3}), mais leur fold d'audit bouge
    autant : \num{0.86} fois l'image exacte pour la V2 (\code{ref4-v2}, trois graines),
    \numrange{0.86}{0.91} pour la V3 retenue (trois graines), contre \numrange{0.87}{0.88} pour la V1
    (\code{ref4-v1}). Côté V1, quand $\tau$ passe de 4 à 12, la dispersion relative des pertes de
    l'élève reste entre \num{3.61} et \num{3.68} (la cible, à $\tau=6$, est à \num{0.28}) : la
    dispersion est fixée par le réseau, ce que la campagne V3 a retrouvé (\cref{sec:campagnev3}).
    Tous les arrêts de la V1 tombent entre les pas 72 et 92, autour de la fin du warmup (88 pas) :
    au taux maximal, le modèle est figé en plein vol.
  \item \textbf{\code{scrub} oublie autrement.} À la même perte moyenne (\num{0.113}), il bat la V1 sur
    les trois composantes, et son fold d'audit ne bouge qu'à \numrange{0.51}{0.60} fois l'image exacte
    (\code{ref4-scrub}, trois graines) : il concentre le mouvement sur moins d'images, avec une
    dispersion relative des pertes deux fois plus forte. Réserves : ses max-steps s'entraînent sur
    les quatre familles d'altérations du banc (le retournement et le JPEG y sont déterministes), donc
    une partie de la moitié \enquote{copies} de $S_{\text{ms}}$ et $S_{\text{pred}}$ lui est déjà
    connue ; seule la comparaison sur les images exactes est propre, et le banc la donne désormais
    à part (\code{bench.exact.*}, hors score). Et son meilleur run (lr \num{2e-4}, \num{0.118}) est
    aussi générique que la V1 (audit/exact \num{0.88}) : son avance vient du terme d'oubli
    ($S_{\text{ms}}=\num{0.745}$), pas de la localité.
\end{itemize}

\begin{table}[ht]
  \centering\small
  \caption{Mécanisme de l'oubli à pleine échelle (\code{bench-v3}, RTX 5080, registre). Audit/exact :
  $\KL(p_\orig\|p_\sub)$ à $T=2$ sur le fold d'audit, rapportée à celle des images exactes
  entraînées ; CV : écart-type sur moyenne des pertes de l'élève sur ces images ; $S_{\text{ms}}$ et
  $S_{\text{pred}}$ sous $s_{40}$, $s_{50}$, $s_{60}$.}\label{tab:diagnostic}
  \sisetup{table-format=1.3}\setlength{\tabcolsep}{3.4pt}
  \begin{tabular}{@{}l S[table-format=1.3] S[table-format=1.2] S[table-format=1.1] c c S S[table-format=1.3]@{}}
    \toprule
    Run & {Perte} & {Audit/exact} & {CV} & {$S_{\text{ms}}$} & {$S_{\text{pred}}$} & {Retain} & {Score}\\
    \midrule
    trust\_region V1, $\varepsilon=\num{0.10}$ & 0.113 & 0.87 & 4.3 & .27/.27/.35 & .31/.27/.22 & 0.830 & 0.059\\
    trust\_region V1, $\varepsilon=\num{0.15}$ & 0.172 & 0.88 & 3.6 & .34/.35/.46 & .32/.26/.18 & 0.771 & 0.065\\
    scrub, lr \num{e-4} & 0.113 & 0.55 & 7.6 & .30/.48/.49 & .37/.35/.32 & 0.928 & 0.102\\
    gold\_soft, 1 époque & 0.096 & 0.81 & 4.6 & .24/.27/.31 & .35/.31/.26 & 0.854 & 0.070\\
    \bottomrule
  \end{tabular}
\end{table}

\begin{retenir}
La région de confiance n'oublie pas \emph{ces} images : elle rend le modèle moins sûr de toutes les
images qu'il a mémorisées. Changer sa cible a amélioré la direction (V2, V3) sans toucher ce
mécanisme ; un terme qui fait payer le mouvement des images mémorisées non visées le peut. C'est
l'objet de la V4 (\cref{sec:trustv4}).
\end{retenir}

\subsection{Campagne V4 : l'ancre, la contrainte et SCRUB conforme}\label{sec:campagnev4}
La V4 (\cref{sec:trustv4}) a été mesurée en \code{bench-v4} par 26 runs, le 22 novembre, sur la RTX
5080 (fiche \chemin{experiments/unlearning/trust-region-v4-campaign.md}, une fiche par
composante : \chemin{scrub-forget-views.md}, \chemin{trust-region-dual.md},
\chemin{trust-region-anchored.md}). Chaque spécification est jugée contre les hypothèses et les
seuils que sa fiche avait posés avant tout run. Les suivis obéissent à une règle écrite elle aussi
avant les runs : confirmer un candidat sur ses graines manquantes, appliquer un remède déjà prévu
par la spécification, ou suivre une dose dont la mesure donne la direction ; chacun est posé dans
sa fiche avant d'être lancé. Références appariées à graine égale : la V3 retenue
(\num{0.0798 +- 0.0012}) et \code{scrub} aux images proxy (\code{ref4-scrub}, non conforme).

\paragraph{Non-régression.} Avant tout run de la V4, \code{ref4-v3} et \code{ref4-scrub} ont été
rejoués sur le code fusionné : leurs 159 métriques communes sont identiques au bit près à celles
des runs d'origine, sur les trois graines (seule la durée diffère). Les options de la V4 sont
neutres à leur défaut. Un premier passage de \code{scrub-forget-views} et de ce rejeu a tourné avec
un fichier de test modifié (drapeau \code{dirty}) ; relancé sur un arbre propre, il redonne les
mêmes chiffres au bit près. Seuls les runs propres sont cités ; les sept autres restent au registre.

\paragraph{SCRUB conforme : l'ancre déplace la dose.} Aux doses de la fiche, celles qui encadrent le
meilleur SCRUB proxy (\num{e-4} et \num{2e-4}), la version conforme est réfutée
(\cref{tab:v4scrub}) : \num{0.017} et \num{0.061} en moyenne, sous la V3. Son mécanisme tient
(audit/exact \numrange{0.22}{0.51}), mais au prix de l'oubli. À \num{e-4}, la KL des images exactes
ne vaut que \numrange{0.27}{0.31} fois celle du proxy : l'ancre, faite de vues des mêmes images,
défait ce que la montée a fait. Elle tient pourtant bien mieux que l'ancre externe : à \num{2e-4},
retain \num{0.96} et KL d'audit \num{0.04}, quand le proxy au même taux en est à \num{0.66} et
\num{0.42}, puis s'effondre à \num{3e-4}. D'où un suivi aux doses supérieures. À \num{3e-4}, la
version conforme vaut \num{0.1085 +- 0.0068} sur trois graines, avec un gain apparié sur la V3, sous
le masquage, de \num{+0.037}, \num{+0.021} et \num{+0.028} (près de $10\sigma$ en moyenne) et un
retain de \num{0.86}. Elle fait jeu égal avec le SCRUB proxy, aux mêmes graines (\num{+0.013},
\num{+0.028}, \num{-0.008}, dans le bruit de ce dernier). À \num{4e-4}, elle monte à
\num{0.121 +- 0.007}, mais avec un retain de \num{0.73}, sous le plancher de \num{0.80} posé avant :
écartée.

\begin{table}[ht]
  \centering\small
  \caption{SCRUB conforme (\code{retain\_source=forget\_views}) selon la dose, \code{bench-v4}, moyennes
  sur $n$ graines ($\pm$ : écart-type de \code{score\_min}). Recadr., Masq. : pire doublure de chaque
  famille ; Aud./ex. : KL du fold d'audit sur celle des images exactes entraînées (à $T=2$). Les deux
  premières doses sont celles de la fiche ; les deux suivantes, celles du suivi.}\label{tab:v4scrub}
  \sisetup{table-format=1.3}\setlength{\tabcolsep}{3pt}
  \begin{tabular}{@{}l c S[table-format=1.4] S[table-format=1.4] S S S S S S[table-format=1.2]@{}}
    \toprule
    Ancre, lr & $n$ & {\code{score\_min}} & {$\pm$} & {Recadr.} & {Masq.} & {Retain} & {KL aud.}
      & {KL exacte} & {Aud./ex.}\\
    \midrule
    conforme, \num{e-4} & 2 & 0.0174 & 0.0052 & 0.018 & 0.018 & 0.993 & 0.007 & 0.031 & 0.23\\
    conforme, \num{2e-4} & 2 & 0.0607 & 0.0129 & 0.064 & 0.061 & 0.959 & 0.040 & 0.082 & 0.49\\
    \textbf{conforme, \num{3e-4}} & 3 & \bfseries 0.1085 & 0.0068 & 0.116 & 0.109 & 0.859 & 0.142
      & 0.167 & 0.85\\
    conforme, \num{4e-4} & 3 & 0.1208 & 0.0069 & 0.132 & 0.121 & 0.733 & 0.284 & 0.271 & 1.05\\
    \midrule
    proxy, \num{e-4} (non conforme) & 3 & 0.0978 & 0.0220 & 0.098 & 0.101 & 0.931 & 0.067 & 0.120 & 0.55\\
    V3 retenue & 3 & 0.0798 & 0.0012 & 0.088 & 0.080 & 0.744 & 0.261 & 0.294 & 0.89\\
    \bottomrule
  \end{tabular}
\end{table}

\paragraph{Pourquoi elle gagne.} Pas par la localité : à la dose retenue, le fold d'audit bouge
\num{0.85} fois autant que les images exactes (\numrange{0.80}{0.93}), comme pour la V3
(\num{0.89}) ; la version conforme n'est locale qu'aux doses où elle n'oublie pas. Elle gagne par
l'efficacité de son oubli. À KL exacte presque deux fois moindre (\num{0.17} contre \num{0.29}), sa
composante des prédictions sous sa pire doublure ($m_{30}$) vaut \numrange{0.35}{0.37}, contre
\numrange{0.23}{0.27} pour la V3 sous la sienne ($m_{15}$). Sa composante des pertes est un peu plus
basse (\numrange{0.33}{0.38} contre \numrange{0.41}{0.44}) ; le produit, l'oubli, est plus haut
(\numrange{0.12}{0.13} contre \numrange{0.10}{0.11}). Son retain, lui, est plus haut de \num{0.12}.
Réserve de circularité, contrôlée : les max-steps de \code{scrub} s'entraînent sur les familles
d'altérations du banc. Mais le gain est le même sous le masquage et sous le recadrage, et sur les
seules images exactes (\code{bench.exact.*}, hors score), la composante des prédictions
(\numrange{0.31}{0.41}) est celle du proxy (\numrange{0.35}{0.42}).

\paragraph{Contrainte active : elle ne lie pas.} Les quatre runs ($\eta\in\{\num{0.25};\num{0.5}\}$,
graines 0 et 1) valent tous \num{0.042}, contre \num{0.078} et \num{0.080} pour la V3 appariée
(\cref{tab:v4tr}). L'historique dit pourquoi. La KL de la sonde à $T=1$ franchit le rayon vers le pas
100, au taux maximal ; $\mu$ monte pendant la deuxième époque, ramène la KL sous le rayon, puis
retombe à 0 et y reste de la troisième à la dixième époque. Les deux valeurs de $\eta$ donnent le même
modèle à \num{e-4} près. Libre, l'objectif de la V3 poursuivi sur tout son calendrier se pose de
lui-même à une KL de \num{0.085} à $T=1$, la moitié du rayon. Il ajuste sa cible bien plus (perte
\num{0.63} contre \numrange{2.1}{2.2} à l'arrêt de la V3), tandis qu'à $T=2$ la KL des images exactes
monte (\num{0.34} contre \numrange{0.28}{0.30}). Au sens de la fiche, les runs sont non
interprétables (KL finale à \numrange{0.53}{0.54} fois le rayon), mais pas pour la raison prévue,
un $\eta$ mal réglé : la contrainte est une borne haute ($\mu\ge0$ ne peut que ramener vers
l'original), et elle ne lie pas. Aucun $\eta$ ne la ferait lier. Ce qui est mesuré réfute en
revanche la prémisse du mode. L'élève apprend bien la forme de sa cible : la composante des
prédictions double (\numrange{0.32}{0.69} contre \numrange{0.23}{0.38}), et le pronostic du
\cref{sec:trustv4} se vérifie (converger vers la cible égalise les pertes, composante des pertes
\numrange{0.12}{0.19} contre \numrange{0.34}{0.47}). Mais le score est divisé par deux, et la KL
d'audit monte de \SIrange{25}{39}{\percent}. Abandonnée.

\paragraph{Ancre croisée : elle freine l'oubli, elle ne le localise pas.} À lr \num{e-5}, aucune
phase n'atteint son rayon. La KL de la sonde plafonne à \numrange{0.02}{0.07}, pour des rayons de
\numrange{0.14}{0.18} : l'ancre freine la montée pendant le court warmup de chaque phase (22 pas),
puis on retrouve le régime de la contrainte active, calendrier complet et pertes égalisées. Le
fold d'audit bouge \numrange{0.87}{0.98} fois autant que les images exactes, et aucun bras ne dépasse
\num{0.0844} : deux critères de réfutation sont atteints. Ces seuils supposaient des phases arrêtées
à leur rayon, et la spécification prévoyait, avant tout run, le remède à ce cas : lr \num{3e-5}. Ce
suivi, avec un témoin sans ancre au même taux, rapproche les phases de leur rayon (\numrange{0.82}{0.89}
fois à $w=1$) sans le leur faire atteindre, sauf la première phase à $K=2$. L'ancre y baisse audit/exact de \num{0.86} à
\numrange{0.67}{0.70} : un cinquième du mouvement commun, pas la moitié visée. Son gain ne vaut que
sous le recadrage (\num{+0.013} à \num{+0.019} sur le témoin) et se paie sous le masquage
(\num{-0.017} à \num{-0.029}). Tripler son poids baisse d'autant l'oubli des images poussées et
celui du fold d'audit : elle freine l'oubli partout sans le localiser. L'auto-audit
(\code{pending\_kl}) classe les réglages exactement comme le fold d'audit, mais sous-estime la
dérive de celui-ci d'un facteur \numrange{1.6}{13} : il mesure des images que l'ancre tient, alors
que le fold d'audit n'est tenu par rien. Abandonnée.

\begin{table}[ht]
  \centering\small
  \caption{Options de \code{trust\_region} de la V4, sur la V3 retenue, \code{bench-v4}, graine 0 sauf
  mention. Recadr., Masq. : pire doublure de chaque famille ; KL ex. : KL des images exactes,
  rapportée à celle de la V3 de même graine ; Aud./ex. : comme au \cref{tab:v4scrub} ; Phases :
  phases arrivées à leur rayon.}\label{tab:v4tr}
  \sisetup{table-format=1.3}\setlength{\tabcolsep}{3pt}
  \begin{tabular}{@{}l c S S S S[table-format=1.2] S[table-format=1.2] c@{}}
    \toprule
    Bras & lr & {\code{score\_min}} & {Recadr.} & {Masq.} & {KL ex.} & {Aud./ex.} & Phases\\
    \midrule
    V3 retenue (3 graines) & \num{e-5} & 0.080 & 0.088 & 0.080 & 1.00 & 0.89 & arrêt\\
    contrainte, $\eta=\num{0.25}$ (graines 0-1) & \num{e-5} & 0.042 & 0.060 & 0.042 & 1.18 & 0.99 & --\\
    \midrule
    ancre, $K=4$, $w=1$, exacte & \num{e-5} & 0.047 & 0.059 & 0.047 & 0.95 & 0.98 & 0/4\\
    ancre, $K=2$ & \num{e-5} & 0.045 & 0.057 & 0.045 & 0.84 & 0.96 & 0/2\\
    ancre sur vues & \num{e-5} & 0.055 & 0.072 & 0.055 & 1.29 & 0.87 & 0/4\\
    ancre, $w=3$ & \num{e-5} & 0.035 & 0.044 & 0.035 & 0.58 & 0.97 & 0/4\\
    \midrule
    ancre, $K=4$, $w=1$, exacte & \num{3e-5} & 0.058 & 0.107 & 0.058 & 1.21 & 0.70 & 0/4\\
    ancre, $K=2$ & \num{3e-5} & 0.060 & 0.110 & 0.060 & 1.18 & 0.67 & 1/2\\
    ancre, $w=3$ & \num{3e-5} & 0.070 & 0.103 & 0.070 & 1.03 & 0.69 & 0/4\\
    témoin sans ancre (V3) & \num{3e-5} & 0.088 & 0.090 & 0.088 & 1.18 & 0.86 & arrêt\\
    \bottomrule
  \end{tabular}
\end{table}

\paragraph{Une leçon commune.} Les deux échecs ont la même signature. L'objectif des V1 à V3 ne tient
son score qu'arrêté en plein vol, au premier franchissement du rayon. Dans les deux façons de ne pas
s'arrêter que nous avons mesurées, contrainte active et phases qui n'atteignent pas leur rayon, il
se pose sous le rayon à $T=1$, s'éloigne de l'original à $T=2$, égalise les pertes et perd la moitié
du score. Les arrêts en fin de warmup du \cref{sec:diagnostic} ne sont donc pas un défaut à corriger
en prolongeant l'entraînement.

\paragraph{La V3 au taux \num{3e-5}.} Le témoin du suivi de l'ancre (la V3 retenue, seul le taux
change) vaut \num{0.0876} à la graine 0, \num{+0.0092} sur la même V3 à \num{e-5}. Ses graines 1 et
2, posées avant, redonnent \num{0.0837} et \num{0.0821} : \num{0.0845 +- 0.0028}, et un gain apparié
moyen de \num{+0.0047}, sous $2\sigma$. Réfutée : comme la V1, la V3 est insensible au taux dans le
bruit. À \num{e-4}, les contrôles cassent (KL d'audit \num{0.38}, au-delà de la porte, ou rayon non
atteint).

\begin{retenir}[Verdict sur la V4]
\textbf{SCRUB conforme confirmé, ancre et contrainte réfutées.} Meilleure méthode conforme mesurée :
\code{scrub}, \code{retain\_source=forget\_views}, lr \num{3e-4} : \num{0.1085 +- 0.0068} en
\code{bench-v4}, \SI{36}{\percent} au-dessus de la V3 sous les deux familles, apparié sur trois
graines, retain \num{0.86}. Sa dose n'est pas celle du SCRUB proxy : aux doses de la fiche, l'ancre
défait l'oubli. L'ancre croisée freine l'oubli partout sans le localiser ; la contrainte active ne
lie pas. La localité, levier visé par la V4, n'a été obtenue par aucune composante à une dose qui
score. Aucun de ces fichiers n'est soumissible : ils n'ont pas oublié le fold d'audit.
\end{retenir}

\subsection{Campagne V5 : sélection réfutée, localité sans score}\label{sec:campagnev5}
La V5 (\cref{sec:v5}) a été jugée le 22 novembre sur la RTX 5080 (fiche
\chemin{experiments/unlearning/trust-region-v5-campaign.md} ; fiches \chemin{band-ascent-step0.md} et
\chemin{band-ascent.md}). Avant tout run, la fiche de l'étape 0 posait deux critères de réfutation,
chacun suffisant : une précision de la sélection sous \num{0.35}, ou une part de la perte des
doublures sous \num{0.5}, sur les trois masquages. Elle fixait aussi la règle de ne pas lancer
\code{band-ascent} sous cette forme si l'un des deux tombait. Les masquages font foi : la sélection
est tirée de recadrages, comme les doublures $s_{40}$ à $s_{60}$. L'étape 0 est un seul run, sans
rétropropagation (\SI{51}{\second}, commit propre).

\paragraph{La sélection vaut le hasard sous le masquage.} Sur les 2800 images entraînées
(\cref{tab:v5step0}), les \SI{25}{\percent} d'images où l'original doute le plus sur quatre vues fortes
ne désignent pas mieux les erreurs des masquages qu'un tirage au hasard. Leur précision
(\numrange{0.087}{0.183}) est égale au taux de base (\numrange{0.92}{1.01} fois). Elles portent
\numrange{0.23}{0.24} de la perte des doublures, la part d'un tirage de \SI{25}{\percent}, quand le
meilleur choix de même taille en porterait \numrange{0.92}{0.99}. Sous le recadrage, famille commune
aux vues et aux doublures, la sélection fait \numrange{2.1}{2.6} fois le hasard, mais ne porte encore
que \numrange{0.41}{0.47} de la perte (oracle \numrange{0.74}{0.92}). Les deux critères sont atteints
sous les trois masquages : \code{band-ascent} n'a pas été lancé. Une réserve sur le premier : à
$\phi=\num{0.25}$, la précision ne peut dépasser le taux de base divisé par $\phi$, soit \num{0.344}
sous $m_{15}$, sous le seuil. Le critère portant sur les trois masquages, et la précision restant au
taux de base sous $m_{30}$ et $m_{20}$ (plafonds \num{0.76} et \num{0.47}), le verdict n'en dépend pas ;
le critère de masse tranche seul de toute façon.

\begin{table}[ht]
  \centering\small
  \caption{Étape 0 de la V5, sans entraînement (\code{bench-v4}, 2800 images entraînées, $\phi=\num{0.25}$).
  Précision : part des images sélectionnées sur lesquelles la doublure se trompe (argmax de son mélange
  de vues différent de la pseudo-étiquette), et taux de base entre parenthèses ; part de la perte :
  part de la perte moyenne de la doublure portée par la sélection, et meilleur choix de même taille ;
  direction : $S_{\text{pred}}$ à logits libres, classes fausses formées comme les secondes de
  l'original (S.) ou comme les vues de la V5 (V.). Les recadrages sont circulaires.}\label{tab:v5step0}
  \sisetup{table-format=1.3}\setlength{\tabcolsep}{4pt}
  \begin{tabular}{@{}l S S[table-format=1.3] S[table-format=1.2] S S[table-format=1.3] S S@{}}
    \toprule
    Doublure & {Précision} & {(base)} & {Préc./base} & {Part perte} & {(oracle)} & {Dir. S.} & {Dir. V.}\\
    \midrule
    $s_{40}$ & 0.144 & 0.068 & 2.13 & 0.406 & 0.741 & 0.714 & 0.802\\
    $s_{50}$ & 0.050 & 0.021 & 2.41 & 0.441 & 0.860 & 0.705 & 0.797\\
    $s_{60}$ & 0.027 & 0.010 & 2.62 & 0.465 & 0.924 & 0.712 & 0.775\\
    \midrule
    $m_{30}$ & 0.183 & 0.190 & 0.96 & 0.236 & 0.919 & 0.752 & 0.637\\
    $m_{20}$ & 0.109 & 0.118 & 0.92 & 0.231 & 0.976 & 0.765 & 0.609\\
    $m_{15}$ & 0.087 & 0.086 & 1.01 & 0.238 & 0.991 & 0.768 & 0.594\\
    \bottomrule
  \end{tabular}
\end{table}

\paragraph{La direction, elle, tient.} À logits libres, une fois la masse de la pseudo-étiquette
égale à celle de la doublure image par image, former les classes fausses comme l'original sur l'image
exacte (ce que fait la montée au premier ordre) donne un $S_{\text{pred}}$ de \numrange{0.75}{0.77}
sous le masquage, contre \numrange{0.59}{0.64} avec la forme des vues. Sous le recadrage, les vues
gagnent (\numrange{0.78}{0.80} contre \num{0.71}), ce qui est circulaire. L'option de
direction (\cref{eq:v5dir}) n'a donc pas lieu d'être codée. Mais ce chiffre suppose de savoir sur
quelles images la doublure doute, et c'est ce que la sélection ne sait pas.

\paragraph{Les deux familles ne s'accordent pas sur les images difficiles.} Une mesure sans
entraînement, posée avant d'être lancée, départage deux lectures : soit seules les vues par
recadrage prédisent mal la difficulté, soit les deux familles de doublures désignent des images
différentes. Le script \chemin{scripts/unlearning_v5_families.py} (sortie
\chemin{experiments/unlearning/measurements/v5_families.json}, \SI{55}{\second}) compare, image par
image, les pertes des six doublures et de deux sélecteurs (\cref{tab:v5familles}). Le premier est
formé des vues de la V5 ; le second, de masquages de \SI{30}{\percent} tirés à des graines que le banc
n'utilise pas. Le sélecteur des vues redonne exactement les chiffres de l'étape 0. La corrélation de
rang de Spearman~\cite{spearman1904} entre tout recadrage et tout masquage vaut de \num{-0.04} à
\num{+0.05}. Dans une même famille, elle vaut \numrange{0.89}{0.93} entre masquages : la difficulté y
est presque une propriété de l'image. Entre recadrages, elle ne vaut que \numrange{0.28}{0.37} : la
difficulté y dépend surtout du tirage. Le sélecteur par masquage retrouve presque toutes les erreurs
des masquages (\numrange{3.7}{4.0} fois le hasard, le plafond étant 4) et en porte
\numrange{0.87}{0.96} de la perte, mais il fait le hasard sous les recadrages (\numrange{0.69}{1.17}
fois, \numrange{0.22}{0.23} de la perte). La lecture posée avant est la première : l'échec ne tient
pas au choix des recadrages. Toute sélection tirée de vues de l'original ne vaut que sous sa propre
famille, où elle est circulaire, et vaut le hasard sous l'autre. Comme \code{score\_min} retient la
pire doublure, le banc ne peut ni valider une telle sélection, ni lui donner de gain. Ce constat est
cohérent avec l'écart, mesuré plus haut, entre les gains des V2 et V3 sous recadrage et sous masquage
(\cref{sec:campagnev2,sec:campagnev3}).

\begin{table}[ht]
  \centering\small
  \caption{La difficulté par image selon la famille de doublures (sans entraînement, 2800 images).
  $\rho$ : corrélation de rang de Spearman des pertes par image (moyenne sur 4 vues, pseudo-étiquette)
  avec les doublures de chaque famille ; Préc./base et part de la perte : sélection des
  \SI{25}{\percent} d'images de plus forte perte, comme au \cref{tab:v5step0}. Pour les doublures
  prises comme sélecteurs, la doublure elle-même est exclue des fourchettes.}\label{tab:v5familles}
  \setlength{\tabcolsep}{3.4pt}
  \begin{tabular}{@{}l cc cc cc@{}}
    \toprule
    & \multicolumn{2}{c}{$\rho$} & \multicolumn{2}{c}{Préc./base} & \multicolumn{2}{c}{Part de la perte}\\
    \cmidrule(lr){2-3}\cmidrule(lr){4-5}\cmidrule(l){6-7}
    Sélecteur & recadr. & masq. & recadr. & masq. & recadr. & masq.\\
    \midrule
    vues de la V5 (recadrages) & \numrange{0.35}{0.37} & \numrange{0.00}{0.01} & \numrange{2.1}{2.6}
      & \numrange{0.92}{1.01} & \numrange{0.41}{0.47} & \numrange{0.23}{0.24}\\
    masquages, graines hors banc & \numrange{-0.03}{0.04} & \numrange{0.89}{0.92} & \numrange{0.69}{1.17}
      & \numrange{3.7}{4.0} & \numrange{0.22}{0.23} & \numrange{0.87}{0.96}\\
    doublure $s_{50}$ & \num{0.34} & \numrange{-0.01}{0.00} & \numrange{1.9}{2.5} & \numrange{0.98}{1.04}
      & \numrange{0.39}{0.44} & \numrange{0.25}{0.26}\\
    doublure $m_{20}$ & \numrange{-0.03}{0.04} & \numrange{0.91}{0.93} & \numrange{0.84}{1.10}
      & \numrange{3.6}{4.0} & \num{0.23} & \numrange{0.86}{0.97}\\
    \bottomrule
  \end{tabular}
\end{table}

\paragraph{La localité seule, mesurée ensuite.} La borne haute des images tenues, qui devait donner
la localité (\cref{eq:v5hold}), ne dépend pas de l'identité des images poussées ; et sous le masquage, la
sélection vaut un tirage au hasard. Sur décision de l'utilisateur, \code{band-ascent} a donc été lancé
tel quel comme \emph{mesure de mécanisme}, critères de localité de sa fiche gardés, critères de score
abandonnés, puis complété par un suivi posé avant d'être lancé : les deux ablations au même calendrier de
quatre niveaux, et une seconde graine (\code{band-locality}, \code{band-ascent-s1} ; huit runs en tout,
\cref{tab:v5bande}). Le rapport audit/exact y est lu \emph{à KL exacte égale}, car il croît avec la dose
pour toutes les méthodes (\cref{fig:v5localite}).

\begin{table}[ht]
  \centering\small
  \caption{La bande de la V5 entraînée (\code{bench-v4}, lr \num{e-4}, une passe par niveau). Niveaux :
  niveaux de dose montés par le run final (CV : choisis par la validation croisée) ; KL ex. : KL des
  images exactes entraînées à $T=2$ ; Aud./ex. : comme au \cref{tab:v4scrub} ; $S_{\text{pred}}$ :
  fourchette sur les six doublures. Le bras sans borne haute et sans validation croisée (dernière ligne)
  s'effondre : la porte de précision n'est pas franchie.}\label{tab:v5bande}
  \sisetup{table-format=1.3}\setlength{\tabcolsep}{3pt}
  \begin{tabular}{@{}l c c S S[table-format=1.2] S S c S@{}}
    \toprule
    Bras & Graine & Niveaux & {KL ex.} & {Aud./ex.} & {Recadr.} & {Masq.} & $S_{\text{pred}}$ & {Retain}\\
    \midrule
    défaut ($\phi=\num{0.25}$) & 0 & 1 (CV) & 0.074 & 0.25 & 0.005 & 0.000 & \numrange{0.00}{0.05} & 0.981\\
    défaut ($\phi=\num{0.25}$) & 1 & 2 (CV) & 0.119 & 0.39 & 0.038 & 0.005 & \numrange{0.01}{0.13} & 0.950\\
    sans sélection ($\phi=1$) & 0 & 1 (CV) & 0.107 & 0.35 & 0.042 & 0.062 & \numrange{0.07}{0.15} & 0.958\\
    sans borne haute & 0 & 0 (CV) & {--} & {--} & 0.000 & 0.000 & -- & 1.000\\
    \midrule
    quatre niveaux ($\phi=\num{0.25}$) & 0 & 4 & 0.213 & 0.54 & 0.077 & 0.017 & \numrange{0.03}{0.27} & 0.878\\
    quatre niveaux ($\phi=\num{0.25}$) & 1 & 4 & 0.271 & 0.74 & 0.117 & 0.046 & \numrange{0.09}{0.40} & 0.809\\
    quatre niveaux, $\phi=1$ & 0 & 4 & 0.418 & 0.90 & 0.122 & 0.040 & \numrange{0.08}{0.52} & 0.698\\
    quatre niveaux, sans borne haute & 0 & 4 & 2.711 & 1.01 & 0.000 & 0.000 & 0 & 0.143\\
    \bottomrule
  \end{tabular}
\end{table}

\begin{figure}[ht]
  \centering
  \begin{tikzpicture}
    \begin{axis}[width=0.86\linewidth,height=6cm,
        xlabel={KL des images exactes entraînées ($T=2$)},ylabel={audit / exact},
        xmin=0,xmax=0.45,ymin=0,ymax=1.2,
        xticklabel style={/pgf/number format/.cd,fixed,precision=2,use comma},
        yticklabel style={/pgf/number format/.cd,fixed,precision=1,use comma},
        legend columns=3,legend style={at={(0.5,-0.24)},anchor=north,font=\small,draw=none,
          /tikz/every even column/.append style={column sep=6pt}},
        grid=major,grid style={gray!18},axis line style={gray!60},tick style={gray!60},
        label style={font=\small},tick label style={font=\small}]
      \addplot[only marks,mark=*,mark size=2pt,color=bleu] coordinates
        {(0.074,0.246) (0.119,0.393) (0.213,0.538) (0.271,0.739)};
      \addlegendentry{bande V5, $\phi=\num{0.25}$}
      \addplot[only marks,mark=o,mark size=2.4pt,thick,color=bleu] coordinates
        {(0.107,0.354) (0.418,0.897)};
      \addlegendentry{bande V5, $\phi=1$}
      \addplot[only marks,mark=square*,mark size=1.8pt,color=sarcelle] coordinates
        {(0.035,0.228) (0.028,0.224) (0.088,0.468) (0.077,0.508) (0.165,0.800) (0.166,0.934)
         (0.171,0.816) (0.277,1.028) (0.261,1.093) (0.275,1.032)};
      \addlegendentry{SCRUB conforme, lr \numrange{e-4}{4e-4}}
      \addplot[only marks,mark=diamond*,mark size=2.4pt,color=ambre] coordinates
        {(0.257,0.878) (0.293,0.876) (0.232,0.869) (0.278,0.856) (0.266,0.864) (0.240,0.859)
         (0.284,0.857) (0.302,0.897) (0.296,0.909)};
      \addlegendentry{V1, V2, V3 retenues}
      \addplot[only marks,mark=triangle*,mark size=2.2pt,color=gray] coordinates
        {(0.128,0.553) (0.088,0.511) (0.144,0.596)};
      \addlegendentry{\code{scrub} proxy, lr \num{e-4}}
    \end{axis}
  \end{tikzpicture}
  \caption{Localité selon la dose, \code{bench-v4} (un point par run, registre). Pour toutes les
  méthodes, le fold d'audit bouge d'autant plus, relativement aux images entraînées, que l'oubli est
  fort. À KL exacte égale, la bande de la V5 est la plus locale, et la même avec ou sans sélection
  ($\phi=1$). Le bras effondré (KL \num{2.7}) est hors cadre.}
  \label{fig:v5localite}
\end{figure}

\paragraph{La localité est obtenue, la plus forte mesurée.} Au niveau gardé par la validation croisée,
le bras par défaut tient le fold d'audit à \num{0.25} et \num{0.39} fois les images exactes (graines 0
et 1), sous le seuil de \num{0.5}, et près des \num{0.22} prédits par \cref{eq:v5hold}. À KL exacte
égale, aucune méthode mesurée ne fait mieux (\cref{fig:v5localite}). Vers \num{0.08}, SCRUB conforme
(lr \num{2e-4}) est à \numrange{0.47}{0.51}. Vers \numrange{0.17}{0.27}, la bande est à
\numrange{0.54}{0.74}, SCRUB conforme à \numrange{0.80}{1.09} et la V1 à \num{0.87}. La validation croisée
tient le run stable (aucun arrêt dur au run final) et la dérive basse (retain \numrange{0.95}{0.98}).

\paragraph{Mais elle ne vient pas de là où le modèle la mettait.} Selon \cref{eq:v5hold}, la localité
devait venir de la retenue des \SI{75}{\percent} d'images non poussées. Or les points sans sélection
($\phi=1$, toutes les images dans leur bande) tombent sur la même courbe que ceux à $\phi=\num{0.25}$ :
à perte exacte comparable et quatre niveaux, le bras à $\phi=\num{0.25}$ n'est qu'à \num{0.60} fois
le bras $\phi=1$, pas \num{0.5}, et à KL exacte égale l'écart disparaît. La localité vient de la montée bornée
elle-même : chaque image s'arrête dans sa bande. La borne haute des images tenues est en revanche
indispensable. Sans elle, le sous-run s'arrête dès le premier niveau (dérive du pli témoin \num{0.15}
au pas 40, contre \num{0.03} au bout du niveau avec la borne). À quatre niveaux, le modèle s'effondre :
KL exacte \num{2.7}, accord \SI{53}{\percent}, porte non franchie. Cette ablation mêle la dose et la
localité ; elle montre seulement que la borne empêche l'effondrement.

\paragraph{Elle ne paie pas.} Aucun des huit runs ne dépasse \num{0.046} (V1 : \num{0.0637}). Le score
s'annule par la composante des prédictions : au bras par défaut, $S_{\text{ms}}$ vaut
\numrange{0.41}{0.70} sous les six doublures, au-dessus de la V1 (\numrange{0.34}{0.43}), mais
$S_{\text{pred}}$ vaut \numrange{0.00}{0.13}. Les doublures sont moins sûres que l'original sur
\emph{toutes} les images, typiques comprises ; ce mouvement commun est celui que la localité refuse.
C'est le prix que la conception annonçait (\cref{sec:v5}), désormais mesuré : presque toute la
composante des prédictions. Sur ce banc, la localité ne rapporte que par le retain. Deux réserves.
Elle porte sur les images d'audit, \emph{mémorisées} par l'original : sur une vue forte de ces images
(\code{bench.audit\_view\_kl}, hors score), la bande les déplace de \numrange{1.6}{2.4} fois la KL
exacte, comme SCRUB conforme et la V1 (\numrange{1.5}{1.9}). Et le point à quatre niveaux n'est pas
reproduit à \SI{20}{\percent} près par la graine 1 (\num{0.74} contre \num{0.54}), mais cette graine
dérive davantage et reste sur la même courbe.

\paragraph{La dose choisie hors pli.} La validation croisée garde un niveau (graine 0) ou deux
(graine 1) : chaque fois, le niveau suivant déclenche un arrêt dur dans le sous-run, au dixième pas. Son
estimateur de dérive, la dérive du pli témoin corrigée de $(5/4)^2$, surestime celle du fold d'audit de
\numrange{1.5}{2.2} fois : le rapport mesuré audit sur témoin vaut \numrange{0.70}{1.04}, et la
correction d'échelle est trop prudente. Le choix n'en dépend pas : au niveau refusé, la dérive brute du
témoin dépasse déjà le budget.

\begin{retenir}[Verdict sur la V5]
\textbf{Abandonnée : la sélection est fausse et la localité ne paie pas.} Sous la famille de doublures
que ses vues ne partagent pas, la sélection vaut un tirage au hasard (étape 0), et les deux familles
du banc ne s'accordent pas sur les images difficiles. Entraînée comme mesure de mécanisme, la bande
donne la localité la plus forte mesurée à KL exacte égale (audit/exact \num{0.25} et \num{0.39} au niveau
choisi), sans rien devoir à la sélection ni à la retenue ; sa borne haute empêche l'effondrement.
Mais son score est nul, faute du mouvement commun que demande la composante des prédictions : la
localité, levier visé depuis la V4, ne rapporte sur ce banc que par le retain. La meilleure méthode
conforme mesurée reste SCRUB conforme (\num{0.1085 +- 0.0068}).
\end{retenir}

\subsection{Campagne V\_Dumb : pas de dose utilisable}\label{sec:campagnevdumb}
V\_Dumb (\cref{sec:vdumb}) a été mesurée le 22 novembre par onze runs sur la RTX 5080, graine 0 (fiche
\chemin{experiments/unlearning/trust-region-vdumb-campaign.md} ; fiches \chemin{v-dumb.md},
\chemin{v-dumb-dose.md}, \chemin{v-dumb-dose2.md}). Sa fiche posait, avant tout run, trois critères de
réfutation (meilleur bras sous \num{0.0577}, la V1 moins $2\sigma$ ; son retain sous \num{0.8} ; l'ancre
inutile si le bras sans elle garde le même retain). La fiche de campagne y ajoutait la comparaison
appariée à SCRUB conforme, qui a la même ancre de vues, et une règle des suivis : redoser si la cible
n'est pas atteinte, chaque suivi étant posé avant d'être lancé.

\paragraph{À sa fiche, le modèle ne bouge pas.} Au taux \num{e-5}, avec l'ancre de poids 1, les trois
cibles donnent le même modèle, à peine déplacé (KL exacte \num{0.0005}, perte \num{0.0013}, score
nul) ; les bras $t=\num{0.31}$ et $\num{0.45}$ sont identiques au bit près, la charnière ne liant
jamais. Sans ancre, la montée démarre à la fin de la première passe et s'emballe à la deuxième : perte
finale \num{12.2}, 27 fois la cible, porte non franchie. Réfutée sur son premier critère ; l'ancre
sert, au sens du troisième.

\paragraph{Entre les deux, une falaise.} La règle des suivis prévoyait de redoser. La mesure en
donnait la direction : le poids de l'ancre mène de l'immobilité à l'effondrement. Deux suivis ont
parcouru ce poids au pas logarithmique, et un bras a monté le taux à \num{e-4} (\cref{tab:vdumb}).
L'oubli croît quand l'ancre faiblit, lentement (KL exacte multipliée par 5 quand le poids est divisé par
10), jusqu'à $w=\num{0.01}$ (\num{0.067}, score \num{0.019}). À $w=\num{0.003}$, le modèle s'effondre :
la KL de l'ancre passe de \num{0.05} à \num{9.3} en deux passes. Au taux \num{e-4}, l'ancre se tend (KL
\numrange{0.14}{0.30} pendant l'entraînement) mais l'emporte. Aucun bras ne dépasse \num{0.0577}, et le
critère d'arrêt posé avant le dernier suivi est atteint : V\_Dumb n'a pas de dose utilisable.

\begin{table}[ht]
  \centering\small
  \caption{V\_Dumb (\code{bench-v4}, graine 0, $t=\num{0.45}$ sauf mention, trois passes). $w$ : poids
  de l'ancre (\code{retain\_weight}) ; atteint : part des images à la cible ($L\ge t$) ; KL ex. : KL
  des images exactes entraînées à $T=2$.}\label{tab:vdumb}
  \sisetup{table-format=1.3}\setlength{\tabcolsep}{3.5pt}
  \begin{tabular}{@{}l S[table-format=1.4] S[table-format=2.2] S[table-format=1.3] S[table-format=1.4] S c@{}}
    \toprule
    Réglage & {KL ex.} & {Perte} & {Atteint} & {\code{score\_min}} & {Retain} & Porte\\
    \midrule
    fiche : lr \num{e-5}, $w=1$, $t\in\{\num{0.31};\num{0.45};\num{0.63}\}$ & 0.0005 & 0.00 & 0.001 & 0.0000 & 1.000 & oui\\
    lr \num{e-4}, $w=1$ & 0.009 & 0.02 & 0.007 & 0.0031 & 0.998 & oui\\
    $w=\num{0.3}$ & 0.006 & 0.01 & 0.009 & 0.0001 & 1.000 & oui\\
    $w=\num{0.1}$ & 0.017 & 0.04 & 0.013 & 0.0008 & 1.000 & oui\\
    $w=\num{0.03}$ & 0.030 & 0.06 & 0.021 & 0.0035 & 0.999 & oui\\
    $w=\num{0.01}$ & 0.067 & 0.12 & 0.033 & 0.0194 & 0.987 & oui\\
    $w=\num{0.003}$ & 3.456 & 5.94 & 0.805 & 0.0000 & 0.052 & non\\
    $w=\num{0.001}$ & 5.496 & 9.94 & 0.954 & 0.0000 & 0.012 & non\\
    $w=0$ (fiche, sans ancre) & 6.531 & 12.23 & 0.977 & 0.0000 & 0.006 & non\\
    \bottomrule
  \end{tabular}
\end{table}

\paragraph{La charnière ne borne rien.} L'idée de V\_Dumb était de monter jusqu'à $t$, pas au-delà.
Avec l'ancre, aucune image ou presque n'atteint $t$ (\SI{3}{\percent} au plus) ; sans elle, ou presque,
la perte dépasse $t$ de 13 à 27 fois. Une image arrivée n'a plus de gradient propre, mais les mises à
jour des autres l'emportent : c'est le mouvement commun déjà mesuré (\cref{sec:diagnostic,sec:campagnev5}).
Le pronostic de la conception, des pertes égalisées, ne se vérifie donc pas. Au meilleur bras, comme pour
la bande de la V5, l'oubli est faible et local (audit/exact \num{0.19} à KL exacte \num{0.067}), et
$S_{\text{pred}}$ reste à \numrange{0.04}{0.12}.

\begin{retenir}[Verdict sur V\_Dumb]
\textbf{Abandonnée.} La méthode la plus simple ne bat pas la V1 : elle n'a pas de régime entre
l'immobilité (ancre) et l'effondrement (sans ancre), et sa charnière ne borne pas l'oubli. Les cibles
des V1 à V3 valent donc quelque chose sur ce banc, au sens posé par sa fiche. SCRUB conforme, qui a la
même ancre de vues, score \num{0.1085 +- 0.0068} : il alterne des pas de montée plafonnée et des pas
d'ancre, quand V\_Dumb les somme dans chaque pas. C'est une différence de conception, pas une cause
mesurée. Une seule graine et un seul taux pour le balayage du poids : une autre combinaison pourrait
trouver un régime, sans indice mesuré qu'elle existe.
\end{retenir}

\subsection{Campagne V5.1 : la marge dépasse la bande}\label{sec:campagnev51}
La V5.1 (\cref{sec:v51}) a été mesurée le 22 novembre sur la RTX 5080, graine 0, par un diagnostic et
onze runs (fiche \chemin{experiments/unlearning/trust-region-v51-campaign.md} ; fiches
\chemin{band-ascent-step0.md}, \chemin{band-ascent-budget.md}, \chemin{band-ascent-v51.md},
\chemin{band-ascent-v51-lr.md}). Sa base, la V5, était déjà abandonnée : elle a été lancée à la demande
de l'utilisateur, aux seuils que ses fiches posaient avant tout run. Le verdict de l'étape 0 de la V5,
une sélection au hasard sous le masquage, vaut aussi pour elle, qui garde la même sélection.

\paragraph{Les trois défauts sont confirmés sur l'original.} L'étape 0 rejouée lit les défauts sans
entraînement :
\begin{itemize}[leftmargin=1.2em,itemsep=1pt]
  \item la famine : \SI{10}{\percent} des images libérées portent \SI{87}{\percent} de la force de la
    montée en perte, et son rapport de participation vaut \num{0.03}, pour 1 si elle était également
    répartie (seuils posés : \num{0.3} et \num{0.5}) ;
  \item la borne $\delta_0$ ne serre pas au budget : une image à cette borne dérive de \num{0.37} en
    KL médiane, pour un budget brut de \num{0.064} ;
  \item la poussée au départ : la force qui rend les images tenues plus sûres vaut \num{1.94} fois
    la montée.
\end{itemize}
Le vote kNN de l'encodeur sélectionne mieux que les vues de la V5, mais peu : à l'entrée de la tête,
bloc choisi par la règle posée avant, sa précision fait \numrange{1.35}{1.59} fois le taux de base
sous les masquages (moyenne \num{0.187} contre \num{0.126}). Il ne porte que
\numrange{0.35}{0.40} de la perte des doublures (critère de masse de la V5 : \num{0.5}), et sa
direction perd sur les trois masquages. L'option n'est pas codée.

\paragraph{Dépenser le budget ne donne pas de score.} Sans code nouveau, \code{band-ascent-budget}
élargit le budget du témoin de la V5 (\cref{tab:v51}). La validation croisée garde alors plus de
niveaux (1, 3, puis 4), et le retain du banc tient (\numrange{0.88}{0.98}). Aucun critère de
réfutation n'est atteint, mais le score reste nul (\num{0.017} au mieux) : seuls
\SIrange{0.6}{4.6}{\percent} des images libérées atteignent leur bande, c'est la famine. Le bras
\num{0.1} redonne la V5 par défaut au bit près (192 métriques communes) : la fusion n'a pas changé la
V5.

\paragraph{La marge fait dépasser la bande.} Avec la force d'ordre 1 de la marge, les trois bras en
marge de \code{band-ascent-v51} s'arrêtent net au premier niveau, à lr \num{e-4} comme à
\num{3e-5}. En 15 à 40 pas, \SIrange{44}{74}{\percent} de la sonde passe au-dessus de sa borne, et
la dérive du témoin dépasse son arrêt dur. Aucun niveau n'est gardé : le run final ne bouge pas le
modèle. La fiche prescrivait de refaire à \num{3e-5} avant d'abandonner, et le suivi ajoutait
\num{e-5}. À \num{e-5} seulement, la V5.1 complète garde un niveau et score \num{0.0796}, au-dessus
de V1 $+2\sigma$. Mais son retain tombe à \num{0.725}, sous le plancher, et son rapport audit/exact
vaut \num{0.84}, comme la V3 : elle est réfutée sur deux critères.

Le témoin, qui annonçait un retain de \num{0.861}, se trompe de \num{+0.136}. Le run final, sur
2800 images, dérive bien plus que le sous-run sur quatre plis. La retenue KL ne tient pas non plus à
pleine échelle : les images tenues finissent à une KL de \num{0.299}, plus du double de la tolérance
de \num{0.128}.

\begin{table}[ht]
  \centering\small
  \caption{Campagne V5.1 (\code{bench-v4}, graine 0, RTX 5080, commit \code{ceae862}). Budget : le
  \code{cv\_max\_drift} de la V5 ; niveaux : niveaux de dose gardés (0 : modèle intact) ; écart du
  témoin : \code{final.cv\_retain} moins \code{bench.retain}.}\label{tab:v51}
  \setlength{\tabcolsep}{3.5pt}
  \begin{tabular}{@{}lcccccccc@{}}
    \toprule
    Réglage & lr & Niveaux & KL ex. & Audit/exact & \code{score\_min} & Retain & Écart du témoin\\
    \midrule
    V5, budget \num{0.1} & \num{e-4} & 1 & \num{0.074} & \num{0.25} & \num{0.000} & \num{0.981} & \num{-0.025}\\
    V5, budget \num{0.2} & \num{e-4} & 3 & \num{0.181} & \num{0.44} & \num{0.006} & \num{0.916} & \num{-0.038}\\
    V5, budget \num{0.3} & \num{e-4} & 4 & \num{0.213} & \num{0.54} & \num{0.017} & \num{0.878} & \num{-0.083}\\
    \midrule
    V5.1 complète & \num{e-4} & 0 & -- & -- & 0 & 1 & --\\
    \quad sans retenue KL & \num{e-4} & 0 & -- & -- & 0 & 1 & --\\
    \quad une passe par niveau & \num{e-4} & 0 & -- & -- & 0 & 1 & --\\
    \quad bande en perte & \num{e-4} & 4 & \num{0.190} & \num{0.54} & \num{0.011} & \num{0.895} & \num{-0.063}\\
    \midrule
    V5.1 complète & \num{3e-5} & 0 & -- & -- & 0 & 1 & --\\
    \quad sans retenue KL & \num{3e-5} & 0 & -- & -- & 0 & 1 & --\\
    \quad bande en perte & \num{3e-5} & 4 & \num{0.067} & \num{0.34} & \num{0.039} & \num{0.975} & \num{-0.015}\\
    V5.1 complète & \num{e-5} & 1 & \num{0.339} & \num{0.84} & \num{0.080} & \num{0.725} & \num{+0.136}\\
    \bottomrule
  \end{tabular}
\end{table}

\begin{retenir}[Verdict sur la V5.1]
\textbf{Abandonnée.} Les trois défauts relevés sans run sont réels. Les corriger ne suffit pas : la
force d'ordre 1 de la marge fait dépasser la bande en quelques pas. Au seul taux où la V5.1 garde une
dose, elle score au niveau de la V3, mais sous le plancher de retain et sans oubli local, et son témoin
ne protège pas le run final. Ce qui marche le moins mal est la retenue KL en espace de perte à
\num{3e-5} : un retain de \num{0.975} et \num{0.039}, sous la V1. La V5 manquait du mouvement commun
que demande la composante des prédictions (\cref{sec:campagnev5}) ; la V5.1 n'y change rien. Une
graine par réglage : les arrêts durs et le plancher manqué sont nets, les petits écarts entre bras ne
le sont pas.
\end{retenir}

\subsection{Campagne V\_Fusion : une ancre qui ne voit pas les dégâts}\label{sec:campagnevfusion}
V\_Fusion (\cref{sec:vfusion}) a été mesurée le 22 novembre sur la RTX 5080, en deux temps, chacun
sous des seuils posés avant (fiche \chemin{experiments/unlearning/trust-region-vfusion-campaign.md} ;
fiches \chemin{v-fusion-step0.md}, \chemin{v-fusion-anchor.md}).

\paragraph{Étape 0 : des images saturées.} La génération a pris \SI{34}{\minute} pour 2000 images, en
une seule tranche ; le job JURECA en prévoit quatre, et ses images ne seraient pas identiques au bit
près. Le forget set compte 45 classes de pseudo-étiquettes, dont 9 mises à part. Le lot A a donc
720 images sur 36 classes, le lot B 1100 sur 55, le lot C 180 ; l'ancre en garde 540 et 825, hors
thermomètre. Le plafond de confiance n'a joué pour aucune image : toutes l'atteignent, mais leur
confiance finale, lue sur les images en uint8 comme le serveur, vaut \numrange{0.998}{1}, au lieu de
\num{0.9}. Le plafond n'est testé qu'au dernier étage, et les images sont déjà saturées après les
étages à 56 et \SI{112}{px}.

Les portes a, b, c et c' passent pour A et B (\cref{tab:vfusionportes}). La porte d est
ininterprétable : le témoin, qui répare sur des vues du forget set, ne remonte qu'à \num{0.911}, puis
à \num{0.921} avec 1000 pas (suivi prévu), sous \num{0.95}. Par la règle posée avant, A et B entrent
dans l'ancre. Lue sans seuil, la porte d dit pourtant l'essentiel. Les images synthétiques réparent
environ trois fois moins que les vues (\SIrange{17}{21}{\percent} des dégâts, contre
\SIrange{55}{60}{\percent}). Un diagnostic l'explique : sur elles, le modèle bruité est déjà d'accord
avec l'original (perte de distillation de départ \numrange{0.04}{0.06}, contre \num{2.17} sur les
vues).

\begin{table}[ht]
  \centering\small
  \caption{V\_Fusion, portes de l'étape 0 (RTX 5080, commit \code{ceae862}, hors registre). d :
  accord sur des vues de la moitié d'évaluation du forget set, après 300 et 1000 pas de réparation ; le
  bruit le ramène à \num{0.800}.}\label{tab:vfusionportes}
  \setlength{\tabcolsep}{4pt}
  \begin{tabular}{@{}lcccccc@{}}
    \toprule
    Lot & Images & a : retourn. / recadr. / flou & b & c & c' & d (300 ; 1000 pas)\\
    \midrule
    A & 720 & 1 / 1 / 1 & \num{1.00} & \num{0.020} & \num{0.000} & \num{0.834} ; \num{0.842}\\
    B & 1100 & 1 / 1 / 1 & -- & \num{0.034} & -- & \num{0.834} ; \num{0.833}\\
    C & 180 & 1 / 1 / 1 & \num{1.00} & \num{0.010} & -- & --\\
    témoin (vues) & 1750 & & & & & \num{0.911} ; \num{0.921}\\
    seuil & & ${}\ge\num{0.8}$ & ${}\ge\num{0.7}$ & ${}\le\num{0.05}$ & ${}\le\num{0.10}$ & ${}\ge\num{0.95}$\\
    \bottomrule
  \end{tabular}
\end{table}

\paragraph{Les bras d'ancre s'effondrent.} Les quatre bras de \code{v-fusion-anchor} (lots AB, A, B à
lr \num{4e-4} ; AB à \num{3e-4} ; \cref{tab:vfusion}) finissent avec un retain de
\numrange{0.010}{0.025}, une KL exacte autour de 4 et la porte du banc non franchie : aucun ne tient
le plancher, V\_Fusion est réfutée. Le mécanisme se lit dans l'historique. La première époque de
montée est la même qu'avec l'ancre de vues, au bit près, puisqu'aucune ancre n'a encore agi. Sur les
vues, les min-steps voient ensuite une perte de \num{6.6} et la ramènent à \num{1.8}. Sur les images
synthétiques, ils ne voient que \numrange{0.28}{0.51}, puis \num{0.03} : le modèle cassé reste d'accord
avec l'original sur ces images, qui ne contraignent donc pas son comportement sur des images réelles.

\begin{table}[ht]
  \centering\small
  \caption{V\_Fusion, bras d'ancre (\code{bench-v4}, graine 0, RTX 5080, commit \code{ba05fcb}) et
  SCRUB conforme apparié (ancre de vues, \cref{sec:campagnev4}). Min-steps : perte de l'ancre, de la
  première à la dernière époque.}\label{tab:vfusion}
  \setlength{\tabcolsep}{4pt}
  \begin{tabular}{@{}lccccccc@{}}
    \toprule
    Ancre & lr & \code{score\_min} & Retain & KL ex. & Audit/exact & Porte & Min-steps\\
    \midrule
    synthétique AB & \num{4e-4} & 0 & \num{0.025} & \num{3.85} & \num{0.99} & non & \num{0.51} $\to$ \num{0.04}\\
    synthétique A & \num{4e-4} & 0 & \num{0.018} & \num{4.18} & \num{1.00} & non & \num{0.44} $\to$ \num{0.02}\\
    synthétique B & \num{4e-4} & 0 & \num{0.010} & \num{4.59} & \num{0.99} & non & \num{0.38} $\to$ \num{0.03}\\
    synthétique AB & \num{3e-4} & 0 & \num{0.022} & \num{3.94} & \num{1.01} & non & \num{0.28} $\to$ \num{0.03}\\
    \midrule
    vues (SCRUB conforme) & \num{4e-4} & \num{0.1279} & \num{0.734} & \num{0.277} & \num{1.03} & oui & \num{6.60} $\to$ \num{1.81}\\
    vues (SCRUB conforme) & \num{3e-4} & \num{0.1151} & \num{0.868} & \num{0.165} & \num{0.80} & oui & \num{5.16} $\to$ \num{1.28}\\
    \bottomrule
  \end{tabular}
\end{table}

\begin{retenir}[Verdict sur V\_Fusion]
\textbf{Abandonnée sous cette génération.} Une ancre qui ne voit pas les dégâts ne les répare pas.
La montée de SCRUB casse le modèle sur les images réelles, pas sur ces images synthétiques saturées.
Ce n'est pas une preuve contre toute ancre synthétique : la génération n'a pas tenu son plafond de
confiance, et une seule a été faite. La suite logique, non lancée (\cref{tab:plan}), commence par une
porte de sensibilité, mesurable en deux minutes sans entraînement. Elle compare, sur le modèle
bruité de la porte d, la perte de distillation sur l'ancre et sur des vues du forget set. Le meilleur
conforme reste SCRUB conforme, à \num{0.1085 +- 0.0068}.
\end{retenir}

\subsection{Meilleure performance au regard de la métrique}\label{sec:meilleure}
Le score officiel n'est pas mesuré : il demande les modèles gold et un control set que nous
n'avons pas, et rien n'a été soumis. Le banc applique ses formules exactes
(\cref{eq:ms,eq:pred,eq:final,eq:retain}) sur la même échelle (0 pour l'original, 1 pour le gold),
en remplaçant le gold par l'original sur des vues fortes. Le \cref{tab:meilleure} décompose les
meilleurs runs : \enquote{oubli} est le produit de la composante des pertes et de celle des
prédictions, sous la doublure la plus défavorable au run.

\begin{table}[ht]
  \centering\small
  \caption{Meilleurs runs du banc, décomposés sous leur pire doublure (RTX 5080, une graine
  chacun). Score $=$ oubli $\times$ retain ; oubli $=$ pertes ($S_{\text{ms}}$) $\times$ prédictions
  ($S_{\text{pred}}$).}\label{tab:meilleure}
  \sisetup{table-format=1.3}\setlength{\tabcolsep}{4.5pt}
  \begin{tabular}{@{}ll c SSSS S[table-format=1.3]@{}}
    \toprule
    Méthode & Réglage & Pire & {\code{score\_min}} & {Oubli} & {$S_{\text{ms}}$} & {$S_{\text{pred}}$}
      & {Retain}\\
    \midrule
    scrub & lr \num{2e-4}, graine 0 & $s_{60}$ & 0.118 & 0.178 & 0.745 & 0.239 & 0.663\\
    scrub & lr \num{e-4}, graine 2 & $s_{40}$ & 0.117 & 0.129 & 0.340 & 0.378 & 0.910\\
    trust\_region V2 & $\tau=1$, $\varepsilon=\num{0.15}$, graine 1 & $s_{50}$ & 0.088 & 0.113 & 0.352 & 0.322 & 0.777\\
    gold\_soft & 1 époque, \num{e-5}, graine 0 & $s_{60}$ & 0.070 & 0.082 & 0.309 & 0.264 & 0.854\\
    trust\_region V1 & $\varepsilon=\num{0.15}$, graine 0 & $s_{60}$ & 0.065 & 0.084 & 0.456 & 0.184 & 0.771\\
    \bottomrule
  \end{tabular}
\end{table}

Sous une seule doublure, le meilleur run atteint \num{0.165} ($s_{40}$), \num{0.164} ($s_{50}$) et
\num{0.156} ($s_{60}$). Trois constats. \textbf{On est à environ un dixième du chemin} de
l'original vers la doublure du gold. \textbf{La composante des prédictions limite} (\numrange{0.18}{0.38}) :
les modèles deviennent moins sûrs d'eux, mais ne se trompent pas vers les mêmes classes que la
doublure ; celle des pertes va plus loin (jusqu'à \num{0.745}). \textbf{Le retain se paie} : le
meilleur run garde \SI{94}{\percent} d'accord sur le fold d'audit mais une divergence de
\num{0.42}, et le réglage qui le produit est à un pas de l'effondrement (\cref{tab:refs}). Sur trois
graines, deux réglages se détachent et ne se départagent pas : \code{scrub} à lr \num{e-4}
(\num{0.098 +- 0.022}, retain \num{0.93}) et la V2 à $\tau=1$ (\num{0.0831 +- 0.0042}, retain
\num{0.78}), le premier meilleur en moyenne, le second cinq fois plus stable. Des deux, seule la V2 est
\emph{conforme} : à leur défaut, \code{scrub} et \code{gold\_soft} ancrent leur rétention sur des images
extérieures au forget set, que le sujet exclut pour une méthode soumise (\cref{sec:pistes}) ; ce
sont des références de banc. La meilleure méthode conforme mesurée est donc la V2 à $\tau=1$ :
\num{0.0831 +- 0.0042} sous recadrage, \num{0.0786 +- 0.0030} sous la pire des six doublures de
\code{bench-v4} (\cref{tab:v4}) ; son meilleur run vaut \num{0.088}. La V3 fait jeu égal en
\code{bench-v4} (\num{0.0798 +- 0.0012}, meilleur run \num{0.0806}, \cref{sec:campagnev3}).

Depuis la V4, la meilleure méthode conforme mesurée est \code{scrub} ancré sur des vues fortes du
forget set, lr \num{3e-4} (\cref{sec:campagnev4}) : \num{0.1085 +- 0.0068} sous la pire des six
doublures de \code{bench-v4}, meilleur run \num{0.115}. Ce run (graine 0), décomposé sous sa pire
doublure ($m_{30}$) : oubli \num{0.133} = pertes \num{0.380} $\times$ prédictions \num{0.348}, retain
\num{0.868}. Sous une seule doublure, il atteint \num{0.185} ($s_{50}$). La composante des prédictions
reste celle qui limite, mais elle y est plus haute que pour la V3 (\num{0.234} sous $m_{15}$). Aucun de
ces fichiers n'est soumissible : ils n'ont pas oublié le fold d'audit.

\subsection{Inventaire des expériences}\label{sec:inventaire}
Les \cref{tab:inventaire,tab:inventaire4} listent toutes les spécifications lancées à pleine échelle
(189 runs d'entraînement et les deux étapes 0 de la V5, RTX 5080, 21 et 22 novembre ; l'étape 0 de
V\_Fusion, hors registre, est décrite au \cref{sec:campagnevfusion}) et leur issue ; les chiffres complets sont dans
\chemin{experiments/unlearning/SCOREBOARD.md}.

\begin{table}[ht]
  \centering\small
  \caption{Inventaire des runs à pleine échelle. \code{tr} : \code{trust\_region} V1 ; sauf mention,
  rayon testé à chaque pas, $\tau=6$, $\varepsilon=\num{0.15}$, graine 0.}\label{tab:inventaire}
  \begin{tabularx}{\linewidth}{@{}>{\raggedright}p{3.5cm}c>{\raggedright}p{4.1cm}X@{}}
    \toprule
    Spécification & Runs & Balayage & Issue (\code{score\_min})\\
    \midrule
    \code{bench-baselines} & 2 & \code{gold\_soft}, \code{scrub} (défauts) & 0 ; \num{0.0027}\\
    \code{trust-region-kl-radius} & 4 & tr, $\varepsilon$ \numrange{0.2}{0.8}, test tous les 10 pas & 0 partout ;
      \numrange{0.2}{0.6} identiques\\
    \code{bench-seed-noise} & 2 & tr, graines 1-2, test tous les 10 pas & 0 (porte cassée)\\
    \code{tr-radius-fine-small}, \code{-large} & 8 & $\varepsilon$ \numrange{0.01}{0.8} & pic \num{0.059}-\num{0.060}
      à \numrange{0.1}{0.2}\\
    \code{tr-seeds} & 3 & graines 0-2 & \num{0.0637 +- 0.0025}\\
    \code{scrub-seeds} & 2 & \code{scrub}, graines 1-2 & \num{0.0025 +- 0.0003} (3 graines)\\
    \code{tr-lr} & 3 & lr \num{3e-6}, \num{3e-5}, \num{e-4} & \numrange{0.062}{0.067}\\
    \code{tr-warmup} & 3 & part 0, \num{0.03}, \num{0.3} & \numrange{0.065}{0.066}\\
    \code{tr-temperature} & 4 & $\tau$ 4, 5, 8, 12 & \num{0.065} partout\\
    \code{tr-blocks} & 3 & blocs 0, 2, 6 & \num{0.0006} ; \num{0.025} ; \num{0.025}\\
    \code{tr-batch} & 3 & lot 8, 16, 64 & \numrange{0.063}{0.067}\\
    \code{tr-confirm-radius} & 4 & $\varepsilon$ \num{0.1} et \num{0.2}, graines 1-2 & \num{0.061 +- 0.003} ; \num{0.058 +- 0.006}\\
    \code{gold-soft-dose} & 4 & 1 époque, lr \num{5e-7} à \num{e-5} & \num{0.070} à \num{e-5}, $\le\num{0.001}$ sinon\\
    \code{scrub-dose} & 3 & lr \num{3e-5}, \num{e-4}, \num{3e-4} & \num{0.014} ; \num{0.102} ; 0\\
    \code{scrub-dose-seeds} & 2 & lr \num{e-4}, graines 1-2 & \num{0.098 +- 0.022} (3 graines)\\
    \code{scrub-dose-fine} & 2 & lr \num{5e-5}, \num{2e-4} & \num{0.037} ; \num{0.118}\\
    \code{gold-soft-dose-seeds} & 2 & 1 époque, \num{e-5}, graines 1-2 & \num{0.059 +- 0.012} (3 graines)\\
    \midrule
    \code{trust-region-views} (V2) & 0 & test tous les 10 pas, $\varepsilon\ge\num{0.2}$ & non lancée
      telle quelle (défaut d'arrêt de la V1) ; remplacée par les \code{tr2-*}\\
    \midrule
    \multicolumn{4}{@{}l}{\emph{Campagne V2} (aire \num{0.45}, 3 vues, $\varepsilon=\num{0.15}$,
      graine 0 sauf mention ; V1 au même point : \num{0.0647})}\\
    \code{tr2-radius-a}, \code{-b} & 6 & $\tau=4$, $\varepsilon$ \numrange{0.05}{0.5} & au-dessus de
      la V1 aux six rayons (\numrange{0.002}{0.006})\\
    \code{tr2-temperature}, \code{-low}, \code{-vlow} & 7 & $\tau$ de \num{0.25} à 6 &
      \num{0.066} à $\tau=6$ ; \num{0.081} à 1 ; \num{0.085} à \num{0.5}\\
    \code{tr2-area-t1}, \code{-t2} & 6 & aire \numrange{0.25}{0.65}, $\tau=1$ et 2 & presque plat\\
    \code{tr2-views} & 3 & 1, 6, 10 vues & \num{0.077} ; \num{0.080} ; \num{0.080}\\
    \code{tr2-radius-best} & 4 & $\varepsilon$ \num{0.1} et \num{0.2}, $\tau$ \num{0.5} et 1 &
      sous $\varepsilon=\num{0.15}$\\
    \code{tr2-seeds} & 4 & $\tau$ \num{0.5} et 1, graines 1-2 & $\tau=1$ : \num{0.0831 +- 0.0042}
      (3 graines)\\
    \bottomrule
  \end{tabularx}
\end{table}

\begin{table}[ht]
  \centering\small
  \caption{Inventaire des runs à pleine échelle (suite) : campagnes en \code{bench-v4}. Sauf
  mention, V3 retenue ($b=0$, $\lambda=\num{0.25}$, $\rho=\num{0.9}$, aire \num{0.5}, 4 vues), rayon
  testé à chaque pas, graine 0.}\label{tab:inventaire4}
  \begin{tabularx}{\linewidth}{@{}>{\raggedright}p{3.5cm}c>{\raggedright}p{4.1cm}X@{}}
    \toprule
    Spécification & Runs & Balayage & Issue (\code{score\_min})\\
    \midrule
    \multicolumn{4}{@{}l}{\emph{Campagne V3}, en \code{bench-v4} (graine 0 sauf mention)}\\
    \code{ref4-v1}, \code{-v2}, \code{-scrub} & 9 & références relancées, graines 0-2 &
      \num{0.0637} ; \num{0.0786} ; \num{0.098} (identiques au recalcul)\\
    \code{trust-region-calibrated} & 4 & $\lambda$ \num{0.25} et \num{0.5}, copies 0 et \num{0.5} &
      \num{0.076} ; \num{0.076} ; \num{0.057} ; \num{0.065}\\
    \code{tr3-lambda}, \code{-hedge} & 4 & $\lambda$ 0, \num{0.1} ; $b$ 0, \num{0.1} & \num{0.060} ;
      \num{0.068} ; \num{0.078} ; \num{0.077}\\
    \code{tr3-views-*} & 4 & 1, 2 et 8 vues & 8 : \num{0.074} ; 1-2 : rayon jamais atteint\\
    \code{tr3-radius}, \code{-fixed} & 3 & $\rho$ \num{0.6}, \num{0.75} ; $\varepsilon$ \num{0.15} &
      \num{0.072} ; \num{0.074} ; \num{0.074}\\
    \code{tr3-b0-*} & 6 & à $b=0$ : $\rho$ \num{0.95}, $\lambda$ \num{0.35}, $\varepsilon$ \num{0.15},
      aires \num{0.4} et \num{0.6}, 2 vues & au mieux \num{0.078} ; un run refusé ($\varepsilon$
      \num{0.2} au-delà de la cible)\\
    \code{tr3-seeds} & 2 & $b=0$, graines 1-2 & \num{0.0798 +- 0.0012} (3 graines)\\
    \midrule
    \multicolumn{4}{@{}l}{\emph{Campagne V4}, en \code{bench-v4} (\cref{sec:campagnev4})}\\
    \code{ref4-v3}, \code{ref4-scrub} & 6 & rejeux, graines 0-2 & identiques au bit près (plus 7
      doublons \code{dirty}, identiques aussi)\\
    \code{scrub-forget-views} & 4 & SCRUB conforme, lr \num{e-4} et \num{2e-4}, graines 0-1 &
      \num{0.017} ; \num{0.061} (réfutée)\\
    \code{scrub-fv-dose}, \code{-seeds} & 6 & lr \num{3e-4} et \num{4e-4}, graines 0-2 &
      \num{0.1085 +- 0.0068} ; \num{0.121 +- 0.007} (retain \num{0.73})\\
    \code{trust-region-dual} & 4 & $\eta$ \num{0.25} et \num{0.5}, graines 0-1 & \num{0.042}
      partout (ne lie pas)\\
    \code{trust-region-anchored} & 4 & ancre : base, $K=2$, vues, $w=3$ & \numrange{0.035}{0.055},
      aucun rayon atteint\\
    \code{tr4-anchored-lr} & 4 & mêmes bras exacts à lr \num{3e-5}, témoin & \numrange{0.058}{0.070} ;
      témoin \num{0.088}\\
    \code{tr3-lr-seeds} & 4 & V3 à lr \num{3e-5} et \num{e-4}, graines 1-2 & \num{3e-5} :
      \num{0.0845 +- 0.0028} (3 graines)\\
    \midrule
    \multicolumn{4}{@{}l}{\emph{Campagne V5}, en \code{bench-v4} (\cref{sec:campagnev5})}\\
    \code{band-ascent-step0} & 1 & sans entraînement, 2800 images & réfutée : sélection au hasard
      sous le masquage\\
    \code{band-ascent} & 4 & mesure de localité : défaut, $\phi=1$, sans borne haute, 4 niveaux &
      audit/exact \num{0.25} ; scores \numrange{0}{0.042}\\
    \code{band-locality}, \code{-ascent-s1} & 4 & ablations à 4 niveaux ; graine 1 & sans borne :
      effondrement ; $\phi=1$ aussi local ; au mieux \num{0.046}\\
    \midrule
    \multicolumn{4}{@{}l}{\emph{Campagne V\_Dumb}, en \code{bench-v4} (\cref{sec:campagnevdumb})}\\
    \code{v-dumb} & 4 & $t$ \num{0.31}, \num{0.45}, \num{0.63} ; sans ancre & immobile ($\approx0$) ;
      sans ancre : effondrement\\
    \code{v-dumb-dose}, \code{-dose2} & 7 & poids d'ancre \num{0.3} à \num{0.001} ; lr \num{e-4} &
      au mieux \num{0.019} ($w=\num{0.01}$) ; effondrement dès $w=\num{0.003}$\\
    \midrule
    \multicolumn{4}{@{}l}{\emph{Campagne V5.1}, en \code{bench-v4} (\cref{sec:campagnev51})}\\
    \code{band-ascent-step0} & 1 & rejouée : encodeur et défauts, sans entraînement & trois défauts
      confirmés ; encodeur non codé\\
    \code{band-ascent-budget} & 3 & V5, budget du témoin \num{0.1}, \num{0.2}, \num{0.3} & budget
      consommé ; au mieux \num{0.017}\\
    \code{band-ascent-v51}, \code{-v51-lr} & 8 & V5.1 et ablations, lr \num{e-4}, \num{3e-5}, \num{e-5} &
      marge : arrêt dur ; à \num{e-5} : \num{0.080}, retain \num{0.725}\\
    \midrule
    \multicolumn{4}{@{}l}{\emph{Campagne V\_Fusion}, en \code{bench-v4} (\cref{sec:campagnevfusion})}\\
    \code{v-fusion-anchor} & 4 & ancre synthétique AB, A, B à lr \num{4e-4} ; AB à \num{3e-4} &
      effondrement (retain \numrange{0.010}{0.025})\\
    \bottomrule
  \end{tabularx}
\end{table}

\subsection{Vérification après la fusion}\label{sec:fusion}
La branche a été fusionnée avec la V2 de \code{trust\_region}, écrite en parallèle de la campagne
(\cref{sec:trustv2}) : son code est repris tel quel, et le réglage \code{radius\_check\_every} de la
campagne s'y ajoute sans changer aucun défaut. Sur le code fusionné, les 503 tests passent ; le
run de référence de la V1 ($\varepsilon=\num{0.15}$, graine 0) redonne exactement le même banc
(\code{score\_min} \num{0.0647}, retain \num{0.7706}) ; un essai réduit de la V2 tourne sur la RTX
5080 et retrouve la cible de l'essai fait sur Apple M2 (divergence \num{0.938}, accord
\SI{92.3}{\percent} sur 52 images). C'est une vérification de la chaîne, pas une mesure de la V2.

La troisième fusion (V3 et \code{bench-v4}, 22 novembre) a été vérifiée de même : 527 tests ; la V2
retenue, réentraînée sur le code fusionné (graine 0), redonne exactement son score sous recadrage
(\num{0.0811}) ; les fichiers V1, V2 et \code{scrub} recalculés en \code{bench-v4} redonnent
exactement leurs scores \code{bench-v3} sous les trois recadrages (\cref{tab:v4}).

\subsection{Mise en œuvre et coûts}\label{sec:couts}
Les mesures de coût ont guidé trois optimisations. Sur processeur, une vue forte coûtait
\SI{1.9}{\milli\second} par image (\SI{61}{\milli\second} pour un lot de 32), soit environ
\SI{250}{\milli\second} d'augmentations par pas de \code{gold\_soft} : les altérations s'exécutent
désormais sur le GPU, tirages aléatoires compris (générateur CPU, mêmes valeurs partout). Le
premier critère d'arrêt de \code{trust\_region} refaisait une passe complète tous les 10 pas, soit
environ trois fois le coût de l'entraînement ; l'échantillon fixe de 256 images le ramène à un
quart environ (estimation pour des lots de 32). Enfin AdamW utilise son implémentation fusionnée sur
GPU, et les runs de réglage n'évaluent qu'une fois, par le banc.

Sur la RTX 5080 (\SI{16}{\giga\byte}, partagée avec l'affichage), mesuré le 21 novembre avec la
config \code{full} : \SI{0.24}{\second} par pas de \code{gold\_soft} (\SI{64}{\second} pour 264
pas, \SI{107}{\second} de bout en bout avec le banc) et \SI{9}{\second} d'entraînement pour
\code{trust\_region} jusqu'au rayon \num{0.6} (\SI{27}{\second} au total). Le job occupe environ
\SI{7.3}{\giga\byte} de mémoire GPU et \SI{4.4}{\giga\byte} de RAM ; la carte monte à
\SI{62}{\celsius} et \SI{303}{\watt} au plus, sans aucun bridage. La config \code{full} tient donc
telle quelle et reste identique à celle de JURECA, ce qui garde les chiffres des deux machines
comparables. Tous les runs locaux passent par une garde (\chemin{scripts/gpu_guard.sh}) : un seul
job à la fois sur la carte, refus de démarrer sous \SI{9}{\giga\byte} libres, suspension du job à
\SI{83}{\celsius} et reprise à \SI{75}{\celsius} (limite de fonctionnement lue par
\code{nvidia-smi} : environ \SI{88}{\celsius}), télémétrie de chaque job. La campagne entière
(54 runs) a pris 37 minutes de GPU, de \SI{25}{\second} à \SI{2}{\minute} par run ; la carte n'a
pas dépassé \SI{60}{\celsius} et la garde n'a jamais dû suspendre un job. Tester le rayon à chaque
pas ajoute environ \SI{7}{\second} à un run de \code{trust\_region} (\SI{34}{\second} au lieu de
\SI{27}{\second}). Un run de la V2 prend \SI{40}{\second} (trois passes de vues sur le fold
entraîné, une douzaine de pas en plus) ; le calcul des vues culmine à \SI{6.8}{\giga\byte} de
mémoire GPU, et la campagne V2 entière a pris 27 minutes, dont trois runs de \SI{3}{\minute} qui
n'atteignaient pas leur rayon.

\section{Bilan par méthode}\label{sec:bilan}
\begin{table}[ht]
  \centering\small
  \caption{État de chaque méthode au 22 novembre, après les campagnes V5, V\_Dumb, V5.1 et V\_Fusion (pleine échelle sauf mention ; protocole indiqué
  pour chaque chiffre, comparaison sous un même protocole au \cref{tab:synthese}).}\label{tab:bilan}
  \begin{tabularx}{\linewidth}{@{}>{\raggedright}p{2.3cm}>{\raggedright}p{3.4cm}>{\raggedright}X>{\raggedright\arraybackslash}X@{}}
    \toprule
    Méthode & Meilleur réglage mesuré & Points forts mesurés & Points faibles mesurés\\
    \midrule
    \code{identity} & -- (0 par construction) & valide la chaîne & aucun oubli\\
    \code{scrub} & lr \num{e-4} : \num{0.098 +- 0.022} ; meilleur run \num{0.118} (lr \num{2e-4}) &
      meilleur score moyen ; meilleure localité (audit/exact \numrange{0.51}{0.60}, KL audit \num{0.067}) &
      bruité (écart-type \num{0.022}) ; effondrement à lr \num{3e-4} ; inutile à ses défauts ;
      \textbf{non conforme} (images proxy, \cref{sec:pistes})\\
    \code{scrub} conforme (V4) & vues du forget set, lr \num{3e-4} : \num{0.1085 +- 0.0068}
      (\code{bench-v4}) & \textbf{meilleure méthode conforme} ; \SI{36}{\percent} au-dessus de la V3 sous
      les deux familles, apparié ; retain \num{0.86} ; jeu égal avec le proxy & pas plus local que la V3
      à cette dose (audit/exact \num{0.85}) ; dose propre à l'ancre (inutile aux doses du proxy) ;
      \num{4e-4} : retain \num{0.73}\\
    \code{gold\_soft} & 1 époque, lr \num{e-5} : \num{0.059 +- 0.012} & jeu égal avec la V1 &
      ses défauts vont trop loin (retain \num{0.40}) ; bruité ; \textbf{non conforme} (images
      proxy)\\
    \code{trust\_region} V1 & $\varepsilon=\num{0.15}$, rayon testé à chaque pas :
      \num{0.0637 +- 0.0025} & reproductible (écart-type \num{0.0025}) ; un seul réglage utile ;
      insensible au taux, au warmup, à $\tau$ et au lot & localité (audit/exact \num{0.87}) ; le test tous
      les 10 pas ne contrôle rien\\
    \code{trust\_region} V2 & $\tau=1$, aire \num{0.45}, 3 vues, $\varepsilon=\num{0.15}$ :
      \num{0.0831 +- 0.0042} ; \code{bench-v4} : \num{0.0786 +- 0.0030} & meilleure méthode
      conforme ; au-dessus de la V1 sous les deux familles de doublures, apparié ; stable ; gain par
      la direction (prédictions) & localité (audit/exact \num{0.86}, comme la V1) ; $\tau\le\num{0.5}$
      instable ; gain réduit à \SI{12}{\percent} sous la famille étrangère à sa cible\\
    \code{trust\_region} V3 & $b=0$, $\lambda=\num{0.25}$, $\rho=\num{0.9}$, aire \num{0.5}, 4 vues :
      \num{0.0798 +- 0.0012} (\code{bench-v4}) & jeu égal avec la V2 ; la plus stable des trois &
      pas mieux que la V2 ; un peu plus de dérive (retain \num{0.74}, audit/exact
      \numrange{0.86}{0.91}) ; réglages couplés au rayon ; 1-2 vues : rayon jamais atteint\\
    \code{trust\_region} V4 & ancre : au mieux \num{0.070} ($w=3$, lr \num{3e-5}) ; contrainte active :
      \num{0.042} & options neutres à leur défaut (V1 à V3 rejouées au bit près) ; l'ancre baisse
      audit/exact d'un cinquième (lr \num{3e-5}) & ancre et contrainte réfutées : phases jamais
      arrêtées à leur rayon, contrainte qui ne lie pas, score divisé par deux ; gain de l'ancre
      sous le recadrage seulement\\
    \code{band\_ascent} (V5) & défaut : \numrange{0.000}{0.005} (2 graines) ; au mieux \num{0.046}
      (4 niveaux) & \textbf{la plus locale} à KL exacte égale (audit/exact \numrange{0.25}{0.39} au
      niveau choisi) ; retain \numrange{0.95}{0.98} ; pari de direction confirmé sous le masquage &
      score nul ($S_{\text{pred}}\approx0$) ; sélection au hasard sous le masquage ; la retenue
      n'explique pas la localité ; sans borne haute, effondrement ; estimateur hors pli
      \numrange{1.5}{2.2} fois trop haut\\
    \code{dumb} (V\_Dumb) & au mieux \num{0.019} ($w=\num{0.01}$, une graine) & la plus simple ;
      conforme au sujet ; retain \numrange{0.99}{1.00} tant que l'ancre tient & aucune dose
      utilisable : immobile avec l'ancre, effondrée dès $w=\num{0.003}$ ; la charnière ne borne pas
      l'oubli ; la cible $t$ sans effet\\
    \code{band\_ascent} V5.1 & lr \num{e-5} : \num{0.080}, retain \num{0.725} (une graine) &
      défauts de la V5 confirmés sur l'original ; bande en perte avec retenue KL (lr \num{3e-5}) :
      retain \num{0.975} & la marge fait dépasser la bande (arrêt dur à \num{e-4} et \num{3e-5}) ;
      sous le plancher et non locale (\num{0.84}) au seul taux qui garde un niveau ; témoin optimiste
      de \num{0.136} ; retenue KL non tenue à pleine échelle\\
    V\_Fusion & aucun réglage utile (quatre bras effondrés) & conforme au sujet ; portes a, b, c, c'
      passées & images saturées (plafond sans effet) ; le modèle cassé reste d'accord avec l'original
      sur elles : l'ancre ne répare rien (retain \numrange{0.010}{0.025})\\
    \bottomrule
  \end{tabularx}
\end{table}
V\_Dumb2 (\cref{sec:vdumb2}), codée sans run, n'a pas encore de mesure et n'entre pas au
tableau : le niveau par image de V\_Dumb dans l'alternance de SCRUB conforme, conforme au sujet, dont
la porte sans entraînement demande un plafond d'oubli des images exactes d'au moins \num{0.155}.


\section{Suite proposée}\label{sec:plan}
Les quatre campagnes ont tourné : V1 et V2 en \code{bench-v3} (recalculées en \code{bench-v4}), V3 et
V4 en \code{bench-v4}. La V5 a été abandonnée : sa sélection est réfutée, et sa localité, obtenue, ne
paie pas (\cref{sec:campagnev5}). V\_Dumb aussi : elle n'a pas de dose utilisable
(\cref{sec:campagnevdumb}). La V5.1 est abandonnée à son tour (\cref{sec:campagnev51}), et
V\_Fusion sous sa première génération (\cref{sec:campagnevfusion}). La meilleure méthode conforme mesurée est toujours \code{scrub} ancré sur des
vues fortes du forget set (\num{0.1085 +- 0.0068}, \cref{sec:campagnev4}). Les prochaines étapes,
proposées et non lancées, sont au \cref{tab:plan}, par valeur d'information. Chacune aura sa fiche,
avec hypothèses et critères de réfutation posés avant tout run, et se comparera, appariée à graine
égale, aux références mesurées en \code{bench-v4} : \code{ref4-scrub-fv} ($\sigma\approx\num{0.007}$ ;
seuil d'amélioration \num{0.122} sur trois graines) et \code{ref4-v3} ($\sigma\approx\num{0.003}$). Un
positif obtenu sur une seule graine se confirme sur trois avant d'entrer dans ce rapport. Pour
reprendre les expériences sur un autre GPU, la fiche \chemin{experiments/unlearning/ETAT.md}
rassemble l'état mesuré, les spécifications de référence à rejouer (\code{ref4-v1}, \code{-v2},
\code{-v3}, \code{-scrub}, \code{-scrub-fv}), les règles du protocole, les pièges mesurés et la recette
pour juger une nouvelle version. Tous les runs de réglage tiennent \SI{20}{\percent} du forget set à
l'écart : leurs fichiers ne doivent pas être soumis, et le run final d'une idée retenue se fera sans
fold d'audit.

\begin{table}[ht]
  \centering\small
  \caption{Suite proposée, par valeur d'information (4 variantes par spécification, un nœud ou la
  RTX 5080).}\label{tab:plan}
  \begin{tabularx}{\linewidth}{@{}>{\raggedright}p{4.2cm}X@{}}
    \toprule
    Expérience & Question et réfutation\\
    \midrule
    1. Run final de SCRUB conforme (décision d'équipe) & Sans fold d'audit, les 3500 images
      donnent plus de pas par époque : la dose \num{3e-4} tient-elle ? Arbitrer entre \num{3e-4} et
      \num{4e-4} (\num{+0.012} au banc, mais retain \num{0.73}), puis soumettre une fois : seul le
      classement public dira ce qui se transfère au vrai gold.\\
    2. V\_Dumb2, codée (\cref{sec:vdumb2}) & L'étape 0 (\code{band-ascent-step0}, clé
      \code{level\_ceiling}) d'abord : plafond d'oubli des images exactes à $t=\num{0.45}$ sous le pire
      masquage $\ge\num{0.155}$, sinon ne pas lancer. Puis \code{v-dumb2} (4 bras : lr \num{3e-4} et
      \num{4e-4}, avec et sans rappel) : réfutée si aucun bras de retain $\ge\num{0.8}$ n'atteint
      \num{0.129}, ou si le retain tombe sous \num{0.8} à \num{4e-4}.\\
    3. La localité de SCRUB conforme (valeur réduite depuis la V5) & Elle n'existe qu'aux doses où l'ancre défait l'oubli
      (audit/exact \numrange{0.22}{0.51} à \num{e-4} et \num{2e-4}) ; un réglage la garde-t-il à une
      dose qui score (ancre plus faible, moins de passes de réparation), ou faut-il un masque propre à
      l'instance (\cref{sec:pistes}) ? Une grille ne suffit pas : à poser à partir de la trajectoire de
      la KL d'audit pendant la réparation.\\
    Écartés par la V5 & Viser la localité pour elle-même : sur ce banc, elle ne rapporte que par le
      retain (\cref{sec:campagnev5}). Une sélection d'images tirée de vues de l'original : le banc ne
      peut pas la juger (ligne suivante).\\
    Un juge pour les sélections & Les deux familles du banc ne s'accordent pas sur les images
      difficiles (\cref{tab:v5familles}) : juger une méthode qui choisit ses images demande une
      troisième famille de doublures, ni recadrage ni masquage, et un sélecteur tiré d'une autre
      famille que ce juge. Laquelle ressemble au vrai gold, seul le serveur le dirait (décision
      d'équipe).\\
    Rayon de la V2 sur trois graines & Le plateau $\varepsilon\in[\num{0.1};\num{0.2}]$ de la V2
      (une graine par voisin aujourd'hui) tient-il ? Faible valeur depuis la V4.\\
    Écartés par la V4 & Ancre croisée et contrainte active sur la V3 (\cref{sec:campagnev4}) ; une
      contrainte d'égalité ($\mu$ de signe libre) serait une autre méthode, sans indice qu'elle aide :
      le score s'effondre dès que l'élève ajuste sa cible.\\
    V\_Fusion, seconde génération (facultatif) & Avec un plafond de confiance actif dès le premier
      étage, des images non saturées. Avant tout bras, une porte de sensibilité, en deux minutes : sur le
      modèle bruité de la porte d, la perte de distillation sur l'ancre, contre celle sur des vues du
      forget set (première génération : \numrange{0.04}{0.06} contre \num{2.17}). Si l'écart reste
      d'un ordre de grandeur, l'ancre ne verra pas les dégâts de la montée, et V\_Fusion s'arrête
      (\cref{sec:campagnevfusion}).\\
    Écartés par la V5.1 & Une bande en marge sans frein : elle dépasse sa bande en quelques pas
      (\cref{sec:campagnev51}). Les bras de protection des poids de V\_Fusion (gel, EWC, L2)
      supposaient une ancre synthétique qui tienne.\\
    \bottomrule
  \end{tabularx}
\end{table}

\section{Limites}\label{sec:limites}
\begin{itemize}[leftmargin=1.2em]
  \item La doublure du gold est une hypothèse, non une mesure, et les recadrages avantagent les
    méthodes qui visent des cibles de même nature (\code{gold\_soft}, la V2 et la V3 de
    \code{trust\_region}). La seconde famille de \code{bench-v4}, le masquage, est moins circulaire
    mais reste \enquote{l'original sur des entrées perturbées} : elle n'est pas indépendante de la
    première par construction. Les deux familles ne s'accordent pourtant pas sur les images
    difficiles (corrélation de rang nulle, \cref{sec:campagnev5}) : au moins l'une des deux, peut-être
    les deux, se trompe sur la difficulté par image du vrai gold, et le banc ne peut pas dire laquelle.
    Le score local mesure une distance à ces hypothèses.
  \item Les gains annoncés pour la V3 sont des plafonds d'un modèle à logits libres, sur les
    seules images exactes ; aucun run à pleine échelle ne les a encore confrontés au réseau, et la
    campagne en affaiblit déjà une prémisse (\cref{sec:trustv3}).
  \item \code{scrub} et \code{gold\_soft} tirent leur rétention d'images proxy extérieures au forget
    set : ce sont des références légitimes sur le banc, mais une méthode soumise ne peut pas en
    dépendre~\cite{sujet2026}.
  \item Le diagnostic (\cref{sec:diagnostic}) dit comment les V1 à V3 oublient ; les runs de la V4
    disent que l'ancre croisée ne corrige pas ce mouvement (\cref{sec:campagnev4}), pas qu'aucune
    ancre tirée du forget set ne le peut.
  \item La dose de SCRUB conforme a été choisie sur le banc lui-même (quatre doses, puis trois graines
    à la dose retenue) : un biais de sélection léger, petit devant l'écart mesuré sur la V3
    (\num{+0.029} apparié, près de $10\sigma$). Le plancher de retain de \num{0.80} est une convention
    posée avant les runs, pas une propriété du score : à lr \num{4e-4}, le banc donne \SI{11}{\percent}
    de plus avec un retain de \num{0.73}.
  \item La V5.1 et V\_Fusion n'ont qu'une graine par réglage. Leurs verdicts reposent sur des écarts
    nets (arrêts durs, plancher manqué de \num{0.075}, effondrements). V\_Fusion n'a qu'une génération,
    aux réglages proposés et non mesurés ; son étape 0 a tourné hors du registre, sa provenance est dans
    son dossier de run.
  \item Les bras de l'ancre et de la contrainte n'ont qu'une ou deux graines. Ils réfutent par des
    écarts de \numrange{0.02}{0.04}, bien au-delà du bruit, mais leurs petits effets (audit/exact
    \SI{-20}{\percent} pour l'ancre à lr \num{3e-5}) reposent sur une seule graine.
  \item Les max-steps de \code{scrub} s'entraînent sur les familles d'altérations du banc ; le gain de
    SCRUB conforme est le même sous les deux familles de doublures et se retrouve sur les seules images
    exactes, mais le banc ne peut pas exclure tout lien entre ces altérations et ses doublures.
  \item La dérive hors score de \code{bench-v4} ne porte que sur une vue forte par image d'audit, et
    ces images sont elles-mêmes mémorisées.
  \item Le fold d'audit est mémorisé par l'original ; si le control set officiel est fait d'images
    jamais vues, la dérive réelle peut différer.
  \item Les copies altérées sont simulées ; une altération d'un autre type pourrait se comporter
    autrement.
  \item Le serveur évalue en float16 et nous entraînons en bfloat16 : l'écart est faible mais non
    nul sur un modèle aux logits aussi grands.
  \item Les balayages de réglages sont à une graine ; seuls les points clés en ont trois, ce qui
    reste peu pour \code{scrub} (écart-type \num{0.022}).
  \item Les références n'ont été dosées que par leur taux (et le nombre d'époques pour
    \code{gold\_soft}) ; un réglage plus fin pourrait encore les améliorer.
  \item La cible de la V2 et les doublures du gold sous recadrage sont tirées de la même famille de
    vues. Le gain de la V2 ne dépend pas de leurs aires, mais il est deux fois et demie plus petit
    sous la famille de masquage (\cref{tab:v4}) : la part du gain qui se transférerait au vrai gold
    reste inconnue.
  \item Les réglages de la V3 ont été balayés à une graine ; seul le meilleur en a trois. Comme son
    rayon dépend de sa cible, ses balayages mêlent l'effet de la cible et celui du rayon.
  \item V\_Dumb n'a été mesurée qu'à une graine, et son poids d'ancre balayé à un seul taux
    (\num{e-5}) : la falaise est mesurée à cette granularité, pas exclue ailleurs.
  \item L'étape 0 de la V5 juge la sélection contre des doublures, pas contre le gold : à une graine de
    vues, sur les 2800 images du banc. Son seuil de précision était inatteignable sous $m_{15}$
    (\cref{sec:campagnev5}) ; le verdict n'en dépend pas. Les runs de la V5 n'ont qu'une ou deux
    graines ; sa localité porte sur des images d'audit mémorisées, pas sur leurs vues, et
    \enquote{la localité ne paie pas} n'est mesuré que pour la bande, contre des doublures.
\end{itemize}

\section*{Journal des versions}
\addcontentsline{toc}{section}{Journal des versions}
\begin{description}[leftmargin=1.2em,font=\sffamily\bfseries\color{bleu}]
  \item[21 novembre 2026] Première version : sujet et score exact, mesures sur le modèle
    (mémorisation, calibration des doublures, oubli propre à chaque image), méthode
    \code{trust\_region}, banc \code{bench-v3} aux formules du sujet, essais réduits, expériences
    prêtes.
  \item[21 novembre 2026, soir] Installation en une commande sur un GPU personnel (RTX 5080) :
    environnement, contrôle de la carte, données et premier passage du code CUDA. Aucun nouveau
    résultat.
  \item[21 novembre 2026, nuit] Mise en service de la RTX 5080 : garde GPU locale (un job à la fois,
    pause thermique, télémétrie), dimensionnement mesuré, modèle de spécification en pleine
    échelle noté par le banc. Deux runs de calibration à pleine échelle : l'oubli n'y est plus
    propre à chaque image (ratio \numrange{1.1}{1.2}), \code{score\_min} nul pour les deux. Le
    registre n'entre plus dans le drapeau \code{dirty} des runs (il change à chaque run).
  \item[21 novembre 2026, nuit, V2] V2 de \code{trust\_region} (\cref{sec:trustv2}) : cible propre à
    chaque image, tirée de vues fortes ; divergence des cibles mesurée sur le vrai modèle
    (\cref{tab:cible}), choix de $\tau=4$ ; spécification \code{trust-region-views} et essai
    réduit (\cref{tab:smoke}).
  \item[21 novembre 2026, nuit, suite] Le balayage des rayons montre que l'arrêt, testé tous les
    10 pas, ne suit pas $\varepsilon$ : $\varepsilon\in\{0{,}2;0{,}4;0{,}6\}$ s'arrêtent au même
    pas et donnent le même modèle. Nouveau réglage \code{radius\_check\_every} (10 par défaut, le
    comportement mesuré jusqu'ici) pour tester le rayon à chaque pas.
  \item[21 novembre 2026, nuit, relecture V2] Relecture statique du code de la V2 : documentation
    précisée (la cible V1 est rendue en logits, coût mémoire du calcul des vues sur tout le
    forget set) ; estimation de durée de \code{trust-region-views} corrigée (les trois vues coûtent
    environ une époque d'entraînement, non \SI{1}{\percent}). Aucun nouveau résultat.
  \item[21 novembre 2026, nuit, campagne] Campagne \code{trust\_region} à pleine échelle (41 runs,
    fiche \chemin{experiments/unlearning/trust-region-campaign.md}, tableau
    \chemin{experiments/unlearning/SCOREBOARD.md}) : rayon testé à chaque pas, rayons, taux,
    warmup, température, blocs, lot, graines.
  \item[21 novembre 2026, nuit, bilan] Résultats de la campagne et des références dosées, 54 runs
    en tout. La région de confiance ne dépend que de son rayon : \num{0.0637 +- 0.0025} à
    $\varepsilon=\num{0.15}$. Dosé, \code{scrub} fait mieux en moyenne, avec
    \num{0.098 +- 0.022} ; dosé, \code{gold\_soft} fait jeu égal, avec \num{0.059 +- 0.012}.
    Résumé, méthode, résultats, suite et limites mis à jour.
  \item[21 novembre 2026, fusion] Fusion de la branche distante (V2 de \code{trust\_region} et sa
    relecture) avec la campagne locale : code de la V2 repris tel quel, \code{radius\_check\_every}
    ajouté à côté de \code{trust\_target}. La fusion reproduit exactement le run de référence de la
    V1 et fait tourner la V2 en essai réduit sur CUDA (\cref{sec:fusion}). Rapport complété :
    détail des références et réglages par défaut, protocole et outillage, meilleure performance,
    inventaire des expériences, bilan par méthode.
  \item[21 novembre 2026, nuit, campagne V2] Préparation de la campagne V2 : cible par vues
    mesurée avant tout entraînement sur les 2800 images entraînées (\chemin{scripts/unlearning\_v2\_targets.py}),
    fiche \chemin{experiments/unlearning/trust-region-v2-campaign.md} (hypothèses posées avant les
    runs), tableau comparatif qui distingue V1 et V2, documentation hors du drapeau \code{dirty}.
  \item[21 novembre 2026, nuit, V3] V3 de \code{trust\_region} (\cref{sec:trustv3}) : cible
    calibrée en forme close, rayon tiré de la cible, copies douces dans les lots ; banc
    \code{bench-v4} : seconde famille de doublures (masquage de patchs, taux mesurés sur le vrai
    modèle, \cref{tab:familles}), pire cas sur les six doublures et par famille, dérive hors score
    sur des vues du fold d'audit ; spécification \code{trust-region-calibrated} et essai réduit
    (\cref{tab:smokev4}). Les V1 et V2 sont inchangées au bit près.
  \item[21 novembre 2026, seconde fusion] Fusion de la V3 et de \code{bench-v4} avec la campagne :
    code et théorie de la V3 repris, résultats de la campagne gardés tels quels. Le tableau
    comparatif sépare les protocoles du banc (\code{bench.protocol\_version}) ; les fiches et
    spécifications déjà lancées restent en \code{bench-v3} ; les spécifications de la V2 et de la V3
    testent le rayon à chaque pas, la V2 sur les rayons \numrange{0.1}{0.3}, comme la campagne le
    recommandait ; la théorie de la V3 est confrontée à la campagne (une prédiction vérifiée, une
    prémisse affaiblie) ; conformité des ancres de rétention au sujet notée. 526 tests ; la V1 garde
    exactement les mêmes poids.
  \item[21 novembre 2026, nuit, bilan V2] Campagne V2, 30 runs en six étapes posées chacune avant
    d'être lancée : la V2 à $\tau=1$ vaut \num{0.0831 +- 0.0042}, \SI{30}{\percent} de plus que la V1,
    apparié ; $\tau$ décide, ni l'aire ni le nombre de vues ; sous $\tau=1$, le rayon n'est parfois
    jamais atteint. Résumé, résultats, meilleure performance, inventaire, bilan, suite et limites
    mis à jour.
  \item[22 novembre 2026, troisième fusion] Fusion de la V3 et de \code{bench-v4} (branche
    distante) avec la campagne V2 : code, théorie et banc de la V3 repris ; résultats des deux
    campagnes gardés, en \code{bench-v3}. Le tableau comparatif sépare les protocoles et étiquette
    V1, V2 et V3. Conformité appliquée aux conclusions : \code{scrub} et \code{gold\_soft} restent des
    références de banc, la V2 est la meilleure méthode conforme mesurée. La spécification
    \code{trust-region-views} révisée correspond à l'étape 1 de la campagne V2, déjà mesurée.
  \item[22 novembre 2026, campagne V3] Préparation de la campagne V3 : cible calibrée mesurée avant
    tout entraînement (560 cibles, \chemin{scripts/unlearning\_v3\_targets.py}), références
    \code{bench-v4} à relancer au lab, fiche \chemin{experiments/unlearning/trust-region-v3-campaign.md}
    (seuils chiffrés et étapes posés avant les runs).
  \item[22 novembre 2026, bilan V3] Campagne V3, 23 runs en sept étapes et 9 références relancées
    en \code{bench-v4} : jeu égal avec la V2 (\num{0.0798 +- 0.0012} contre \num{0.0786 +- 0.0030}) au
    meilleur réglage (sans couverture) ; gain, mécanisme par $S_{\text{ms}}$ et copies douces
    réfutés ; tous les réglages de la cible agissent par le rayon. Résumé, résultats, inventaire,
    bilan, suite et limites mis à jour.
  \item[22 novembre 2026, synthèse] Rapport réorganisé autour d'un chemin unique : vue d'ensemble
    de la V1 à la V3 en tête des résultats, comparaison finale sous un même protocole
    (\cref{tab:synthese,fig:synthese}) ; fiche de reprise des expériences pour l'équipe
    (\chemin{experiments/unlearning/ETAT.md}).
  \item[22 novembre 2026, V4] Écrite en parallèle de la campagne V3, sans la connaître. Étude du
    blocage des V1--V3, en trois volets : revue de littérature (sources A* et A du classement CORE,
    venues vérifiées), relecture des 54 runs de la V1 et du code, coordination. Diagnostic tiré du
    registre (\cref{sec:diagnostic}) : environ \SI{88}{\percent} du mouvement de la V1 est commun à
    toutes les images mémorisées ; le \enquote{ratio 1,4} vaut \num{1.14} sur les images exactes. Un
    essai intermédiaire de la V3 sur M2 (\cref{sec:intermediaire}). V4 (\cref{sec:trustv4}), codée
    et couverte par les tests unitaires, sans run : ancre croisée entre plis du forget set
    (\code{anchor\_folds}, \code{anchor\_weight}, \code{anchor\_views}), contrainte active au lieu de
    l'arrêt (\code{radius\_dual\_lr}), rétention par vues fortes du forget set pour \code{scrub} et
    \code{gold\_soft} (\code{retain\_source}) ; douze défauts relevés par la relecture, corrigés. Le
    banc donne aussi $S_{\text{ms}}$ et $S_{\text{pred}}$ sur les seules images exactes
    (\code{bench.exact.*}, hors score, protocole inchangé) ; le tableau comparatif sépare les
    réglages texte (\code{retain\_source}).
  \item[22 novembre 2026, V5] Étude d'une méthode par montée de gradient, en quatre volets :
    revue de littérature (venues A* et A du classement CORE, journaux Q1), lecture du registre et du
    code avec une sonde exploratoire (Apple M2, 320 images, une graine), vérification indépendante
    des sources, coordination. V5 (\cref{sec:v5}), conçue, non codée, aucun run : montée bornée par
    image en bande, réservée à la fraction des images où les vues doutent le plus, dosée par niveaux
    choisis sur un pli témoin jamais touché, corrigée de l'échelle, avec un moniteur préquentiel et
    un arrêt dur. Modèle linéarisé de la fuite (\cref{eq:v5lin,eq:v5ratio,eq:v5hold}), dérivation
    de l'équipe, non testée. Étape 0 et critères de réfutation posés avant tout run ; la projection
    hors du gradient moyen est écartée.
  \item[22 novembre 2026, V5 codée] \code{band\_ascent}, nouvelle méthode (\cref{sec:v5}) : bande,
    sélection, niveaux de dose, validation croisée à 5 plis, moniteur préquentiel, arrêt dur avec
    retour au dernier niveau achevé ; les V1 à V4 sont inchangées. Diagnostic de l'étape 0
    (sous-commande \code{diagnose}), sans entraînement. Deux spécifications et leurs fiches
    (\code{band-ascent-step0}, \code{band-ascent}), seuils posés avant tout run ; l'option de
    direction n'est pas codée. Huit tests unitaires (vingt cas) et un test du diagnostic ; aucun run.
  \item[22 novembre 2026, quatrième fusion] Fusion de la V4 et de la V5 (conçue, puis codée) avec les
    campagnes V2 et V3 : code, résultats et fiches des trois côtés gardés. Les résultats de la V3 priment sur le pronostic de la
    V4, qui la plaçait au niveau de la V1 : elle la dépasse, par la direction de sa cible. Ils
    renforcent en revanche son diagnostic : relu sur le registre fusionné, le fold d'audit de la V2
    et de la V3 retenues bouge autant que celui de la V1 (audit/exact \num{0.86} à \num{0.91}), la
    localité est bien le défaut commun. Le rayon pris sur la sonde ne vaut plus que pour les options
    de la V4 : la V1, la V2 et la V3 se rejouent exactement (\code{ref4-*}). Les spécifications de la
    V4 partent de la V3 retenue ($b=0$, sans copies douces, qui n'apportaient rien), se comparent
    appariées à ses runs mesurés et abandonnent $\lambda=\num{0.5}$, déjà réfuté. 558 tests du kit ; V1, V2
    et V3 gardent exactement les mêmes poids.
  \item[22 novembre 2026, V\_Dumb] \code{dumb}, nouvelle méthode (\cref{sec:vdumb}) : la plus simple
    possible, une montée de la perte de chaque image oubliée coupée à une cible (\num{0.45}, perte
    mesurée de l'original sur une vue forte d'aire \num{0.50}) et une ancre KL sur les vues fortes du
    forget set ; un réglage nouveau (\code{dumb\_target\_loss}), les autres méthodes inchangées. Deux
    tests unitaires, une spécification \code{v-dumb} et sa fiche, seuils posés avant tout run ; aucun
    run.
  \item[22 novembre 2026, campagne V4] Campagne V4 sur la RTX 5080 (\cref{sec:campagnev4}), fiche
    \chemin{experiments/unlearning/trust-region-v4-campaign.md} : 26 runs, 6 rejeux de références
    identiques au bit près (les options de la V4 sont neutres à leur défaut), jugés contre les seuils
    posés avant, quatre suivis posés chacun avant son lancement. SCRUB conforme confirmé une fois
    redosé (\num{0.1085 +- 0.0068}, meilleure méthode conforme mesurée) ; ancre croisée et
    contrainte active réfutées ; V3 à lr \num{3e-5} pas mieux. Résumé, vue d'ensemble
    (\cref{tab:synthese,fig:synthese}), meilleure performance, inventaire, bilan, suite et limites mis
    à jour ; la théorie de la V4 et de la V5 est inchangée. Le tableau comparatif compte chaque graine
    une fois (rejeux et doublons \code{dirty}) ; deux tests du kit en plus (560).
  \item[22 novembre 2026, cinquième fusion] Fusion de V\_Dumb (branche distante, codée, sans run)
    avec la campagne V4 : code, théorie, fiche et spécification de V\_Dumb repris tels quels,
    résultats de la V4 gardés. Son ancre est celle de SCRUB conforme ; la suite proposée la compare
    aussi à \code{ref4-scrub-fv} (\cref{tab:plan}).
  \item[22 novembre 2026, campagne V5] Campagne V5 sur la RTX 5080 (\cref{sec:campagnev5}), fiche
    \chemin{experiments/unlearning/trust-region-v5-campaign.md} : étape 0 sans entraînement
    (\code{band-ascent-step0}), jugée contre les seuils posés avant. Réfutée sur ses deux critères,
    sous les trois masquages ; \code{band-ascent} non lancé, comme le prévoyait la fiche. Mesure
    annexe, posée avant d'être lancée (\chemin{scripts/unlearning_v5_families.py}) : les deux familles
    de doublures ne s'accordent pas sur les images difficiles. Résumé, vue d'ensemble
    (\cref{tab:synthese,fig:synthese}), inventaire, bilan, suite et limites mis à jour ; la théorie de
    la V5 est inchangée, son paragraphe d'état renvoie à la campagne. Le tableau comparatif distingue
    désormais une option mise à \code{None} (\code{band\_floor=none}) du défaut ; un test du kit en
    plus (563). Puis, sur décision de l'utilisateur, \code{band-ascent} lancé comme mesure de la
    localité seule, et un suivi posé avant (\code{band-locality}, \code{band-ascent-s1}) : huit runs.
    Localité obtenue, la plus forte mesurée à KL exacte égale (\cref{fig:v5localite}), sans rien
    devoir à la sélection ; borne haute indispensable ; score nul ; V5 abandonnée.
  \item[22 novembre 2026, campagne V\_Dumb] Campagne V\_Dumb sur la RTX 5080
    (\cref{sec:campagnevdumb}), fiche \chemin{experiments/unlearning/trust-region-vdumb-campaign.md} :
    onze runs, \code{v-dumb} jugée contre ses seuils posés avant, puis deux suivis de dose posés chacun
    avant d'être lancés. Réfutée à sa fiche (immobile), sans dose utilisable (falaise entre
    l'immobilité et l'effondrement), charnière sans effet de borne. Résumé, vue d'ensemble,
    inventaire, bilan, suite, limites mis à jour ; la théorie de V\_Dumb est inchangée, son paragraphe
    d'état renvoie à la campagne.
  \item[22 novembre 2026, V5.1] Objectif fixé par le responsable : un retain d'au moins \num{0.8},
    budget dépensé vers le gold. Section V5.1 (\cref{sec:v51}), conçue sans code ni run à partir d'une
    revue de la littérature et d'une exploration du code : trois défauts de la V5 (montée affamée,
    borne des images tenues qui ne serre pas au budget, poussée des images tenues dès le premier
    pas, ce qui corrige une phrase de la section V5), trois correctifs (bande en marge, retenue en KL
    à $T=2$, critère en retain estimé) et une option d'encodeur (vote kNN ajusté par plis), à n'activer
    que si l'étape 0 la justifie. Chiffres du registre relus ; simulation du modèle linéarisé
    (tendances seulement). Spécification \code{band-ascent-budget} et sa fiche, sans code, seuils
    posés avant tout run.
  \item[22 novembre 2026, V5.1 codée] Trois options de \code{band\_ascent}, neutres par défaut
    (\code{band\_space}, \code{band\_hold\_kl}, \code{cv\_min\_retain}) ; \code{diagnose} lit en
    plus le vote kNN de l'encodeur et les trois défauts de la V5 ; onze tests unitaires, 210 tests de
    la tâche au vert ; spécification \code{band-ascent-v51} (V5.1 et trois ablations) et sa fiche,
    lectures V5.1 ajoutées à la fiche de l'étape 0, seuils posés avant tout run ; aucun run.
  \item[22 novembre 2026, sixième fusion] Fusion de la V5.1 (locale, codée, sans run) avec les
    campagnes V4, V5 et V\_Dumb (distantes) : résultats, fiches et code des deux côtés gardés. La
    V5.1 a été conçue avant le verdict de la V5 ; sa base étant abandonnée, elle n'est pas lancée,
    et ses lignes du plan (étape 0 étendue, \code{band-ascent-budget}, \code{band-ascent-v51}) en
    sortent. Bilan, suite et résumé l'indiquent.
  \item[22 novembre 2026, V\_Fusion] Section V\_Fusion (\cref{sec:vfusion}) : SCRUB conforme dont
    l'ancre devient un lot d'images synthétiques tirées du seul original, en deux lots ; gel, L2
    et EWC rangés en bras. Code : générateur, cinq portes de qualité, sous-commandes
    \code{synth-generate} et \code{synth-gates}, job \chemin{jureca/synthetic.sbatch} ; neuf tests
    unitaires sur un petit ViT ; aucun run, aucune image synthétique n'existe encore.
  \item[22 novembre 2026, ancre V\_Fusion] Nouvelle ancre de \code{scrub}, faite d'images
    synthétiques : le réglage \code{retain\_source} prend la valeur \code{synthetic}, avec deux
    réglages de plus (\code{synthetic\_path}, \code{synthetic\_pools}) et deux tests unitaires
    (589 tests du kit au vert). Spécification
    \code{v-fusion-anchor} et sa fiche, seuils posés avant tout run de bras. Code figé avant ces
    runs ; résultats au commit suivant.
  \item[22 novembre 2026, campagnes V5.1 et V\_Fusion] V5.1 mesurée à la demande de l'utilisateur
    (un diagnostic, onze runs, \cref{sec:campagnev51}) : défauts confirmés, la marge dépasse la bande,
    abandonnée. V\_Fusion (\cref{sec:campagnevfusion}) : 2000 images synthétiques sur la RTX 5080, portes
    a, b, c, c' passées, d ininterprétable ; quatre bras d'ancre effondrés, abandonnée sous cette
    génération. Résumé, vue d'ensemble, inventaire, bilan, suite et limites mis à jour. Meilleur
    conforme inchangé : SCRUB conforme.
  \item[22 novembre 2026, V\_Dumb2] Section V\_Dumb2 (\cref{sec:vdumb2}) : le niveau par image de
    V\_Dumb, écrit en bande sur la marge et borné des deux côtés, à la place du plafond moyen des
    max-steps de SCRUB conforme. Option \code{scrub\_target\_loss} de \code{scrub}, neutre par défaut
    (V4 rejouée au bit près sur le petit modèle des tests) ; plafond \code{level\_ceiling} ajouté à
    l'étape 0 (\code{diagnose}) ; quatre tests unitaires ; spécification \code{v-dumb2} et sa fiche,
    porte de l'étape 0 et seuils posés avant tout run ; aucun run.
  \item[22 novembre 2026, septième fusion] Fusion de V\_Dumb2 (locale, codée, sans run) avec les
    campagnes V5.1 et V\_Fusion et l'ancre synthétique de \code{scrub} (distantes) : code, résultats
    et fiches des deux côtés gardés. La section V\_Dumb2 cite désormais la mesure de la V5.1 : la
    bande en marge dépasse son plafond en quelques pas dès lr \num{3e-5}, un risque pour V\_Dumb2,
    qui l'emploie à \num{3e-4} et \num{4e-4} sans arrêt dur. Seuils de \code{v-dumb2} inchangés.
\end{description}

\printbibliography[heading=bibintoc,title={Références}]

\end{document}
